% PHP Gallery Úplná dokumentace produktu
% Project: PHP Gallery
% Repository: https://github.com/klusik/PHP_gallery
% File: docs/PHP_Gallery_Manual_CZ.tex
% Module Type: Produktová příručka v~LaTeXu
% Purpose: Vysvětlit pracovní postupy správce, provoz a~údržbu PHP Gallery.
% Responsibilities:
%   - Udržovat uživatelskou příručku a~technickou referenci při zachování reprodukovatelné sazby PDF.
% Author: Rudolf Klusal
% Sestavte v~tomto adresáři pomocí:
%   pdflatex PHP_Gallery_Manual_CZ.tex
%   makeindex PHP_Gallery_Manual_CZ.idx
%   pdflatex PHP_Gallery_Manual_CZ.tex
%   pdflatex PHP_Gallery_Manual_CZ.tex
%
% Soubory Markdown a~implementace uvedené v~odkazech představují
% redakční a~technické zdroje tohoto vydání. Zkontrolujte místní projekt
% při každé změně verze aplikace.

\documentclass[11pt,a4paper]{article}

\usepackage[utf8]{inputenc}
\usepackage[T1]{fontenc}
\usepackage{lmodern}
\usepackage[czech]{babel}
\usepackage{geometry}
\usepackage{microtype}
\usepackage{xcolor}
\usepackage{graphicx}
\usepackage{booktabs}
\usepackage{array}
\usepackage{ragged2e}
\usepackage{longtable}
\usepackage{tabularx}
\usepackage{enumitem}
\usepackage{listings}
\usepackage{fancyhdr}
\usepackage{titlesec}
\usepackage[bottom,hang,flushmargin]{footmisc}
\usepackage{makeidx}
\usepackage{hyperref}
\usepackage{bookmark}

\geometry{top=23mm,bottom=25mm,left=24mm,right=24mm,headheight=15pt}
\setlength{\parindent}{0pt}
\setlength{\emergencystretch}{1.5em}
\setlength{\parskip}{0.65em}
\setcounter{tocdepth}{3}
\setcounter{secnumdepth}{3}
\makeatletter
% Zachovejte odstup dvouciferných čísel pododdílů od jejich názvů v~obsahu.
\renewcommand*\l@subsection{\@dottedtocline{2}{1.5em}{4.1em}}
\makeatother
\setlist{itemsep=0.3em,topsep=0.45em}
\renewcommand{\arraystretch}{1.2}

\definecolor{GalleryBlue}{HTML}{245A86}
\definecolor{GalleryLightBlue}{HTML}{EAF2F8}
\definecolor{GalleryDark}{HTML}{263238}
\definecolor{GalleryGray}{HTML}{F3F5F7}
\definecolor{GalleryCode}{HTML}{1F2933}
\definecolor{GalleryGreen}{HTML}{2E6B45}
\definecolor{GalleryRed}{HTML}{8A2F2F}

% Metadata vydání jsou zde, takže aktualizace vydání vyžaduje jedinou úpravu verze.
\newcommand{\version}{0.113.2}
\newcommand{\manualdate}{1.~října~2026}

\hypersetup{
  colorlinks=true,
  linkcolor=GalleryBlue,
  citecolor=GalleryGreen,
  urlcolor=GalleryBlue,
  pdftitle={PHP Gallery Úplná dokumentace produktu},
  pdfauthor={Rudolf Klusal},
  pdfsubject={Referenční příručka instalace, správy, architektury, zabezpečení, údržby a~vývoje},
  pdfkeywords={PHP Gallery, galerijní CMS, správa, architektura, dokumentace},
  pdfdisplaydoctitle=true,
  bookmarksopen=true,
  bookmarksnumbered=true
}

\pagestyle{fancy}
\fancyhf{}
\fancyhead[L]{\small PHP Gallery}
\fancyhead[R]{\small Úplná dokumentace produktu}
\fancyfoot[L]{\small Verze \version}
\fancyfoot[C]{\thepage}
\fancyfoot[R]{\small \manualdate}
\renewcommand{\headrulewidth}{0.4pt}
\renewcommand{\footrulewidth}{0.2pt}

\titleformat{\section}{\Large\bfseries\color{GalleryBlue}}{\thesection}{0.75em}{}
\titleformat{\subsection}{\large\bfseries\color{GalleryDark}}{\thesubsection}{0.75em}{}
\titleformat{\subsubsection}{\normalsize\bfseries\color{GalleryDark}}{\thesubsubsection}{0.75em}{}
\pretocmd{\section}{\clearpage}{}{}

\lstdefinestyle{gallery}{
  basicstyle=\ttfamily\small\color{GalleryCode},
  backgroundcolor=\color{GalleryGray},
  frame=single,
  rulecolor=\color{GalleryLightBlue},
  breaklines=true,
  breakatwhitespace=false,
  columns=fullflexible,
  keepspaces=true,
  showstringspaces=false,
  tabsize=4,
  xleftmargin=0.8em,
  xrightmargin=0.8em,
  framexleftmargin=0.4em,
  aboveskip=0.8em,
  belowskip=0.8em
}
\lstset{style=gallery}

\newcommand{\product}{\textbf{PHP Gallery}}
\newcommand{\setting}[1]{\nolinkurl{#1}}
\newcommand{\filepath}[1]{\path{#1}}
\newcommand{\ui}[1]{\textbf{#1}}
\newcommand{\codeword}[1]{\nolinkurl{#1}}
\newcommand{\important}[1]{%
  \begin{center}\fcolorbox{GalleryBlue}{GalleryLightBlue}{\parbox{0.91\linewidth}{\textbf{Důležité.} #1}}\end{center}}
\newcommand{\securitynote}[1]{%
  \begin{center}\fcolorbox{GalleryRed}{GalleryGray}{\parbox{0.91\linewidth}{\textbf{Poznámka k~zabezpečení.} #1}}\end{center}}
\newcommand{\performancenote}[1]{%
  \begin{center}\fcolorbox{GalleryGreen}{GalleryGray}{\parbox{0.91\linewidth}{\textbf{Poznámka k~výkonu.} #1}}\end{center}}

\makeindex

\begin{document}

\pagenumbering{gobble}
\begin{titlepage}
  \centering
  \vspace*{24mm}
  {\Huge\bfseries\color{GalleryBlue} PHP Gallery\par}
  \vspace{6mm}
  {\LARGE Úplná dokumentace produktu\par}
  \vspace{14mm}
  {\large Referenční příručka instalace, správy, architektury, zabezpečení, údržby a~vývoje\par}
  \vfill
  \begin{tabular}{rl}
    \textbf{Verze aplikace:} & \version\\
    \textbf{Vydání příručky:} & \manualdate\\
    \textbf{Formát dokumentu:} & A4, PDF s~rejstříkem a~odkazy\\
    \textbf{Autor projektu:} & Rudolf Klusal\\
    \textbf{Licence:} & Licence MIT\\
  \end{tabular}
  \vfill
  {\small Průvodce vytvářením, uspořádáním, prezentací, ochranou, sdílením a~uchováváním fotografické sbírky v~PHP Gallery verze \version.\par}
\end{titlepage}

\pagenumbering{roman}
\begingroup
\small
\setlength{\parskip}{0pt}
\tableofcontents
\endgroup
\clearpage
\pagenumbering{arabic}

\section{Účel a~průvodce čtením}
\index{d00033@dokumentace}
\index{p00189@příručka!r00196@rozsah}

\product{} je fotografická galerie provozovaná na vlastním serveru. Fotografie zůstávají na úložišti, které máte pod kontrolou; aplikace k~nim přidává katalog, navigaci, pravidla přístupu, náhledy, mapy, vyhledávání, stahování a~nástroje správy. Jednu instalaci lze využít jako veřejné portfolio, rodinný archiv, pro předávání zakázek klientům, fotografie z~akcí, cestovatelské sbírky nebo rozsáhlou osobní knihovnu.

Příručka je uspořádána podle činností, které vlastník skutečně provádí. Úvodní oddíly se věnují instalaci a~první galerii; střed příručky pojednává o~zveřejňování, médiích, soukromí, prezentaci a~údržbě. Informace o~architektuře a~vývoji jsou zařazeny ke konci, aby nepřerušovaly běžnou správu.

Při nové instalaci začněte oddílem~\ref{sec:first-hour}. Plánování sbírky popisuje oddíl~\ref{sec:planning}; běžné úpravy galerií a~fotografií najdete v~oddílech~\ref{sec:create-edit-gallery} a~\ref{sec:photo-workflow}. Oddíl~\ref{sec:audiences} je užitečný před zveřejněním, protože porovnává několik běžných skupin návštěvníků. Obsah a~rejstřík umožňují přejít přímo k~nastavení nebo pracovnímu postupu, když web již používáte.

\subsection{Jak tuto příručku používat}
\index{p00189@příručka!p00151@použití}

Názvy v~administraci, například \ui{Vzhled} a~\ui{Stav systému}, jsou vyznačeny tučně. Neproporcionální písmo označuje názvy souborů, nastavení, hodnoty nebo příkazy. Postupy, u~nichž záleží na pořadí, jsou číslované; kontrolní seznamy a~referenční informace používají odrážky nebo tabulky. Technické poznámky k~implementaci jsou uvedeny pouze tam, kde vysvětlují viditelné chování, rozhodnutí při obnově nebo požadavek na údržbu.

\subsection{Další projektové zdroje}

Zdroje pro vývojáře a~provozovatele hostingu zůstávají oddělené od této příručky pro vlastníka. Oddíl se zdroji uvádí dokumenty repozitáře použité pro architekturu, databázi, vlastnictví zdrojového kódu, testování a~pravidla přispívání \cite{readme,architecture,database,codemap,testing,agents}.\footnote{Tyto soubory jsou vhodnějším zdrojem při změnách kódu, zkoumání podrobností schématu nebo přípravě vydání.}

\subsection{Terminologie}

\begin{description}
  \item[Galerie] Pojmenovaná sbírka fotografií. Galerie může obsahovat fotografie i~menší galerie.
  \item[Fotografie nebo obrázek] Jedna položka galerie spolu s~názvem, popisem, štítky, viditelností a~dostupnými údaji o~fotoaparátu.
  \item[Náhled] Menší kopie vytvořená pro rychlé zobrazení v~mřížkách galerií, výsledcích vyhledávání, titulních obrázcích a~náhledech.
  \item[Návštěvník] Osoba prohlížející veřejnou část webu.
  \item[Správce] Důvěryhodná osoba, která se přihlašuje, aby vytvářela a~spravovala obsah a~nastavení.
  \item[Vybraná galerie] Galerie právě otevřená na obrazovce.
  \item[Větev galerií] Galerie společně se všemi menšími galeriemi, které obsahuje.
\end{description}

\section{Přehled produktu}
\label{sec:overview}
\index{p00172@přehled produktu}
\index{s00205@sdílený hosting}

\product{} mění sbírku složek na fotografický web, který lze pohodlně procházet. Návštěvníci vidí titulní obrázky, názvy, popisy, mřížky fotografií, celoobrazovkový prohlížeč, mapy, štítky, vyhledávání a~stahování podle voleb vlastníka. Správci mají neveřejné nástroje pro přidávání fotografií, uspořádání sbírek, řízení soukromí, zlepšování popisů a~udržování instalace v~dobrém stavu.

Hlavním cílovým prostředím je běžný webhosting PHP/MySQL. Vlastník uchovává zdrojové fotografie v~rozpoznatelných složkách, zatímco databáze poskytuje katalog a~stav aplikace potřebný pro procházení.\footnote{Přehled projektu obsahuje stručné shrnutí instalace a~funkcí \cite{readme}.}

\subsection{Základní možnosti}

\begin{itemize}
  \item Vytvářejte galerie uvnitř dalších galerií, například \emph{2026 / Itálie / Řím} nebo \emph{Klienti / Northwind / Konečný výběr}.
  \item Každé galerii přiřaďte název, popis, rozsah dat, titulní obrázek, veřejnou adresu, nastavení soukromí a~vlastní vzhled.
  \item Nahrávejte fotografie prostřednictvím prohlížeče nebo vyhledávejte soubory zkopírované do složek galerií.
  \item Přidávejte k~fotografiím názvy, příběhy, štítky, viditelnost, pořadí, data a~údaje o~poloze.
  \item Umožněte návštěvníkům vyhledávat, procházet štítky, zobrazovat fotografie na celou obrazovku, používat mapy, hlasovat a~stahovat povolené sbírky.
  \item Chraňte rodinné, klientské nebo citlivé sbírky pomocí řízení přístupu a~způsobů sdílení.
  \item Vyhledávejte pravděpodobné duplicity fotografií a~posuzujte je pomocí přímých odkazů na obě umístění.
  \item Udržujte instalaci aktuální pomocí řízených aktualizací, kontrol stavu, záloh a~nástrojů údržby.
  \item Automatizujte nahrávání s~vymezeným rozsahem, přenášejte galerie mezi kompatibilními instalacemi a~připravujte stažení vybraných fotografií nebo galerií.
  \item Vytvářejte volitelná metadata pomocí AI a~návrhy překladů; každou vygenerovanou hodnotu lze před zveřejněním zkontrolovat.
  \item Publikujte výstupy robots a~obrazové mapy webu vhodné pro roboty pro veřejný obsah, ke kterému mají návštěvníci skutečně přístup.
\end{itemize}

\subsection{Spolupráce souborového systému a~databáze}
\index{n00104@návrh založený na souborovém systému}
\index{d00024@databáze}

PHP Gallery ukládá trvalá data do souborového systému i~databáze. Adresáře galerií obsahují zdrojové fotografie v~rozpoznatelném stromu. Databáze uchovává informace, které nepatří přímo do souboru: názvy, popisy, štítky, pořadí, nastavení přístupu, hlasy, data, záznamy o~generovaných médiích a~další stav aplikace.

Ani jedna část sama o~sobě nepředstavuje úplnou zálohu. Obnovení pouze databáze ponechá položky katalogu bez zdrojových souborů; obnovení pouze stromu galerií ztratí stav správy a~prezentace spojený s~těmito soubory. Zálohování a~obnova proto musí zacházet s~databází, adresáři galerií a~místní konfigurací jako s~jedním celkem. Oddíl~\ref{sec:user-data} se tomuto tématu věnuje podrobněji.

\section{Požadavky a~prostředí}
\label{sec:requirements}
\index{p00156@požadavky}
\index{p00134@PHP}
\index{m00091@MySQL}
\index{m00079@MariaDB}

Tento oddíl použijte při výběru hostingového tarifu nebo kontrole existujícího účtu. Tabulku požadavků můžete také zaslat přímo poskytovateli hostingu.

\subsection{Minimální běhové prostředí}

\begin{longtable}{>{\bfseries}>{\hspace{0pt}}p{0.27\linewidth}>{\hspace{0pt}}p{0.65\linewidth}}
\toprule
Součást & Požadavek\\
\midrule
\endhead
PHP & Verze 8.1 nebo novější.\\
Databáze & MySQL 5.7 nebo novější, případně MariaDB 10.2 nebo novější.\\
PDO & Rozšíření PDO MySQL pro dotazy aplikace a~migrace.\\
Zpracování obrázků & GD pro běžné generování náhledů. Další místní nástroje mohou podporovat specializované formáty obrázků.\\
Archivy & ZipArchive pro pracovní postupy s~balíčky, stahováním a~přenosem, které vytvářejí nebo zpracovávají data ZIP.\\
Souborový systém & Zapisovatelná umístění pro data aplikace, galerie, mezipaměť a~generovaná média.\\
Webový server & Pro čisté adresy URL se doporučuje přepisování ve stylu Apache; směrování pomocí parametrů dotazu podporuje kompatibilní prostředí bez přepisování.\\
\bottomrule
\end{longtable}

Odchozí přístup HTTPS podporuje aktualizace aplikace a~vzdálená metadata vydání. Místní provoz a~ručně nasazená vydání mohou fungovat i~s~přísnějším omezením odchozích spojení.\footnote{Síťová pravidla by měla odpovídat požadavkům vlastníka instalace na aktualizace a~soukromí. Aktualizační kód ověřuje balíčky a~uchovává data pro obnovu podle místní implementace aktualizátoru.}

\subsection{Oprávnění}
\index{o00127@oprávnění!s00214@souborový systém}

Webový proces potřebuje přístup pro čtení běhového kódu a~pro zápis do určených měnitelných umístění. Pokud to hosting umožňuje, ponechte běhový kód pouze pro čtení, s~výjimkou řízené aktualizace. Zdrojové soubory galerií si zaslouží stejné zálohování jako databáze, protože úplný stav aplikace zahrnuje obě oblasti.

\important{Zálohujte společně databázi, soubory galerií a~konfiguraci instalace. Záloha samotné databáze uchovává metadata, ale vynechává rozhodující strom médií.}

\section{Instalace a~počáteční nastavení}
\label{sec:installation}
\index{i00056@instalace}
\index{n00098@nastavení}

\subsection{Instalace v~prohlížeči}

\subsubsection{Pořadí instalace}

Instalátor v~prohlížeči shromáždí nastavení databáze, inicializuje schéma pomocí migrací, vytvoří prvního správce, připraví zapisovatelná umístění a~zapíše místní konfiguraci. Pokud požadovaná konfigurace chybí, nový požadavek přesměruje na instalaci.\footnote{Stav instalace a~vygenerovaná konfigurace vyžadují stejnou ochranu jako přihlašovací údaje. Zdrojové archivy musí nadále vynechávat místní \filepath{config.php}.}

Obvyklý postup:

\begin{enumerate}
  \item Vytvořte prázdnou databázi MySQL nebo MariaDB v~ovládacím panelu hostingu.
  \item Nahrajte balíček aplikace do zamýšleného webového adresáře.
  \item Otevřete web a~postupujte podle instalátoru.
  \item Zadejte údaje připojení k~databázi a~přihlašovací údaje prvního správce.
  \item Potvrďte zapisovatelné adresáře a~kontroly rozšíření.
  \item Dokončete instalaci a~vstupte do administrace.
\end{enumerate}

\subsection{Nastavení z~příkazového řádku}

\subsubsection{Místní příkazy}

Repozitář poskytuje prosté skripty PHP pro migrace a~vytvoření správce:

\begin{lstlisting}[language=bash]
php scripts/migrate.php
php scripts/create_admin.php <username> <password>
\end{lstlisting}

Nastavení z~příkazového řádku používá stejný model trvalého schématu jako nastavení v~prohlížeči. Prostudujte \filepath{config.example.php}, vytvořte místní konfiguraci, připravte databázi, spusťte migrace a~vytvořte prvního správce.

\subsection{První kontrola správy}

Po instalaci zkontrolujte:

\begin{itemize}
  \item Název webu, jazyk, vzhled, šířku stránky a~volby veřejného vykreslování.
  \item Oprávnění kořenového adresáře galerií a~výsledky vyhledání souborů.
  \item Formáty, rozměry, kvalitu a~stav údržby náhledů.
  \item Volby přístupu, obnovy hesla, e-mailu, telemetrie, aktualizačního kanálu a~záloh.
  \item Kompatibilitu přepisování a~kanonické veřejné adresy URL.
  \item Diagnostiku Stavu systému a~čekající migrace.
\end{itemize}

\subsection{Čisté adresy URL a~kompatibilita přepisování}
\index{p00177@přepisování URL}
\index{s00212@směrování pomocí parametrů dotazu}
\index{c00022@čisté adresy URL}

PHP Gallery podporuje čisté veřejné cesty, pokud hostingové prostředí používá pravidla přepisování projektu. Čisté adresy URL jsou volbou prezentace a~směrování, nikoli požadavkem datového modelu galerie. Trasy s~parametry dotazu zůstávají kompatibilním řešením a~fungují i~tehdy, když je přepisování vypnuté nebo je nelze spolehlivě ověřit.

Nastavení přepisování URL řídí podobu odkazů generovaných PHP Gallery. Když je přepisování vypnuté, pomocné funkce používají adresy s~parametry dotazu pro galerie, fotografie, náhledy a~další trasy aplikace místo předpokladu, že webový server přeloží čistou cestu. Při změně tohoto nastavení není nutné přejmenovávat identifikátory v~URL, složky galerií, databázové řádky ani zdrojové fotografie.

Pokud na novém hostingu čisté cesty nefungují, nejprve v~administraci vypněte přepisování URL a~ověřte, že funguje podoba s~parametry dotazu, například \codeword{?page=gallery}. Na Apache nebo LiteSpeed poté ověřte, že se uplatňuje projektový \filepath{.htaccess} a~že je dostupná podpora přepisování. Na jiných webových serverech ponechte režim kompatibility, dokud nenastavíte odpovídající pravidla směrování.

\section{První hodina s~PHP Gallery}
\label{sec:first-hour}
\index{p00167@první hodina}
\index{z00272@začínáme}
\index{v00256@vlastník galerie}
\index{p00170@přehled administrace}

Při prvním použití zvolte malý rozsah: vytvořte jednu galerii, přidejte několik fotografií a~otevřete výsledek v~prohlížeči, ve kterém nejste přihlášeni jako správce. Úpravy vzhledu a~podrobnější uspořádání mohou počkat, dokud nebude tento základní postup ověřený.

\subsection{Otevřete administraci}
\index{p00183@přihlášení správce}

Ve veřejném záhlaví zvolte \ui{Přihlášení správce} a~zadejte účet vytvořený při instalaci. Přehled je výchozím bodem každodenní správy. Shrnuje galerie, fotografie, stav systému a~dostupné činnosti údržby.\footnote{Přesný obsah přehledu závisí na povolených funkcích, dokončených migracích a~vývojových nastaveních. I~zdravá instalace může zobrazovat informační oznámení, která správce vybízejí k~dokončení volitelných úkolů.}

Věnujte několik minut seznámení s~těmito oblastmi:

\begin{itemize}
  \item \ui{Galerie}, kde se sbírky vytvářejí, vyhledávají, upravují a~uspořádávají;
  \item \ui{Vzhled}, kde se vybírá vzhled veřejného webu a~výchozí volby prezentace;
  \item \ui{Štítky}, kde se udržují opakovaně využitelná témata a~kategorie;
  \item \ui{Stav systému}, kde se kontrolují migrace, náhledy, úložiště a~konzistence;
  \item \ui{Účet}, kde se spravují přihlašovací údaje, e-mail pro obnovu a~doručování pošty.
\end{itemize}

\subsection{Vytvořte první galerii}
\index{g00040@galerie!v00266@vytvoření první galerie}

Otevřete správu galerií a~vytvořte galerii s~krátkým názvem, například \emph{Letní procházka}. Zkontrolujte navrženou veřejnou cestu a~název složky a~poté přidejte jednu větu popisující obsah sbírky.

Pro tuto první kontrolu nastavte galerii jako veřejnou a~ponechte pokročilé vlastní značky nebo přepsání zděděných voleb na výchozích hodnotách. Přidejte tři až deset fotografií. To stačí k~ověření pořadí, generování náhledů, názvů a~lightboxu, aniž by se cvičení změnilo v~rozsáhlý import.

\subsection{Náhled očima návštěvníka}
\index{a00006@anonymní náhled}
\index{n00093@náhled pro návštěvníka}

Použijte odkaz na veřejnou galerii nebo ovládací prvek náhledu pro návštěvníka. Pro realistickou kontrolu otevřete veřejnou adresu v~soukromém okně prohlížeče bez relace správce. Ověřte, že se název, popis, fotografie a~navigace zobrazují podle záměru. Otevřete fotografii, přejděte v~prohlížeči dopředu i~zpět a~vraťte se do galerie.

Zobrazení pro správce může obsahovat ovládací prvky úprav a~další podklady. Anonymní náhled nejlépe ukazuje prostředí, které uvidí rodina, klienti nebo veřejnost.\footnote{Chráněné galerie a~obsah s~věkovým omezením vyžadují samostatnou anonymní zkoušku, protože aktivní relace správce již poskytuje širší přístup.}

\subsection{Dokončete kontrolní seznam první hodiny}

\begin{enumerate}
  \item Nastavte veřejný název webu a~jazyk.
  \item Vytvořte jednu galerii s~názvem a~popisem.
  \item Nahrajte malou skupinu fotografií.
  \item Přidejte smysluplný název alespoň k~jedné fotografii.
  \item Změňte pořadí dvou fotografií a~ověřte nové pořadí na veřejném webu.
  \item Otevřete galerii jako anonymní návštěvník.
  \item Zkontrolujte galerii na obrazovce velikosti telefonu.
  \item Ověřte, že Stav systému nehlásí naléhavý problém s~migracemi nebo oprávněními.
  \item Po dokončení počátečního nastavení vytvořte souběžnou zálohu všech souvisejících částí.
\end{enumerate}

\section{Plánování sbírky galerií}
\label{sec:planning}
\index{p00137@plánování galerií}
\index{n00103@návrh sbírky}
\index{h00046@hierarchie galerií!p00136@plánování}

Promyšlená struktura pomáhá návštěvníkům pochopit sbírku a~vlastníkovi ji dlouhodobě udržovat. Začněte tím, jak budou lidé fotografie hledat. Názvy složek a~podrobnosti databáze mohou z~tohoto přirozeného uspořádání vycházet.

\subsection{Zvolte srozumitelnou strukturu}

Běžné struktury zahrnují:

\begin{description}
  \item[Podle roku a~události] Rok na nejvyšší úrovni obsahuje svatby, výlety, oslavy a~projekty. Hodí se pro rodinné a~historické archivy.
  \item[Podle místa] Země obsahují oblasti a~města. Hodí se pro cestovatelské, letecké, krajinářské a~dokumentární fotografie.
  \item[Podle klienta nebo projektu] Každý klient obsahuje zakázky a~galerie pro předání. Hodí se pro profesionální práci s~oddělenými pravidly přístupu.
  \item[Podle tématu] Příroda, architektura, portréty, letadla a~podobná témata tvoří trvalé kategorie. Štítky mohou poskytnout druhý způsob jejich vzájemného procházení.
  \item[Podle pracovního postupu] Galerie příchozích, vybraných, předaných a~archivních fotografií představují jednotlivé fáze. Nastavení viditelnosti ponechává pracovní sbírky soukromé, zatímco hotové výběry jsou veřejné.
\end{description}

Galerie používejte pro výrazné hranice navigace a~štítky pro vztahy, které tyto hranice překračují. Fotografie z~Prahy může být v~\emph{Cestování / Česká republika / Praha} a~nést štítky \emph{architektura}, \emph{noc} a~\emph{tramvaj}.\footnote{V~běžném modelu obsahu má fotografie jedno umístění v~galerii, zatímco opakovaně použitelné štítky poskytují několik významových vazeb. Nástroje pro kopírování a~přesouvání podporují záměrné změny fyzického uspořádání.}

\subsection{Rozhodněte o~hloubce hierarchie}
\index{h00046@hierarchie galerií!h00049@hloubka}

Vnořené galerie podporují hluboké sbírky, ale většina návštěvníků se lépe orientuje, když má každá úroveň jasný účel. Mnoha webům dobře poslouží dvě až čtyři úrovně. Hlubší archiv zůstává praktický, pokud čtenáře vedou drobečková navigace, smysluplné názvy a~jednotná data.

Vytvořte podřízenou galerii, pokud si zaslouží vlastní název, popis, titulní obrázek, pravidla přístupu, rozsah dat nebo veřejnou adresu. Ponechte fotografie pohromadě, pokud by nová úroveň přidala pouze jedno kliknutí navíc.

\subsection{Připravte názvy a~popisy}
\index{n00106@název galerie}
\index{p00141@popis galerie}

Název galerie by měl být srozumitelný i~mimo administraci. Popis může odpovědět na tři otázky:

\begin{itemize}
  \item Co tato sbírka zobrazuje?
  \item Kde a~kdy vznikla?
  \item Jaké souvislosti dávají fotografiím význam?
\end{itemize}

Popisy mohou být stručné u~běžných alb a~rozsáhlé u~dokumentární tvorby. První větu použijte jako zásadní shrnutí, protože návštěvníci často nejprve zběžně prohlížejí a~teprve potom čtou celý text.

\subsection{Naplánujte soukromí před nahráváním}
\index{p00138@plánování soukromí}
\index{v00251@viditelnost galerie}

Rozhodněte, zda je sbírka veřejná, nezveřejněná pro přípravu ve správě, nebo chráněná pro známé publikum. Pravidla nadřazené galerie mohou ovlivňovat potomky, proto pokud možno umísťujte sbírky pro podobné publikum společně.

Příklady:

\begin{itemize}
  \item veřejné portfolio s~pečlivě vybranými fotografiemi;
  \item rodinná větev galerií chráněná heslem;
  \item nezveřejněná galerie pro klientský výběr během přípravy;
  \item veřejná galerie z~akce s~vybranými jednotlivými fotografiemi označenými jako soukromé;
  \item větev galerií s~ověřením věku pro obsah vyžadující potvrzení návštěvníka.
\end{itemize}
\section{Vytváření a~úpravy galerií}
\label{sec:create-edit-gallery}
\index{g00040@galerie!v00264@vytváření}
\index{g00040@galerie!u00244@úpravy}
\index{e00035@editor galerie}

\subsection{Vytvoření galerie v~administraci}

Ovládacím prvkem \ui{+} ve veřejné galerii nebo volbou \ui{Vytvořit galerii zde} v~jejím editoru zahajte vytvoření podřízené galerie. Pravý postranní panel administrace požaduje pouze nový název; nadřazenou galerii, z~níž byl ovládací prvek vyvolán, již zná. Nová galerie začíná jako \ui{Nezveřejněná}. Po vytvoření tentýž panel otevře úplný editor popisu, data, jazyka, štítků, přístupu, zobrazení, médií a~fotografií. Během tohoto rozšířeného postupu veřejná stránka nikam nepřechází.

Na samostatné stránce vytváření galerie v~administraci zadejte název a~před odesláním případně doplňte údaje SimBrief, popis, zdrojový jazyk a~štítky. Sbalený oddíl \ui{Další možnosti} obsahuje název složky, počáteční viditelnost, rozsah dat, nadřazenou galerii, hlasování, zobrazení názvů souborů a~odznak počtu obsažených fotografií. Krátké shrnutí pod ním ukazuje aktuální viditelnost a~nadřazenou galerii. Samostatný formulář pro vytvoření a~nahrání zůstává dostupný, pokud chcete fotografie nahrát již při vytváření. Při chybě ověření zůstávají zadané neutajované hodnoty i~kontext nadřazené galerie k~dispozici pro opravu.

Jakmile do pole názvu zadáte alespoň dva znaky, administrace může zbývající text odpovídající existující galerie zobrazit jako nenápadný návrh přímo v~poli. Přednost mají nedávné názvy pod vybranou nadřazenou galerií; shody jinde se zvažují až po vyčerpání omezeného hledání mezi sourozenci. Starší názvy mohou být vynechány, takže nepřítomnost návrhu nedokazuje jedinečnost názvu. Formulář požaduje nejvýše osm kandidátů místo načtení úplného katalogu názvů.

Je-li kurzor na konci a~není vybrán žádný text, přijměte návrh stisknutím samotné klávesy \codeword{Tab} nebo \ui{Šipka doprava}, případně výběrem navrženého textu. Klávesové zkratky s~modifikátory a~skládání znaků vstupní metodou si zachovávají běžné chování. První stisknutí \codeword{Escape} zavře doplnění; další stisknutí se může dostat k~okolnímu panelu. Přeložená nápověda a~nenaléhavé oznámení zpřístupňují návrhy bez změny názvu pole.

Porovnávání Unicode přijímá pouze bezpečné hranice celých znaků. Některé návrhy mimo ASCII se vynechají, pokud nejsou dostupná serverová rozšíření nebo podpora segmentace v~prohlížeči. Pomalé, neúspěšné nebo zastaralé vyhledávání nikdy nebrání běžnému psaní. Přijetí mění pouze pole názvu: neodesílá formulář, nepřejmenovává existující galerii ani nemění její pravidla přístupu. Stejné chování platí i~po nahrazení formuláře v~postranním panelu. Omezení porovnávání a~technické podrobnosti najdete v~\filepath{docs/TITLE_COMPLETION.md}.

Po vytvoření pouze zadáním názvu je galerie již skutečnou sbírkou ve stromu galerií, i~před uložením dalších polí editoru. Její fyzická složka může přijímat fotografie přes prohlížeč, přenosem v~souborovém systému, vyhledáním souborů nebo jiným podporovaným postupem nahrávání. Před zveřejněním vyplňte a~uložte pole editoru. Databáze uchovává popisná a~prezentační nastavení; o~tom, co návštěvníci uvidí, rozhodují běžná pravidla přístupu.

Každý vygenerovaný formulář vytváření galerie nebo klasického nahrávání nese také soukromou identitu operace. Pokud spojení selže poté, co server mohl práci dokončit, zkuste znovu tentýž formulář místo otevření nového: stejný správce, operace a~obsah obdrží původní dokončený výsledek. Nově otevřený formulář představuje novou záměrnou operaci a~nadále může vytvořit úmyslnou duplicitu. Tato identita není heslem ani náhradou přihlášení a~ochrany CSRF.

\subsubsection{Postup vytváření}

\begin{enumerate}
  \item Otevřete zamýšlenou nadřazenou galerii a~zvolte její ovládací prvek \ui{+}, nebo začněte v~administraci vytvářet kořenovou galerii.
  \item Zadejte název srozumitelný návštěvníkům; případně přijměte návrh názvu.
  \item Vytvořte nezveřejněnou galerii. Je-li dostupný JavaScript, úplný editor se otevře ve stejném panelu.
  \item V~oddílu \ui{Identita} přidejte popis, datum nebo rozsah dat, zdrojový jazyk, štítky a~volitelné údaje SimBrief.
  \item Zkontrolujte \ui{Přístup}, \ui{Zobrazení}, \ui{Média} a~podle potřeby také hodnoty složky a~nadřazené galerie v~pokročilých volbách.
  \item Uložte galerii, přidejte fotografie a~po vložení vhodných fotografií zvolte titulní obrázek.
  \item Až budete připraveni, galerii zveřejněte a~poté anonymně ověřte výsledek.
\end{enumerate}

\subsubsection{Zapamatované výchozí hodnoty a~návrhy SimBrief}
\index{v00265@vytváření galerie!z00276@zapamatované výchozí hodnoty}
\index{s00207@SimBrief!n00102@návrh před vytvořením}

Editor si může pro přihlášeného správce zapamatovat aktuální identifikátor SimBrief a~zdrojový jazyk. Sousední volbu \ui{Zapamatovat pro budoucí galerie} zaškrtněte pouze u~hodnoty, kterou chcete používat opakovaně. Pozdější galerie zobrazí tuto hodnotu již vybranou s~krátkou poznámkou, že pochází z~uložených výchozích hodnot. Hodnotu lze nadále upravit; ponechání volby nezaškrtnuté zachová předchozí uloženou volbu. Výchozí hodnoty náleží jednomu správci a~nezveřejňují přihlašovací údaje ani nekopírují pravidla přístupu. Aby bylo možné uložit novou výchozí hodnotu, musí běžná migrační akce použít \codeword{202609250001\_gallery\_creation\_preferences.php}; vytváření galerií funguje i~bez této volitelné tabulky.

Viditelné pole SimBrief přijímá ID pilota nebo jeho jméno. Samotné číslice se považují za ID pilota, ostatní text za jméno. V~přímém formuláři vytváření může volba \ui{Vygenerovat návrh popisu} načíst nejnovější OFP ještě před vznikem galerie. Před uložením navržený popis zkontrolujte a~upravte. Import uchovává soukromý odkaz po dobu 30~minut, omezený na stejného správce a~relaci; změna identifikátoru tento odkaz zneplatní. OFP a~dostupné souřadnice trasy se připojí až po vytvoření galerie. Pokud po vytvoření selže volitelné připojení nebo uložení výchozí hodnoty, galerie zůstane vytvořena a~panel oznámí konkrétní upozornění. Při vypršení návrhu nebo změně relace proveďte import znovu.

\subsubsection{Práce v~kompaktním editoru}
\index{e00035@editor galerie!p00146@postranní panel}

Běžné pracovní oblasti jsou \ui{Identita}, \ui{API}, \ui{Přístup}, \ui{Zobrazení} a~\ui{Média}; podle dostupnosti následují nástroje pro obrázky a~specializované úlohy. Pravý postranní panel ponechává tlačítko \ui{Uložit galerii} viditelné dole, zatímco formulář nastavení se posouvá. Změna karty automaticky neukládá. Po úpravách použijte tlačítko a~chcete-li ověřit uložené hodnoty, galerii znovu otevřete. Delší vysvětlení jsou dostupná přes ovládací prvky s~otazníkem přístupné z~klávesnice a~méně časté volby jsou pod rozbalovacími prvky \ui{Pokročilé}. Přímý editor v~administraci používá stejná pole; prohlížeč bez JavaScriptu odesílá formulář běžným způsobem.

\ui{Identita} ponechává snadno dostupný název, popis, datum, zdrojový jazyk, SimBrief a~štítky. Její pokročilý oddíl obsahuje identifikátor v~URL, složku, nadřazenou galerii, pořadí a~text galerie od AI, je-li povolen. \ui{API} zobrazuje kopírovatelný koncový bod nahrávání, označení a~akci klíče a~nově vygenerovaný klíč pouze jednou; existující klíče se zobrazují, jen pokud nějaké jsou. Obnova metadat pomocí AI, migrace a~globální správce API jsou pod \ui{Pokročilé nástroje API}. \ui{Přístup} seskupuje viditelnost, heslo, ovládání odkazu pro sdílení a~zaškrtávací pole NSFW; delší vysvětlení pravidel jsou pod otazníky. \ui{Zobrazení} začíná mřížkou a~běžnými přepínači hlasování a~názvů souborů; EXIF/GPS, formát karty galerie, odznak počtu, režim lightboxu, hra s~fotografiemi, meze náhledů a~text letové mapy jsou pod \ui{Pokročilá nastavení zobrazení}.

\subsection{Úprava identity a~souvislostí}
\index{m00084@metadata galerie}
\index{r00197@rozsah dat}

Editor umožňuje galerii průběžně zlepšovat. Názvy a~popisy se mohou měnit s~vývojem sbírky. Pole dat mohou představovat jeden den, cestu, roční období nebo delší historické období. Návrhy dat z~EXIF mohou prozkoumat fotografie a~potomky a~poté nabídnout upravitelný rozsah ke schválení.

Úplný editor zaznamenává přesnou revizi galerie, ze které pocházejí jeho pole. Pokud galerii nejprve změní jiná karta, správce, nahrávání, přesun, skenování nebo údržba, starší uložení se odmítne, aby tiše nepřepsalo novější stav. Zobrazení konfliktu uchová zadaný neutajovaný text a~ukáže bezpečné porovnání s~aktuálním stavem. Zkontrolujte a~zkopírujte zamýšlené změny do nejnovějšího editoru; hesla zadejte znovu a~soubory znovu vyberte, protože soukromé hodnoty se do podrobností konfliktu nikdy nekopírují.

Návrhy dat používejte jako užitečný výchozí bod. Skenované fotografie, snímky obrazovky, upravené exporty a~fotoaparáty s~nesprávně nastavenými hodinami mohou vytvářet data vyžadující lidskou kontrolu.\footnote{Použití návrhu odvozeného z~EXIF aktualizuje navržená pole galerie prostřednictvím stávajícího editoru. Za konečná data zobrazená návštěvníkům nadále odpovídá správce.}

\subsubsection{Kontrola návrhů dat napříč galeriemi}
\index{d00023@data galerií!k00066@kontrola EXIF}

Samostatná stránka údržby \ui{Data galerií} používá stejné podklady EXIF napříč mnoha galeriemi. Může zkontrolovat celou instalaci nebo jednu vybranou větev galerií. Každý řádek ukazuje aktuální ruční rozsah dat, navržený rozsah, počet fotografií EXIF a~uzlů galerií, které poskytly podklady, a~upravitelná pole \ui{Od}/\ui{Do}. Řádky bez existujícího ručního data jsou ve výchozím stavu vybrány k~použití; řádky s~již zadaným ručním datem zůstávají nezaškrtnuté, aby hromadná kontrola tiše nepřepsala chronologii zadanou vlastníkem.

Ukládají se pouze zaškrtnuté řádky. Nezaškrtnuté návrhy se ignorují a~správce může před použitím výběru upravit obě data. Pokud starší importované fotografie dosud nemají načtená data pořízení EXIF, nejprve spusťte běžné skenování/import obrázků galerie. Stejná cílená akce návrhu je dostupná v~editoru jednotlivé galerie a~může jej aktualizovat přímo na místě.

Pole dat v~administraci ponechávají nativní datumové pole prohlížeče jako odesílanou hodnotu, ale přidávají kompaktní tlačítko kalendáře a~pomocné volby \ui{Dnes} a~\ui{Vymazat}. \ui{Dnes} vloží aktuální místní kalendářní datum prohlížeče; \ui{Vymazat} pole vyprázdní. Toto vylepšení se znovu použije na formuláře načtené do postranního panelu administrace, zatímco základní nativní datumové pole zůstává použitelné i~bez dodatečných ovládacích prvků JavaScriptu.

\subsection{Přehled ovládacích prvků editoru galerie}
\index{e00035@editor galerie!o00131@ovládací prvky}
\index{g00040@galerie!v00250@viditelnost}
\index{g00040@galerie!p00178@přepsání zobrazení}

Editor galerie je rozhodujícím místem pro nastavení náležející jedné fyzické galerii. Běžné karty jsou Identita, API, Přístup, Zobrazení a~Média; následují Obrázky a~volitelné nástroje Organizátor/Přejmenování, pokud jsou dostupné. Rozbalovací pokročilé oddíly uchovávají méně častá nastavení ve společném formuláři ukládání. Následující tabulka shrnuje aktuální ovládací prvky jednotlivých galerií ve verzi~\version{}.

\begin{longtable}{>{\RaggedRight\arraybackslash}>{\hspace{0pt}}p{0.21\textwidth} >{\RaggedRight\arraybackslash}>{\hspace{0pt}}p{0.27\textwidth} >{\RaggedRight\arraybackslash}>{\hspace{0pt}}p{0.43\textwidth}}
\toprule
\textbf{Oblast} & \textbf{Ovládací prvek} & \textbf{Současné chování}\\
\midrule
\endfirsthead
\toprule
\textbf{Oblast} & \textbf{Ovládací prvek} & \textbf{Současné chování}\\
\midrule
\endhead
Identita & Název a~popis & Veřejný redakční text. Udržované překlady lze uložit odděleně bez změny zdrojového textu.\\
Identita & Ruční datum a~rozsah dat & Volitelné datum události nebo fotografování a~podporovaný rozsah od/do. Návrhy odvozené z~EXIF lze před uložením upravit.\\
Identita & Identifikátor v~URL & Řídí identitu veřejné adresy URL galerie. Pokud možno zachovejte zavedené identifikátory, protože na nich mohou záviset záložky a~externí odkazy.\\
Identita & Název složky & Přejmenuje fyzickou složku galerie na disku prostřednictvím běžného postupu změn. Nejde pouze o~zobrazený popisek.\\
Identita & Nadřazená galerie a~pořadí & Mění hierarchii a~pořadí sourozenců. Přesun galerie fyzicky přemístí strom jejích složek, pokud to nakonfigurovaný kořen galerií dovoluje.\\
Identita & Štítky & Opakovaně použitelné štítky oddělené čárkami. Návrhy se váží podle okolních galerií, fotografií a~kontextu složek, nejde o~prostý náhodný seznam.\\
Identita & Identifikátor SimBrief a~výchozí hodnoty & Jediné viditelné pole přijímá číselné ID pilota nebo jeho jméno. Správce může identifikátor a~zdrojový jazyk zapamatovat pro budoucí galerie; předvyplněné hodnoty jsou označené.\\
API & Koncový bod nahrávání a~klíče & Zkopírujte koncový bod omezený na galerii, vygenerujte klíč s~označením a~ihned zkopírujte jeho novou nezpracovanou hodnotu. Aktivní klíče se zobrazí pouze tehdy, pokud existují; nezpracovaná hodnota se znovu nezobrazuje. Obnova pomocí AI a~migrace jsou v~pokročilých volbách.\\
Přístup & Viditelnost & \ui{Veřejná} se běžně uvádí v~seznamech. \ui{Nezveřejněná} je skrytá z~běžných seznamů, ale lze ji otevřít z~normální adresy URL, pokud to pravidla přístupu dovolují. \ui{Soukromá} je dostupná pouze správci s~výjimkou podporovaného přímého přístupu tokenem.\\
Přístup & Uzamčení heslem & Nezávisí na viditelnosti. Veřejná, nezveřejněná i~soukromá galerie může mít vlastní pravidla hesla, je-li dostupné schéma přístupu. Prázdné pole nového hesla zachovává stávající heslo.\\
Přístup & Odkaz pro sdílení & Lze jej vytvořit, znovu vygenerovat, zrušit a~volitelně mu nastavit dobu platnosti. Uložené údaje tokenu chrání hash; starší token nebo token, který nelze dešifrovat, může vyžadovat opětovné vygenerování, aby bylo možné znovu zobrazit kopírovatelný odkaz.\\
Přístup & NSFW / 18+ & Označuje větev galerií pro anonymní potvrzení věku. Ochrana se vztahuje i~na podgalerie a~poskytovaná média podle účinných pravidel přístupu.\\
Zobrazení & Hra s~fotografiemi & Pokročilá volba porovnávací hry v~této větvi galerií. Povolení hry vyžaduje také dostupné hlasování o~obrázcích.\\
Zobrazení & Hlasování o~obrázcích & Povolí veřejná označení oblíbených obrázků v~této galerii. Vypnutí uchová uložené skóre, ale odstraní ovládací prvky nových hlasů a~zablokuje nové příspěvky.\\
Zobrazení & Zobrazovat názvy souborů & Ve výchozím stavu vypnuto. Původní názvy nahraných souborů zůstávají skryté, dokud není volba povolena; vlastní názvy a~popisy fotografií se nadále zobrazují.\\
Zobrazení & Zobrazení EXIF/GPS & \ui{Použít globální výchozí hodnotu}, \ui{Vynutit zapnutí} nebo \ui{Vynutit vypnutí}. Účinný stav řídí veřejné značky GPS, mapy EXIF galerie a~zobrazení souřadnic GPS.\\
Zobrazení & Formát popisu karty galerie & Ve výchozím stavu dědí Vzhled a~tam, kde je to podporováno, může přepsat rozvržení karty této galerie.\\
Zobrazení & Odznak počtu fotografií & Dědí Vzhled nebo přepisuje ikonu naskládaných fotografií a~počet obsažených obrázků zobrazovaný na kartách galerií.\\
Zobrazení & Režim procházení lightboxu & Dědění ze Vzhledu, klasický jednotlivý obrázek, pás fotografií nebo 3D karusel. Celoobrazovkové zobrazení a~prezentace zůstávají dostupné ve všech režimech.\\
Zobrazení & Vlastní mřížka & První viditelný ovládací prvek Zobrazení: volitelné sloupce a~řádky této galerie. Bez výslovné mřížky se použije nejbližší mřížka nadřazené galerie povolující dědění do podgalerií a~poté Vzhled.\\
Zobrazení & Meze responzivních náhledů & Volitelné minimální a~maximální meze velikosti náhledů. Uložení může meze galerie rekurzivně zkopírovat potomkům; rekurzivní uložení je ve výchozím stavu vypnuté.\\
Zobrazení & Letová trasa & Pokročilý ruční text trasy podporuje místní vyhledávání identifikátorů a~výslovné body \codeword{NAME@latitude,longitude}. SimBrief může uložit OFP, návrh popisu a~vyřešené souřadnice trasy.\\
Média & Titulní fotografie & Automatická nebo vybraná fotografie, včetně způsobilých fotografií potomků, pokud je výběr nabízí.\\
Média & Nahraný náhled galerie & Samostatné médium karty galerie uložené nezávisle na běžných fotografiích galerie.\\
Média & Zdroj pozadí & Dědění ze Vzhledu, nahrání nového pozadí, výběr existujícího obrázku galerie nebo vytvoření koláže z~veřejných galerií.\\
Média & Podklady vlastní značky & Volitelné podklady banneru, loga a~oddělovače pro konkrétní galerii přepisují náhradní hodnoty Vzhledu pouze v~té galerii, kde jsou nastavené.\\
Obrázky & Skenování/import & Znovu prohledá fyzickou složku galerie a~synchronizuje podporovaná média do katalogu.\\
Obrázky & Vytvořit všechny náhledy & Znovu vytvoří náhledy této galerie. Chcete-li zahrnout potomky, vyberte \ui{Zahrnout podgalerie}; ostatní větve galerií zůstávají mimo úlohu. Je-li dostupné generování v~prohlížeči, je vybrané jako výchozí, s~možností generování na serveru.\\
Nástroje & Organizátor, Přejmenování, API & Organizátor podle data, fyzické přejmenování médií, správa klíčů API pro nahrávání, obnova pomocí AI a~migrace galerie se zobrazují pouze tehdy, jsou-li dostupná jejich schémata a~funkce.\\
\bottomrule
\end{longtable}

\important{Ovládací prvky konkrétních galerií záměrně nejsou sloučeny do globálního centra Nastavení. Dědění je součástí datového modelu, takže výsledek globálního vyhledávání může odkazovat na editor galerie, aniž by hodnotu náležející galerii změnil na nastavení celého webu.}

\subsection{Výběr titulního obrázku a~vizuální identity}
\index{t00235@titulní obrázek galerie}
\index{v00254@vlastní značka galerie}

Titulní obrázek je vizuální pozvánkou do galerie. Vyberte obrázek, který zůstává rozpoznatelný jako náhled a~věrně představuje sbírku. Obličeje, výrazné siluety, jasné hlavní motivy a~nepřeplněné kompozice často dobře fungují i~v~malých velikostech.

Vlastní značka galerie může podle dostupných ovládacích prvků poskytnout logo, banner, pozadí, oddělovač nebo přepsání prezentace. Začněte zděděným stylem vzhledu a~poté přidejte podklady konkrétní galerie tam, kde má sbírka skutečnou vlastní identitu, například svatba, výstava, klient nebo dlouhodobý projekt.

\subsection{Přesouvání a~změna pořadí galerií}
\index{g00040@galerie!p00179@přesouvání}
\index{g00040@galerie!z00280@změna pořadí}

Přesun galerie mění její místo v~hierarchii a~může ovlivnit zděděný přístup nebo prezentaci. Před potvrzením přesunu zkontrolujte cíl. Změna pořadí mění posloupnost sourozenců a~umožňuje sestavit promyšlené pořadí prohlížení.

Přesouvání fotografií mezi galeriemi používá trvalý žurnál operací. Organizátor fyzicky přesune každý originál a~před změnou vlastnictví v~katalogu ověří jeho zaznamenanou identitu. Pokud pracovní proces nebo spojení skončí během operace, Stav systému a~Diagnostika běhového prostředí zobrazí omezený stav vyžadující akci, aniž by odhalily soukromé cesty. Neúspěšný přesun uvede na obrazovce bezpečný důvod; důvěryhodné protokoly administrace uchovávají podrobnější soukromou diagnostiku. Zastavte souběžnou práci s~úložištěm, zachovejte zálohu a~postupujte podle \filepath{docs/IMAGE_MOVE_RECOVERY.md}; obnova změněné nebo neověřitelné soubory uchovává místo jejich přepsání.

Pořadí může být chronologické, geografické, vyprávěcí, abecední nebo ručně sestavené. V~rámci jedné nadřazené galerie zachovejte dostatečně jednotnou volbu, aby návštěvníci mohli předvídat další položku.

Přihlášený správce může také měnit pořadí karet přímých podgalerií na veřejné stránce galerie. Kompaktní veřejný panel změny pořadí je omezen na karty viditelné na aktuální stránce stránkování. Požadavek uložení obsahuje posun této stránky a~počet viditelných položek; server změnu odmítne, pokud odeslaná ID již neodpovídají stejnému aktuálnímu výřezu. Mění pouze hodnoty \codeword{sort_order} sourozenců; nemění nadřazené galerie ani nevnořuje galerie.

Přihlášený správce prohlížející veřejnou galerii může také dočasně seřadit přímé podřízené galerie s~vyplněným datem \ui{Od} od nejstarších po nejnovější nebo opačně. Podřízené galerie bez použitelného data zůstávají na svých místech. Náhled je dočasný, dokud nezvolíte \ui{Uložit}; uložení zapíše výsledek jako skutečné pořadí sourozenců zobrazované všem návštěvníkům. Pro uložení pořadí podle data jsou nutné alespoň dvě přímé podřízené galerie s~datem.

\subsection{Promyšlené odstranění galerie}
\index{g00040@galerie!o00121@odstranění}

Odstranění galerie může ovlivnit zdrojové soubory, potomky, metadata, náhledy, odkazy a~záložky návštěvníků. Zkontrolujte přesně vybranou galerii a~její podgalerie, vytvořte zálohu a~použijte běžný postup odstranění v~administraci. Dočasná změna viditelnosti může poskytnout bezpečnější období přípravy, pokud galerii ještě můžete potřebovat.

Koš galerií je ve výchozím stavu povolen. Odstranění galerie běžnou akcí správce přesune vybranou složku a~celý strom potomků mimo aktivní kořen galerií, zaznamená metadata potřebná pro obnovu a~uvede položku v~\ui{Údržba > Koš}. Veřejné stránky přestanou tuto podvětev okamžitě vyhledávat, ale správce ji může obnovit do původní cesty, pokud je stále volná. Obnovení nikdy nepřepíše novou aktivní galerii na stejném místě.

V~panelu Koše zvolte \ui{Obnovit}, \ui{Trvale odstranit} nebo \ui{Vyprázdnit Koš}. Trvalé odstranění a~vyprázdnění Koše nelze vrátit. Vyprázdnění Koše probíhá v~omezených dávkách, takže velké množství položek nevyžaduje jediný neomezený požadavek. Odstraňování jednotlivých fotografií není součástí této obnovy na úrovni galerií a~zůstává trvalé. Záloha je nadále nezbytnou ochranou proti selhání úložiště, trvalému vymazání správcem nebo poškození mimo PHP Gallery.
\section{Chytré galerie}
\index{h00053@chytré galerie}
\index{v00252@virtuální galerie}

Chytrá galerie je uložený dotaz, nikoli složka. Otevřete \ui{Administrace > Galerie > Chytré galerie}, přidejte podmínky a~kombinujte je pomocí vnořených skupin \ui{AND}, \ui{OR} a~\ui{NOT}. Pole zahrnují zdrojové galerie a~potomky, štítky, data pořízení, EXIF/GPS, text, metadata AI, stav duplicit, vlastnosti souborů a~soukromé redakční hodnocení.

Pomocí \ui{Náhled} zobrazte aktuální počet shod a~skutečné karty odpovídajících fotografií, poté galerii uložte a~případně zveřejněte. Obrázky i~soubory zůstávají ve svých fyzických galeriích, takže změny metadat nebo hodnocení okamžitě mění jejich příslušnost. Redakční hodnocení je soukromé a~nezávislé na hlasování návštěvníků. Správci mohou také ve veřejném vyhledávání použít \ui{Uložit vyhledávání jako chytrou galerii}.

Zvolte \ui{Neuvedená v~seznamech} pro přístup pouze přes URL, \ui{Kořenová galerie} pro zobrazení chytré galerie na domovské stránce nebo \ui{Umístění jako podgalerie}, aby ji mohly připojit fyzické galerie. V~editoru každé fyzické galerie vyberte, zda připojení patří \ui{Nad obsah galerie} nebo \ui{Pod obsah galerie}, a~přiřaďte mu hodnotu pořadí. Nejprve se vykreslí připojení nad obsahem, poté běžné podgalerie a~fotografie a~nakonec připojení pod obsahem. Shodné hodnoty pořadí se jednotně rozlišují podle ID chytré galerie. Existující připojení mají ve výchozím stavu umístění pod obsahem a~pořadí~0.

Stejnou virtuální galerii lze připojit pod několik fyzických galerií, s~nezávislým umístěním a~pořadím u~každého rodiče. Oddíl \ui{Použito ve fyzických galeriích} v~editoru chytré galerie ukazuje všechny rodiče a~umožňuje uložit jejich umístění a~pořadí nebo použít \ui{Odpojit}, aniž se změní ostatní rodiče. Je-li editor otevřený v~postranním panelu administrace, tyto akce zůstávají v~panelu; při vypnutém JavaScriptu je dostupné běžné odeslání formuláře. Veřejná chytrá galerie umístěná pod soukromou nebo jinak omezenou nadřazenou galerií nadále pro svou kartu dědí hranici jejího veřejného přístupu.

Veřejné výsledky a~počty vždy uplatňují stejná pravidla přístupu ke zdrojům jako běžný veřejný obsah galerií. Odpovídající fotografie se započítá pouze tehdy, když je její fyzická galerie veřejně uvedená a~přístupná aktuálnímu návštěvníkovi a~fotografie samotná je veřejná. Stávající pravidla hesel, odkazů pro sdílení, zděděné viditelnosti, nezveřejněných galerií a~věkových omezení tedy nadále platí. Trasy s~parametry dotazu fungují bez přepisování a~čisté adresy URL \texttt{/smart/<slug>} jsou volitelné.

\subsection{Bezpečnost vztahů a~opravy}
\index{h00053@chytré galerie!o00125@ochrana před cykly}

Chytrá galerie se nesmí vyhodnotit zpět do fyzické galerie, která ji připojuje. Server při každé změně pravidel, připojení nebo hierarchie galerií kontroluje úplný graf vztahů; filtrování v~prohlížeči se nepovažuje za bezpečnostní opatření. To zahrnuje nepřímé cesty přes fyzické potomky a~jiné připojené chytré galerie. Synchronizace hierarchie odvozené ze souborového systému a~oprava rodičů podle veřejných cest ověřují úplnou navrženou mapu rodičů před zápisem jakéhokoli nového rodičovského vztahu. Současný editor pravidel může odkazovat na fyzické galerie, ale nenabízí přímou podmínku odkazu na chytrou galerii.

Úplný strom administrace upravený přetažením se ověří jednou ve své konečné podobě. Následné přesuny složek používají tuto předem ověřenou mapu a~odkládají synchronizaci hierarchie, dokud nevznikne konečný strom, takže dočasný mezistav nemůže odmítnout jinak platnou dávku. Každá potvrzená změna fyzické hierarchie vymaže mezipaměť grafu chytrých galerií platnou pro aktuální požadavek před pokračováním další práce závislé na grafu.

Pokud starší nebo ručně upravená databázová data již obsahují cyklus, veřejné vykreslování tento neplatný vztah zastaví místo rekurzivního pokračování. Seznam připojení v~administraci označí vztah k~opravě a~volba \ui{Odpojit} zůstane dostupná. Ověřování i~průchod za běhu používají konzervativní limity hloubky a~počtu uzlů a~odstraňování opakování v~rámci požadavku, takže chybná data nemohou opakovaně rozvíjet stejnou chytrou galerii. Běžný dotaz výsledků zůstává stránkovaný a~nekopíruje všechny odpovídající fotografie do PHP jen kvůli detekci cyklů.

Migrace \texttt{202608170002\_smart\_gallery\_attachment\_ordering.php} přidává pole umístění a~pořadí pro každého rodiče. Dokud tato migrace není dostupná, stávající připojení zachovávají historické umístění pod obsahem, zatímco změny umístění v~administraci jsou odmítnuty místo částečného uložení.

\subsection{Nastavení prezentace}
\index{h00053@chytré galerie!p00159@prezentace}

Chytrá galerie ve výchozím stavu používá aktuální Vzhled a~prezentaci webu. Povolte \ui{Přepsat výchozí nastavení Vzhledu pro tuto chytrou galerii}, pokud tato virtuální sbírka potřebuje vlastní sloupce, řádky, stránkování, meze velikosti náhledů, vykreslovač náhledů, rozvržení karty galerie, překryvy metadat, chování lightboxu, prezentaci, stahování nebo hlasování. Řazení a~směr zůstávají součástí definice chytré galerie.

Chytrá galerie nedědí prezentaci z~fyzické galerie, v~níž se zobrazuje její karta. Jedna chytrá galerie se může objevit pod několika fyzickými galeriemi, takže neexistuje jediná nadřazená prezentace s~jasnou předností. Pokud přepsání neexistuje nebo uložená data prezentace chybí či jsou neplatná, náhradou zůstávají aktuální výchozí hodnoty Vzhledu a~webu. Meze náhledů fyzické galerie a~fotografie mají nadále přednost, pokud s~nimi mez chytré galerie koliduje.

\subsection{Lightbox napříč stránkami výsledků}
\index{h00053@chytré galerie!l00075@lightbox}

Velký prohlížeč není omezen na fotografie vykreslené na aktuální stránce HTML. Karty chytré galerie si pamatují polohu každé fotografie v~úplném povoleném pořadí výsledků. Když návštěvník přechází dopředu nebo zpět, prohlížeč žádá server o~blízká metadata v~malých autorizovaných skupinách. Server opakuje stejná pravidla chytré galerie, kontroly přístupu ke zdrojům, řazení a~stabilní rozlišení shodných hodnot jako stránka.

Díky tomu zůstávají velké virtuální sbírky praktické: otevření chytré galerie předem nenačítá všechny originály ani metadata všech výsledků. Stávající ovládání klávesnicí, dotykem, celou obrazovkou, přiblížením, zavřením a~opětovným otevřením a~prezentací je společné s~běžnými galeriemi. Pokud je prezentace pro danou chytrou galerii vypnutá, její ovládací prvky se vynechají, zatímco běžné galerie zachovávají obvyklé výchozí hodnoty.

\subsection{Stahování a~bezpečnostní limity}
\index{h00053@chytré galerie!s00224@stahování}

Jsou-li povoleny stahování na webu i~volba stahování chytré galerie, server může připravit ZIP z~aktuálně odpovídajících povolených originálů. Znovu vyhodnotí uložená pravidla a~pravidla přístupu místo odhalení cest souborového systému prohlížeči. Vytváření archivu odmítá sbírky s~více než 5,000 odpovídajícími obrázky a~vynucuje také souhrnný limit bajtů původních souborů. Existující instalace používají limit 2~GiB, pokud není \texttt{smart\_gallery\_zip\_max\_source\_bytes} nastaveno v~\texttt{config.php}. Souběžné požadavky na stejný podpis archivu se zpracovávají postupně pomocí souborového zámku; po získání zámku se znovu zkontroluje mezipaměť, zapíše se jedinečný částečný ZIP a~dokončený archiv se zveřejní atomickým přejmenováním. Protokoly chyb uchovávají pouze omezenou diagnostiku důvodu a~třídy namísto nezpracovaných zpráv výjimek nebo soukromých cest serveru.

Pravidla jsou čitelný verzovaný JSON, nikdy uživatelské SQL: pole a~operátory jsou na seznamu povolených, hodnoty jsou připravené parametry, hloubka je omezena na pět a~přijímá se nejvýše 50 podmínek. Odstraněné odkazy neodpovídají ničemu a~chybné či neznámé verze pravidel se odmítají. Běžný migrační postup přidává úložiště prezentace chytrých galerií; existující chytré galerie dědí výchozí hodnoty Vzhledu a~webu, dokud není uloženo přepsání.

\section{Kooperativní galerie}
\index{kooperativní galerie}
Kooperativní galerie propojují nezávislé instalace s~veřejně dostupným HTTPS; každý správce si ponechává vlastní alba a~originály. Použijte běžný migrační postup a~zapněte kooperativní galerie ve Funkcích na každém serveru. Funkce je ve výchozím stavu vypnutá. Vypnutí zachovává identity, přátelství, návrhy i~nastavení.

V~Nastavení > Pokročilé > Spřátelené galerie přímo spárujte každou dvojici účastníků. Přátelství samo fotografie nesdílí. Ve spolupráci alb vyberte veřejné, uvedené album bez hesla a~bez NSFW a~vytvořte jeho referenční kód. Jeden účastník spojí kódy do neměnného návrhu; každý správce zkontroluje členy a~oprávnění a~výslovně schválí účast na vlastní instalaci. Doručení a~ověření probíhá v~omezených krocích v~bočním panelu. E-mailová pozvánka používá nastavené odesílání pošty; otevření odkazu není souhlas.

Po aktivaci otevřete sdílené fotografie z~panelu nebo odkazu v~místním albu. Každý zdroj má vlastní označení a~stránkování. Sdílí se pouze existující autorizované náhledy JPEG/WebP; originály, chráněná alba a~vzdálený zápis nejsou podporovány. Odkazy fungují i~bez JavaScriptu. Chybějící náhledy vytvořte běžným postupem. Nové fotografie získávají identitu celé cesty včetně přípony; staré názvy zůstávají zachovány bez hromadné obnovy, nejednoznačné vlastnictví se však odmítá.

Odvolání účasti, změna ochrany nebo místní zrušení přátelství blokuje další autorizovaná čtení. Nedostupný server je omezen nejvýše 120 sekundami platnosti ověření; již stažené pixely nelze odvolat. Přidání účastníka vytváří samostatný návrh vyžadující nový souhlas všech. Obnova ověření může probíhat při používání; plánovač je volitelný. Chybějící nebo neznámé povinné schéma odmítá sdílení i~mutace. Podrobnosti a~omezení nasazení uvádí \filepath{docs/COOPERATIVE_GALLERIES.md}.

\subsection{Poloha simulátoru ve Windows}
WinApp 0.2.0 automaticky používá světovou polohu kamery MSFS 2024 a~při omezeném selhání přechází na polohu letadla. MSFS 2020 používá přímo polohu letadla přes stejný SimConnect. Souřadnice jsou volitelné; nedostupný simulátor neblokuje nahrávání. Diagnostika rozlišuje zdroj polohy. Instalátor požaduje ukončení aplikace před výměnou souborů; před aktualizací dokončete rozpracované úlohy. Skutečný simulátor a~aktualizaci ověřte na cílovém PC.

\section{Přidávání, popisování a~uspořádání fotografií}
\label{sec:photo-workflow}
\index{p00157@práce s fotografiemi}
\index{m00085@metadata obrázků}
\index{u00245@uspořádání fotografií}

\subsection{Výběr způsobu nahrávání}
\index{z00283@způsoby nahrávání}

Nahrávání v~prohlížeči se hodí pro každodenní přírůstky. Vyberte galerii, zvolte soubory, zkontrolujte dávku a~sledujte průběh, zatímco server ověřuje média a~připravuje záznamy a~náhledy. Vyhledání souborů v~souborovém systému se hodí pro velké existující stromy složek nebo přenosy pomocí FTP a~hostingových nástrojů. Automatizace a~WebDAV se hodí pro opakované či mobilní postupy po nastavení přihlašovacích údajů s~vymezeným rozsahem.

U~velké akce nahrávejte po smysluplných skupinách. Menší skupiny usnadňují identifikaci chyb a~umožňují průběžnou kontrolu názvů, pořadí a~viditelnosti.

\subsection{Psaní užitečných názvů fotografií}
\index{n00105@název fotografie}
\index{p00142@popisky}
\index{a00004@alternativní text}

Název umožňuje fotografii rychle rozpoznat. Popis nebo popisek může vysvětlit osoby, místo, činnost, technické souvislosti nebo příběh mimo záběr. Užitečné příklady:

\begin{itemize}
  \item \emph{Karlův most před východem slunce};
  \item \emph{Příjezd prvního vlaku do obnovené stanice};
  \item \emph{Rodinná skupina na zahradě, červen 1998};
  \item \emph{Závěrečné přiblížení během sobotní ukázky}.
\end{itemize}

Smysluplný text zlepšuje vyhledávání, přístupnost, pozdější identifikaci a~zážitek čtenářů, kteří chtějí více než vizuální mřížku. Názvy souborů mohou zůstat provozními podrobnostmi, pokud nastavení veřejného zobrazení upřednostňuje pečlivě připravené názvy.

\subsection{Štítky jako druhá vrstva uspořádání}
\index{s00230@štítky!p00158@praktické použití}
\index{s00228@strategie štítkování}

Než přidáte mnoho téměř duplicitních štítků, vytvořte malý slovník. Upřednostněte jeden stabilní výraz, například \emph{letadla}, před několika náhodnými variantami. Štítky mohou představovat témata, osoby, místa, techniky, vybavení, roční období nebo redakční stav.

V~cestovatelském archivu mohou galerie představovat cesty a~štítky opakující se témata. V~rodinném archivu mohou galerie představovat roky a~události, zatímco štítky určují osoby. V~profesionálním portfoliu mohou galerie představovat klienty a~štítky typ výstupu, místo a~specializaci.

Pole zadávání štítků mohou při psaní navrhovat existující štítky. Opětovné použití návrhu předchází náhodným pravopisným variantám a~udržuje cílové stránky štítků jednotné. Návrh je pouze pomůckou při úpravách; o~tom, které štítky se uloží pro galerii či fotografii, nadále rozhoduje běžná akce uložení.

\subsection{Uspořádání pořadí prohlížení}
\index{p00145@pořadí fotografií}
\index{v00263@vyprávění příběhu}

Pořadí mění sbírku v~posloupnost. Začněte lákavým obrázkem, představte místo nebo souvislosti, střídejte široké a~blízké pohledy, ponechte úzce související detaily pohromadě a~zakončete přirozeným závěrem. Chronologické pořadí dobře funguje u~událostí, zatímco vizuální rytmus se hodí pro portfolia a~výstavy.

Veřejné ovládací prvky pořadí pomáhají správcům upravovat viditelné karty. Změna pořadí fotografií na veřejné stránce je také omezená na aktuální viditelný výřez stránkování a~ukládá se přes ověřený koncový bod pořadí obrázků. Je-li aktivní Správce fotografií, přetažení vybraných karet fotografií nad viditelnou kartu podřízené galerie mění gesto ze změny pořadí na ověřený přesun do této podřízené galerie. Stránkování může velkou sbírku rozdělit, proto kontrolujte přechody na hranicích stránek i~první stránku.

\subsection{Kontrola EXIF a~údajů o~poloze}
\index{e00036@EXIF!u00247@uživatelská kontrola}
\index{s00216@soukromí polohy}

EXIF může fotografii obohatit o~datum pořízení, fotoaparát, objektiv, expozici a~souřadnice GPS. Tyto informace podporují mapy, návrhy dat, podklady pro duplicity a~technické souvislosti. Před zveřejněním fotografií pořízených doma, ve škole, na soukromém pracovišti nebo v~citlivých přírodních lokalitách zkontrolujte viditelnost GPS.

\section{Zveřejňování, sdílení a~příklady použití pro různá publika}
\label{sec:audiences}
\index{p00149@postup zveřejňování}
\index{p00190@publikum}
\index{p00184@příklady použití}

\subsection{Veřejné portfolio}
\index{p00144@portfolio}

Veřejnému portfoliu prospívá malý počet výrazných galerií nejvyšší úrovně, jednotné titulní obrázky, stručné popisy a~pečlivé pořadí. Pro specializace napříč projekty používejte štítky. Zdrojové archivy uchovávejte v~soukromých nebo nezveřejněných větvích a~tam, kde postup vyžaduje samostatný výběr, kopírujte vybranou tvorbu do veřejných prezentačních galerií.

Vyzkoušejte portfolio na telefonu, displeji s~vysokou hustotou bodů a~pomalejším připojení. Responzivní výběr náhledů poskytuje výchozí rovnováhu. Postupné zostřování umožňuje nejprve zobrazit malé náhledy v~instalacích upřednostňujících rychlou vizuální odezvu.

\subsection{Soukromý rodinný archiv}
\index{r00195@rodinný archiv}

Uspořádejte obsah podle roku, rodinné větve nebo události. Pomocí nadřazených galerií chráněných heslem vytvořte srozumitelnou soukromou oblast a~umístěte do ní související potomky. Přidávejte jména a~data, dokud jsou vzpomínky dostupné. Skenovaným fotografiím obzvlášť prospívají lidské popisy, protože vložená metadata fotoaparátu mohou chybět.

Chráněnou adresu a~heslo sdílejte vhodným soukromým kanálem. Vysvětlete, že příjemci mohou procházet podřízené galerie, aniž by dostali účet správce. Pravidelně kontrolujte, zda staré odkazy a~hesla stále slouží zamýšlenému publiku.

\subsection{Předávání klientům}
\index{k00063@klientská galerie}
\index{p00147@postup klientského výběru}

Vytvořte jednu chráněnou galerii pro každého klienta nebo zakázku. Nahrajte podklady k~výběru, uspořádejte je do užitečných skupin a~v~popisech vysvětlete očekávání ohledně výběru nebo stahování. Hlasování může podporovat jednoduché zjišťování preferencí, je-li pro tuto galerii povoleno. Galerie konečného předání může obsahovat schválené soubory v~zamýšleném pořadí.

Odkazy pro sdílení nabízejí další řízený způsob vstupu, pokud je nastaven. Před odesláním adresy otestujte přesný průchod klienta v~soukromém okně prohlížeče.

\subsection{Galerie akcí a~komunity}
\index{g00041@galerie akce}

Galerie akce může začít jako nezveřejněná, zatímco fotografie přicházejí. Podřízené galerie použijte pro dny, scény, týmy, místa nebo činnosti. Nadřazenou galerii zveřejněte, jakmile je připraven první užitečný výběr, a~další skupiny přidávejte bez změny známé veřejné adresy.

Veřejné hlasování může zvýraznit oblíbené fotografie publika. Štítky propojují opakující se účastníky nebo témata. Stahování poskytuje pohodlný přenos sbírky, pokud to pravidla akce dovolují.

\subsection{Cestovatelský a~místopisný archiv}
\index{c00018@cestovatelský archiv}
\index{m00078@mapy!p00153@použití při cestování}

Galerie zemí, oblastí, měst a~cest poskytují přirozenou navigaci. Rozsahy dat určují chronologii, zatímco mapy GPS vytvářejí geografickou cestu prohlížení. Zkontrolujte souřadnice ubytování a~soukromých míst. Popisný úvod může vysvětlit trasu, účel a~historické souvislosti cesty.

\subsection{Historický a~institucionální archiv}
\index{h00047@historický archiv}
\index{i00057@institucionální archiv}

Historickým sbírkám prospívají pečlivé poznámky o~zdrojích, formulace nejistých dat, jména, místa a~údaje o~původu. V~popisech rozlišujte potvrzená fakta od rodinných či institucionálních vzpomínek. Štítky mohou propojit osoby, budovy, organizace a~období napříč mnoha galeriemi.

Před rozsáhlou prací s~metadaty vytvořte zálohy a~při plánování externí archivace exportujte informace dostupnými nástroji. Galerie může poskytovat přístupnou prezentační vrstvu, zatímco původní archivační postupy pokračují souběžně.

\section{Porozumění datům galerie}
\label{sec:user-data}
\index{v00255@vlastnictví dat}
\index{d00024@databáze!p00160@pro vlastníky galerií}
\index{s00214@souborový systém!p00160@pro vlastníky galerií}

Vlastníci galerií pracují se dvěma spolupracujícími úložišti: běžnými soubory a~databází. Pochopení jejich rolí usnadňuje zálohování, migrace a~řešení problémů.

\subsection{Co je ve složkách galerií}

Složky galerií obsahují fotografie a~hierarchii složek, takže základní sbírka zůstává rozpoznatelná přes FTP, hostingového správce souborů, zálohovací nástroj nebo místní kopii. Přesunutí či přejmenování souborů mimo aplikaci může vyžadovat vyhledání nebo synchronizaci, aby databáze poznala nový stav.

\subsection{Co je v~databázi}

Databáze uchovává názvy, popisy, data, viditelnost, hesla a~konfiguraci přístupu, veřejné cesty, pořadí, štítky, hlasy, rozměry obrázků, indexy EXIF, kontrolní součty, záznamy náhledů, účty správců, nastavení, auditní informace a~stav pracovních postupů.\footnote{Hesla používají hashe nebo chráněnou konfiguraci odpovídající jejich účelu. Export databáze přesto obsahuje citlivé provozní informace a~zaslouží si bezpečné uložení.}

Databáze zrychluje vyhledávání a~správu. Umožňuje také přiřadit souboru přívětivý název, bohatý popis, štítky, zvolené pořadí a~pravidla přístupu bez změny původních bajtů obrázku.

\subsection{Jak databáze pomáhá uživatelům}
\index{d00024@databáze!p00185@přínosy}

Každodenní přínosy zahrnují:

\begin{itemize}
  \item vyhledávání slov v~popisech galerií a~fotografií;
  \item otevírání galerie pomocí stabilní veřejné cesty;
  \item zobrazování fotografií v~pečlivě sestaveném pořadí;
  \item uchovávání štítků a~hlasů návštěvníků;
  \item ochranu galerie pomocí zděděných pravidel přístupu;
  \item vyhledávání přesně duplicitního obsahu pomocí uložených kontrolních součtů;
  \item přípravu vhodných náhledů bez opětovného čtení každého zdrojového souboru na každé stránce;
  \item zaznamenávání migrací a~rozhodnutí údržby s~auditní stopou.
\end{itemize}

\subsection{Záloha jako úplný celek}

Úplná záloha zahrnuje strom galerií, export databáze, místní konfiguraci a~vlastní podklady instalace. Přiřaďte těmto částem stejné datum zálohy, aby popisovaly jeden soudržný stav. Kopii uchovávejte mimo webový server a~obnovu vyzkoušejte v~samostatném umístění.
\section{Každodenní správa a~péče}
\label{sec:daily-care}
\index{k00062@každodenní správa}
\index{m00081@měsíční údržba}
\index{k00069@kontrolní seznam vlastníka}

\subsection{Po každém důležitém nahrávání}

Zkontrolujte cílovou galerii, počet fotografií, pořadí, názvy, viditelnost a~vzhled náhledů. Otevřete několik fotografií v~lightboxu a~zkontrolujte galerii anonymně. U~chráněné sbírky zopakujte přístup pomocí hesla nebo odkazu pro sdílení z~nové soukromé relace.

\subsection{Měsíční kontrola}

Měsíční kontrola údržby by měla zahrnovat:

\begin{enumerate}
  \item Zkontrolujte Stav systému a~čekající migrace.
  \item Zkontrolujte nedávné auditní události administrace.
  \item Prohlédněte dostupné aktualizace aplikace a~informace o~jejich vydáních.
  \item Ověřte, že plánované nebo ruční zálohy byly dokončeny.
  \item Otevřete reprezentativní veřejné a~chráněné galerie.
  \item Zkontrolujte údržbu náhledů, pokud se výrazně změnily zdrojové soubory nebo nastavení vykreslování.
  \item Odstraňte nebo zrušte přihlašovací údaje a~způsoby sdílení, které již splnily svůj účel.
\end{enumerate}

\subsection{Před významným veřejným spuštěním}
\index{k00068@kontrolní seznam spuštění}

Zkontrolujte názvy, popisy, titulní obrázky, pravopis, mobilní rozvržení, soukromí, zveřejnění GPS, stahování, mapy, hlasování, výsledky vyhledávání a~očekávání ohledně kontaktu. Pokud publikum může používat mobilní připojení, otestujte web s~profilem pomalé sítě. Bezprostředně před spuštěním vytvořte zálohu připraveného stavu.

\subsection{Když se správou pomáhá další osoba}

Každému správci přidělte vlastní účet tam, kde model účtů podporuje zamýšlený postup. Vysvětlete rozsah galerií, očekávání soukromí, pravidla pojmenování, slovník štítků, odpovědnost za zálohy a~postupy odstraňování. Samostatné účty zlepšují dohledatelnost odpovědnosti v~auditních záznamech a~oddělují osobní záznamy kontroly duplicit.

\section{Prostředí veřejného návštěvníka}
\label{sec:public}
\index{v00249@veřejná galerie}
\index{s00211@směrování}

\subsection{Co návštěvník vidí}
\index{z00278@zážitek návštěvníka}

Návštěvník otevře domovskou stránku webu nebo přejde přímým odkazem do galerie či na fotografii. Stránka představuje název webu, navigaci, název galerie, popis, drobečkovou navigaci, štítky, podřízené galerie a~fotografie, které návštěvník smí vidět. Chráněný obsah před odhalením soukromých podkladů vyžádá příslušný krok přístupu.

Každá galerie má veřejnou adresu, kterou lze uložit do záložek nebo sdílet. Přímá adresa fotografie otevře kontext galerie a~dovede návštěvníka k~dané fotografii. Smysluplné názvy a~stabilní veřejné cesty pomáhají udržet odkazy srozumitelné ve zprávách, záložkách a~výsledcích vyhledávání.

\subsubsection{Typická návštěva}

Jeden běžný průchod návštěvníka:

\begin{enumerate}
  \item Otevřete domovskou stránku a~vyberte zajímavý titulní obrázek galerie.
  \item Přečtěte si úvod galerie a~projděte podřízené galerie nebo štítky.
  \item Procházejte menší náhledové obrázky posouváním stránky.
  \item Otevřete jednu fotografii ve velkém prohlížeči.
  \item Přecházejte mezi okolními fotografiemi pomocí tlačítek, kláves, pásu nebo vybraného režimu prohlížení.
  \item Otevřete mapu, hlasujte, vyhledávejte nebo stahujte, pokud galerie tyto akce nabízí.
  \item Pomocí drobečkové navigace se vraťte do širší sbírky.
\end{enumerate}

\subsection{Hierarchie galerií a~stránkování}
\index{s00227@stránkování}
\index{h00046@hierarchie galerií}

\subsubsection{Přímý obsah a~stabilní pořadí}

Galerie může nejprve zobrazovat menší galerie a~poté fotografie. Například \emph{Itálie 2026} může obsahovat \emph{Řím}, \emph{Florencii} a~\emph{Benátky} a~současně zobrazovat několik přehledových fotografií přímo na stránce Itálie.

Stránkování rozděluje dlouhou sbírku na zvládnutelné stránky. Pomáhá rychle zobrazit první stránku a~udržuje posouvání pohodlné. Malá výstava může nejlépe vypadat na jediné stránce. Svatbě s~1,500 fotografiemi obvykle prospějí stránky nebo podřízené galerie. Velký prohlížeč fotografií může pokračovat širší posloupností galerie a~podle potřeby načítat další informace.\footnote{Stránkování mění počet karet zobrazených najednou. Samotné fotografie i~jejich pořadí ponechává beze změny.}

Dlouhé veřejné seznamy obsahují plovoucí ovládací prvek \ui{Nahoru}, jakmile se návštěvník posune do výpisu galerie. Používá plynulé posouvání prohlížeče, mimo aktivní výpis mizí a~při otevřeném lightboxu nebo celoobrazovkovém prohlížeči zůstává skrytý, aby nekonkuroval ovládání fotografií.

\subsection{Vyhledávání, štítky a~navigace}
\index{v00261@vyhledávání}
\index{s00230@štítky}

Vyhledávání může procházet celý veřejný web nebo zůstat v~aktuální galerii a~jejích potomcích. Návštěvník v~\emph{Rodinném archivu} může hledat jméno osoby pouze v~této větvi, zatímco návštěvník na domovské stránce může prohledávat všechny veřejné sbírky. Současná služba vyhledávání začíná na dvou viditelných znacích dotazu, vrací omezenou směs výsledků galerií a~fotografií a~řadí silnější shody názvů, jmen a~štítků před slabší shody textu. Výchozí požadavek vrací až 12 sloučených výsledků a~služba omezuje požadavek nejvýše na~30 výsledků.

Prohledávané podklady zahrnují veřejné názvy a~popisy galerií a~fotografií, názvy souborů obrázků, štítky galerií a~fotografií a~jejich veřejné popisy, aktivní obsah udržovaných jazyků, je-li dostupné jeho schéma, a~interní prohledávatelný text AI, jsou-li připravena metadata AI. Filtrování přístupu probíhá jako součást dotazu: soukromý, nepřístupný nebo neveřejný obsah obrázků se nevrací jen proto, že odpovídá jeho text. Vyhledávání ve větvi také omezuje kontext výpisu na daný podstrom galerie místo prohledání celého katalogu a~následného skrývání nesouvisejících shod.

Štítky nabízejí další cestu archivem. Štítek \emph{noční fotografie} může propojit Prahu, Londýn a~Tokio, i~když tyto fotografie jsou v~samostatných cestovatelských galeriích. Drobečková navigace ukazuje aktuální místo v~hierarchii a~oblíbené zkratky umísťují důležité galerie přímo do záhlaví webu.
\section{Správa médií a~galerií}
\label{sec:admin}
\index{s00219@správa}

\subsection{Pracovní prostředí administrace}

Po přihlášení správce vidí neveřejné ovládací prvky pro přehled, galerie, fotografie, vzhled, štítky, údržbu, aktualizace, protokoly a~nastavení účtu. Mnohé úpravy se otevírají v~pravém postranním panelu, takže při změně a~uložení jedné položky zůstává galerie viditelná.\footnote{Úplná stránka administrace zůstává dostupná pro postupy vyžadující více prostoru nebo pro případy, kdy není dostupné rozšířené ovládání v~prohlížeči.}

Samotná veřejná galerie také zpřístupňuje kontextové ovládací prvky administrace, je-li aktivní běžná relace správce. Ikony tužky u~galerií a~fotografií otevírají stejné úplné editory v~postranním panelu. Kompaktní odstranění galerie používá obnovitelný postup přes Koš, je-li povolen; vypnutí Koše vrací budoucí odstranění galerií ke starší službě trvalého odstranění podstromu. Kompaktní odstranění fotografie je záměrně označeno \ui{Odstranit fotografii ... z~CMS}: tato rychlá akce na veřejné stránce odstraní řádek obrázku z~katalogu a~nespouští úplnou službu odstranění zdrojového souboru. Chcete-li odstranit také originál a~vygenerované odvozeniny, použijte běžný postup odstranění fotografie nebo hromadného odstranění. Režim anonymního náhledu tyto kontextové ovládací prvky administrace skrývá.

Tabulka galerií v~administraci může filtrovat řádky podle viditelnosti: všechny, nezveřejněné, veřejné nebo soukromé. Filtrování je pomůcka zobrazení v~prohlížeči, nikoli změna pravidel přístupu. U~řádku skrytého filtrem se také zruší zaškrtnutí a~vymaže se hlavní výběr, aby se skryté zastaralé zaškrtnutí nemohlo omylem dostat do další hromadné akce. Shrnutí uvádí počet zobrazených a~způsobilých galerií.

Strom galerií podporuje rozbalení a~sbalení jednotlivých větví i~volby \ui{Sbalit vše} a~\ui{Rozbalit vše}. Pamatuje si, které větve správce sbalil. Prohlížeč odesílá ID sbalených galerií ověřenou akcí chráněnou CSRF a~aplikace tento stav rozhraní ukládá, takže návrat do stromu administrace zachová stejné větve zavřené místo opětovného rozbalení celé hierarchie. U~potomků skrytých sbalenou větví se před hromadnou prací také zruší zaškrtnutí. Změna pořadí obnovuje ovládací prvky stromu, když galerie získá nebo ztratí potomky.

\subsection{Spolehlivé změny přímo v~postranním panelu}
\index{s00219@správa!z00281@změny v postranním panelu}
\index{p00146@postranní panel!d00032@dokončení změny}

Je-li dostupný JavaScript, trvalá akce zahájená v~pravém postranním panelu administrace se dokončí v~tomto panelu. Panel zůstává vložený a~otevřený, adresa v~prohlížeči se nemění a~stránka se nenačítá znovu jen kvůli zobrazení výsledku. To platí pro vytvoření, úpravu, přesun a~odstranění galerie, nahrávání, úpravy, odstranění, přesun, pořadí, viditelnost, NSFW a~nastavení titulního obrázku u~fotografií, skenování/import, návrhy dat EXIF, vytváření náhledů, Přejmenování médií, Organizátor metadat, kontrolu a~odstraňování v~Detektoru duplicitních fotografií, přihlašovací údaje Automatizace nahrávání, stahování do cíle při Migraci galerií, úpravy štítků a~vytvoření, úpravu, umístění, odpojení, duplikování a~odstranění chytré galerie. Samostatná stránka administrace a~běžný postup POST/přesměrování zůstávají dostupné při vypnutém JavaScriptu nebo přímém otevření trasy.

Obsah panelu smí nahradit pouze nejnovější požadavek na jeho otevření; zavření nebo otevření jiné položky ruší oprávnění starší odpovědi. Fokus klávesnice zůstává uvnitř otevřeného panelu a~vrací se k~ovládacímu prvku, který jej otevřel. Neuložený běžný text může být uchován v~omezené paměti prohlížeče pro aktuální stránku a~obnoven při návratu stejného formuláře. Hesla, tokeny, klíče API, vstupy souborů, identity operací a~další soukromé hodnoty jsou vyloučeny a~nic se neukládá trvale do úložiště prohlížeče.

Uložený stav databáze a~souborového systému zůstává rozhodující. Po uložení server popíše změněnou entitu, fragment panelu, který může vyžadovat nové vykreslení, a~všechny ovlivněné veřejné kontexty. Používají se stabilní identifikátory galerií, obrázků, štítků a~chytrých galerií místo předpokladu, že stará URL stále označuje uložený objekt. To je důležité po přejmenování složky či identifikátoru galerie v~URL, změně rodiče galerie nebo identifikátoru štítku v~URL: aplikace může načíst nový kanonický zdroj vykreslení a~ponechat viditelnou adresu prohlížeče beze změny. U~přesunů se popisují zdrojová i~cílová galerie; při změně rodiče starý a~nový rodič i~přesunutá galerie; změny umístění chytré galerie popisují každý ovlivněný kořenový nebo fyzický rodičovský kontext.

Úspěšný požadavek na fragment sám o~sobě nedokazuje, že je změna viditelná. Obnovené HTML vykreslené serverem musí splnit pozorovatelnou podmínku specifickou pro operaci, pokud taková existuje. Příklady zahrnují přítomnost či nepřítomnost galerie nebo karty a~celkové počty, přesnou účinnou viditelnost, uloženou revizi galerie, nepřítomnost, viditelnost či stav NSFW obrázku, revizi obrázků celé galerie pro změny mimo stránku, identitu titulního obrázku, uložené pořadí obrázků, identitu štítku a~umístění či pořadí chytré galerie. Kontroly zohledňující stránkování přijmou přímo viditelnou odpovídající kartu před použitím souhrnných počtů, zatímco výsledek mimo stránku nadále čeká na shodu rozhodujícího celkového počtu. Oprava mezipaměti náhledů je záměrnou výjimkou: obnovuje ovlivněné kontexty, ale nemá stabilní podmínku veřejného HTML podloženou databází.

Nezávislé požadavky na vykreslení mohou na sdíleném hostingu krátce vracet zastaralá data. PHP Gallery proto provádí jednu centralizovanou omezenou posloupnost opakovaných pokusů pouze tehdy, když samotná změna uspěla a~deklarovaná podmínka ještě není pozorovatelná. Každé dokončení získává novou generaci operace; novější změna zruší nebo nahradí starší načítání, probuzení opakovaných pokusů a~odpovědi panelu pro stejný kontext. Strukturálně platné, ale zastaralé HTML se před nahrazením ověřuje, takže starší odpověď nemůže přepsat novější viditelný stav. Obnovy obcházejí mezipaměti prohlížeče a~zachovávají příslušný stav stránkování, řazení, jazyka, filtru, aktivní karty, posunutí panelu a~posunutí stránky.

Pokud uložení uspěje, ale obnovené zobrazení nelze ověřit v~rámci omezeného počtu opakovaných pokusů, panel oznámí problém synchronizace a~zachová uloženou změnu na serveru. Neoznačí změnu chybně za neúspěšnou, tiše znovu nenačte stránku, nepřejde jinam ani nepřepíše historii prohlížeče. Selhání autentizace, oprávnění, CSRF, připravenosti schématu, bezpečnosti cest a~ověření vracejí rozšířenému postupu očekávanou strukturovanou chybu a~neopakují se. Dynamicky obnovené ovládací prvky zůstávají aktivní, protože panel používá delegovanou obsluhu bezpečnou z~hlediska životního cyklu místo vazeb, které po nahrazení fragmentu mizí.

Vytváření a~nahrávání s~pomocí prohlížeče má další ochrany. Výběr zpracování v~prohlížeči je výslovnou volbou provedení: chybějící schopnost prohlížeče nebo neúspěšná příprava pracovního procesu zastaví postup před uložením místo tichého přepnutí na generování náhledů serverem. PHP před uložením originálů ověří přítomnost všech požadovaných připravených velikostí a~formátů náhledů. Jakmile začne ukládání z~prohlížeče, metadata náhradního přímého postupu se nikdy nepovažují za svolení zopakovat klasický požadavek vytvoření či nahrání, čímž se předchází duplicitním galeriím a~fotografiím. Nastavená velikost dávky ZIP je měkkým cílem balení: jeden nedělitelný balíček fotografie, který cíl překročí, může putovat samostatně v~jedné dávce, ale účinný limit nahrávání PHP zůstává tvrdým omezením. Automatizace nahrávání pro Windows může odesílat úplnou kompatibilní matici JPEG+WebP; moderní server pouze s~WebP přijímá požadované varianty WebP a~bezpečně přeskočí platné dodatečné varianty JPEG, zatímco chybné, neodpovídající, velikostně neznámé nebo nepodporované datové obsahy náhledů zůstávají nepřekročitelnými chybami.

\subsection{Centrální centrum Nastavení}
\index{n00099@Nastavení!c00015@centrální centrum}

Navigace administrace obsahuje centrální stránku \textbf{Nastavení} pro konfiguraci celého webu. Je rozdělena na Obecné, Adresa webu, Veřejný vzhled, Obsah, Média a~prohlížení, Nahrávání a~automatizace, Soukromí a~diagnostika a~Pokročilé. Každý oddíl má vlastní stabilní adresu, takže obnovení, záložky, historie prohlížeče i~náhradní postup bez JavaScriptu se vracejí do stejné kategorie.

Pole ve stylu Spotlight vyhledává nastavení během psaní. Porovnává názvy, popisy, kategorie a~technické názvy včetně ovládacích prvků vlastněných specializovanými stránkami administrace. Výběr výsledku otevře příslušný oddíl nebo specializovanou stránku a~zaměří odpovídající ovládací prvek. Šipky, Enter, Escape a~tlačítko vymazání lze používat bez opuštění klávesnice.

Často používaná globální nastavení lze upravovat přímo v~centru, včetně názvu webu, veřejného jazyka, přepisování URL, veřejného vyhledávání, kde je dostupné, vykreslování náhledů, globální výchozí hodnoty zobrazení EXIF/GPS a~vývojové diagnostiky. Další nastavení zůstávají na specializovaných stránkách, pokud potřebují více souvislostí, přihlašovací údaje, nahrávání souborů, destruktivní akce nebo větší editory. Příklady zahrnují rozvržení Vzhledu, nahrávání, telemetrii, zabezpečení účtu, konfiguraci Google/OpenAI, vlastní CSS, vlastní značku, jazykové balíčky a~údržbu databáze.

Centrum neruší stávající dědění. Rozvržení stránek štítků může nadále dědit globální mřížku a~návrh karet, zatímco volby karty galerie, lightboxu a~EXIF/GPS pro konkrétní galerii zůstávají u~galerie, která je vlastní. Kontextové hodnoty galerií a~fotografií se nemění na globální nastavení.

Tajné údaje se nikdy nezobrazují ve výsledcích vyhledávání ani shrnutích. Hesla, tokeny, přihlašovací údaje API a~podobné hodnoty se objevují pouze jako bezpečné stavy, například Nastaveno, Nenastaveno, nebo jako odkaz na specializovanou stránku.

\subsubsection{Průvodce nastavením}
\index{n00099@Nastavení!p00164@průvodce nastavením}
\index{p00164@průvodce nastavením}

V~horní části \ui{Nastavení} zvolte \ui{Průvodce nastavením} a~zahajte nebo obnovte návrh řízené konfigurace. Průvodce prochází stejnými osmi oddíly jako centrum. Krátké pododdíly udržují související volby pohromadě; \ui{Pokročilé} a~\ui{Expertní a~specializovaná nastavení} odhalují méně běžné možnosti. Při vypnutém JavaScriptu zůstává veškerý obsah pododdílů čitelný a~běžná navigace formulářem funguje.

Každé podporované nastavení má vysvětlení a~příklad. Ponechte zaškrtnuté pole zahrnutí, chcete-li připravit hodnotu ke změně, nebo je zrušte, chcete-li zachovat původní nastavení. \ui{Přeskočit} zachovává celý oddíl. \ui{Další}, \ui{Zpět} ani navigace pododdíly nastavení aplikace neukládají. Podporované volby Vzhledu sdílejí stávající živý náhled; náhled je dočasný a~nemění veřejný web.

Průvodce podporuje bezpečné hodnoty centrálního Nastavení, základní barvy a~rozvržení Vzhledu, volby mřížky, karet a~hlavních štítků, omezené volby nahrávání a~telemetrie, předběžnou přípravu náhledů, ochranu SEO, Koš galerií, automatické aktualizace a~volby Plánované údržby. Přihlašovací údaje, akce s~klíči API, nahrané podklady značky, surové CSS, přemístění úložiště galerií a~destruktivní údržba zůstávají informativní. Jejich specializované stránky se z~návrhu neotevírají; před použitím těchto postupů průvodce opusťte. Přepsání jednotlivých galerií zůstávají beze změny.

Závěrečná kontrola nejprve ukazuje navržené změny. Rozbalte nezměněné podrobnosti nebo použijte \ui{Zobrazit všechna nastavení} a~prohlédněte zachované a~informativní položky. Zaškrtněte pole souhlasu a~změny použijte až po kontrole shrnutí. \ui{Zrušit} nebo \ui{Začít znovu} návrh zahodí. Návrh náleží přihlášenému správci a~aktuální relaci prohlížeče; není sdíleným konfiguračním dokumentem.

Uložení znovu ověřuje aktuální nastavení. Pokud jiná relace změnila ovlivněnou hodnotu, průvodce ji odmítne přepsat. Požadované databázové úložiště musí být ověřeno; chybějící nebo neznámé úložiště zablokuje příslušné uložení. Řádky aplikace a~telemetrie používají transakci, zatímco změny adresy webu mají vratný zápis konfigurace pro obnovu při selhání. Volby Koše a~Plánované údržby se ukládají prostřednictvím jejich stávajících vlastníků jako souvislé skupiny. Tajné údaje se nikdy nevkládají do návrhu, shrnutí ani protokolu průvodce.

\subsubsection{Oprava adresy webu}
\index{n00099@Nastavení!a00002@adresa webu}
\index{k00065@konfigurace!z00274@základní URL}

Pro opravu veřejné adresy URL instalace po nastavení otevřete \ui{Nastavení > Adresa webu}. Zadejte absolutní adresu HTTP nebo HTTPS, například \url{https://example.com}. Cestu zahrňte pouze tehdy, pokud je galerie skutečně instalovaná v~daném podadresáři. Pokud detekce hostingu přidala nežádoucí adresář, odstraňte jej z~tohoto pole a~uložte.

Uložení mění pouze řetězec nejvyšší úrovně \codeword{base_url} v~\filepath{config.php}; nevytváří ani nemění databázové nastavení. Ostatní hodnoty konfigurace, komentáře a~formátování se zachovávají. Soubor i~adresář, který jej obsahuje, musejí být zapisovatelné. Aplikace odmítá adresy s~přihlašovacími údaji, parametry dotazu nebo fragmenty a~odmítá přepsat duplicitní nebo vypočítávané hodnoty konfigurace. Vlastní konfiguraci upravte ručně, pokud ji nelze bezpečně aktualizovat.

Po úspěšném uložení prohlížeč otevře stejný oddíl Nastavení na nové adrese. Následující požadavky používají aktualizovanou konfiguraci. Toto nastavení opravuje generované adresy; nepřesouvá soubory galerií ani nenastavuje DNS a~směrování webového serveru.

\subsection{Správa metadat štítků}
\index{s00230@štítky!s00219@správa}
\index{s00230@štítky!m00082@metadata}

\ui{Administrace > Upravit štítky} je centrálním editorem opakovaně použitelného slovníku štítků. Seznam lze řadit podle celkového využití nebo abecedně a~ukazuje aktuální počet použití každého štítku v~galeriích a~fotografiích dohromady. Výběr štítku zpřístupní jeho kanonický název malými písmeny, identifikátor čisté veřejné adresy URL a~volitelný veřejný popis pro cílovou stránku štítku. Stejné zobrazení uvádí galerie a~fotografie, které štítek aktuálně používají, s~odkazy zpět na příslušný veřejný kontext nebo editor fotografie.

Přejmenování štítku nebo změna jeho identifikátoru v~URL aktualizuje společný záznam štítku, takže všechna přiřazení galeriím a~fotografiím nadále odkazují na tentýž opakovaně použitelný štítek. Názvy a~identifikátory se při uložení normalizují podle bezpečných pravidel štítků a~duplicitní identifikátor se odmítá. Odstranění štítku je odlišné: aplikace jej v~transakci odpojí od galerií i~fotografií a~poté odstraní jeho záznam. Volbou \ui{Zobrazit veřejný štítek} zkontrolujte cílovou stránku a~volbou \ui{Nastavit zobrazení štítků} otevřete ovládací prvky Vzhledu pro stránky a~karty štítků a~prezentaci hlavních štítků.

\subsection{Přepínače funkcí}
\index{p00176@přepínače funkcí}
\index{f00039@Funkce!a00001@administrace}

\ui{Vzhled > Funkce} poskytuje globální přepínače volitelných částí aplikace. Registrované funkce jsou ve výchozím stavu povolené, aby upgrade tiše neodstranil chování, které již existující instalace používá. Správce může vypnout jednotlivé funkce pro návštěvníky, mapy a~nástroje letecké simulace, integrace AI a~nahrávání nebo specializované nástroje administrace, pokud upřednostňuje jednodušší instalaci.

Vypnutá funkce je skrytá z~běžné navigace a~její chráněné trasy odmítají příslušný postup. Vypnutí přepínače neodstraňuje kód funkce, fotografie, databázové záznamy, přihlašovací údaje ani konfiguraci. Opětovné zapnutí proto obnoví běžné vstupní body za stejných kontrol schématu, přístupu a~konfigurace jako dříve. Příklady přepínatelných funkcí zahrnují veřejné vyhledávání, režimy procházení lightboxu, Správce fotografií, přímou správu na veřejných stránkách, chytré galerie, hlasování, hru s~fotografiemi, stahování, mapy galerií a~letů, navigační data, SimBrief, pomoc OpenAI, místní metadata AI, nahrávání API, mobilní WebDAV, migrace galerií, Přejmenování médií a~anonymní telemetrii.

\subsection{Správa galerií}

Správa galerií pokrývá celý život sbírky: vytvoření, vyhledání, pojmenování, popis, data, štítky, soukromí, hesla, vlastní značku, titulní obrázky, rozvržení, pořadí, přesun a~pozdější odstranění. Podřízená galerie může dědit užitečné volby od rodiče. Například každá galerie v~soukromé rodinné větvi může následovat ochranu rodiče a~jednotlivý potomek může používat vlastní volbu prezentace tam, kde editor nabízí přepsání. Obrazovka hromadné správy galerií poskytuje odpovídající operace s~více galeriemi, pokud je stejná změna vhodná pro několik sbírek najednou.

Popisy galerií mohou zobrazovat bezpečné externí odkazy zapsané jako běžné odkazy Markdown, \codeword{[link=URL]label[/link]}, \codeword{[url=URL]label[/url]} nebo zkrácené podoby, jejichž viditelný text je zároveň cílem. Odkazují se pouze cíle HTTP a~HTTPS; ostatní schémata zůstávají neaktivním textem. Rozpoznané služby používají přibalené místní symboly značek. U~ostatních veřejných hostitelů může uložení galerie načíst malou ověřenou ikonu webu do místní mezipaměti galerie. Načítání odstraňuje opakování podle názvu hostitele, má limity velikosti, času, přesměrování a~opakovaných pokusů, odmítá soukromé nebo vyhrazené síťové cíle a~nikdy neprobíhá při vykreslování veřejného popisu návštěvníkovi. Pokud mezipaměť nebo její databázová tabulka nejsou dostupné, odkaz zůstává použitelný bez ikony.

Je-li dostupné úložiště čistých veřejných cest, volba \ui{Znovu vygenerovat čisté veřejné cesty} přestaví uložené cesty galerií a~obrázků podle aktuální hierarchie a~oznámí samostatné počty galerií a~obrázků. Změna pořadí a~vybraná uložení editoru také obnovují cesty, pokud to jejich změny vyžadují. Tato akce údržby je užitečná po opravě hierarchie, migraci nebo starších importech; sama nezapíná přepisování URL, takže při vypnutém nebo hostingem nepodporovaném přepisování zůstává náhradou směrování pomocí parametrů dotazu.
\subsection{Organizátor metadat}
\index{o00128@organizátor metadat}
\index{e00036@EXIF!o00129@organizátor podle data}
\index{g00040@galerie!a00012@automatické uspořádání}

Editor galerie obsahuje postup \ui{Organizátor}, který převádí plochou sadu přímých fotografií do podřízených galerií podle data. Současná implementace seskupuje podle data pořízení EXIF již uloženého v~databázi obrázků. Druhotné seskupování podle polohy je vyhrazeno pro budoucí použití a~ve verzi~\version{} není aktivní strategií organizátoru.

Vyberte nejstarší a~nejnovější datum EXIF a~poté sestavte náhled. Náhled čte databázová metadata v~malých dávkách a~neskenuje, nepřejmenovává ani nepřesouvá zdrojové soubory. Návrh ukazuje počet přímých fotografií, odpovídající kandidáty, navržené podřízené galerie a~fotografie ignorované kvůli chybějícímu použitelnému datu EXIF nebo datu mimo zvolený rozsah. Nejstarší datum je zvlášť užitečné pro vyloučení fotografií z~fotoaparátu s~nenastavenými hodinami.

Pro každé přijaté kalendářní datum je navrženým názvem podřízené složky stabilní datum ISO \codeword{YYYY-MM-DD}. Zobrazený název používá formátovač zobrazení data galerie a~pokud formátovač nedokáže dodat popisek, použije čitelnou hodnotu den/měsíc/rok. Před navržením nového potomka organizátor hledá existujícího přímého potomka se stejným názvem organizátoru a~znovu jej použije. Opakovaný cyklus náhledu a~použití tak doplňuje stejný denní cíl místo vytváření další odpovídající galerie dne.

Použití návrhu vyžaduje samostatné potvrzení. PHP Gallery vytvoří nebo znovu použije navržené podřízené galerie podle data a~přesune odpovídající fotografie stejnou službou fyzického přesunu jako běžný hromadný přesun obrázků. Originály a~známé vygenerované odvozeniny se tedy přesouvají společně s~aktualizací vlastnictví v~databázi. Rozsáhlé operace jsou rozděleny do omezených požadavků řízených prohlížečem místo jediného dlouhého požadavku PHP. Před reorganizací velké galerie návrh zkontrolujte a~udržujte zálohu.

\subsection{Přejmenování médií}
\index{p00174@přejmenování médií}
\index{n00108@názvy souborů!p00173@přejmenování}
\index{u00239@údržba médií}

\ui{Galerie > Přejmenování médií} a~panel přejmenování jednotlivé galerie mohou normalizovat fyzické názvy souborů obrázků podle kontextu galerie a~aktuálního pořadí fotografií. Jde o~operaci souborového systému, nikoli kosmetickou změnu názvu, proto nástroj vždy začíná plánem nanečisto. Náhled uvádí původní a~navržený název každého souboru a~označuje soubory, které již odpovídají, chybí, kolidovaly by nebo musejí být přeskočeny kvůli bezpečnostnímu pravidlu.

Způsobilé jsou pouze přímé soubory obrázků uvnitř vybrané složky galerie. Chybějící zdroje, cesty mimo galerii a~nebezpečné cíle se nepřejmenovávají. Je-li preferovaný cílový název obsazený, ale kolizi může vyřešit bezpečná deterministická přípona, náhled tuto úpravu ukáže. Použití plánu vyžaduje potvrzení, že jste náhled zkontrolovali.

Výchozí vzor je \codeword{\{gallery_path\}_\{seq4\}}. Vlastní vzor může používat \codeword{\{gallery_path\}}, \codeword{\{gallery_title\}}, \codeword{\{photo_title\}}, \codeword{\{old_name\}}, \codeword{\{old_stem\}}, \codeword{\{image_id\}} a~zástupné značky pořadového čísla \codeword{\{seq\}} až \codeword{\{seq5\}}. Číslované varianty doplňují číslo na dvě, tři, čtyři nebo pět číslic. Textové složky se normalizují na malé znaky bezpečné pro souborový systém ve stylu ASCII se slovy oddělenými podtržítky, zdrojová přípona se zachovává v~očištěné podobě a~prázdný vlastní vzor použije výchozí hodnotu.

Po úspěšném fyzickém přejmenování PHP Gallery zachovává soulad databázového řádku a~odkazů veřejných cest, aktualizuje odvozené názvy tam, kde je spravuje nástroj přejmenování, podle potřeby přesune nebo zneplatní generované odvozeniny a~odstraní zastaralé záznamy archivů ZIP v~mezipaměti odkazující na původní identitu zdroje. Upozornění na úklid odvozenin nejsou důvodem vracet jinak dokončené přejmenování zdroje, protože reprodukovatelné odvozeniny lze znovu vygenerovat. U~výběrů napříč webem prohlížeč postupuje po omezených částech a~ukazuje výsledek každé galerie.

\subsection{Správa obrázků}
\label{sec:media}
\index{o00115@obrázky}
\index{n00097@nahrávání}

Fotografie lze nahrát v~prohlížeči, zkopírovat do složek a~vyhledat, přijmout nakonfigurovanou automatizací, přenést z~jiné kompatibilní instalace nebo přidat mobilním postupem s~vymezeným rozsahem. Po příchodu fotografie galerie zaznamená její velikost a~dostupné údaje fotoaparátu a~připraví menší zobrazovací kopie používané na celém webu.

Editor jednotlivé fotografie nyní zahrnuje název, popis, varianty udržovaných jazyků, viditelnost, pořadí, štítky, stav NSFW / 18+, soukromé redakční hodnocení pro pravidla chytrých galerií, volitelné meze responzivních náhledů pro fotografii a~údaje EXIF/GPS pouze pro čtení. Existují-li místní metadata AI, editor ukazuje použitý model a~zdroj, prohledávatelný interní text a~nezpracovaný JSON metadat odděleně od veřejného popisku. Ovládací prvky OpenAI vkládají kontrolovatelné návrhy do upravitelných textových polí a~automaticky nenahrazují uložený redakční text.

Hromadné operace s~obrázky jsou záměrně užší než individuální editor a~znovu provádějí serverové kontroly vlastnictví. Současné hromadné akce zahrnují:

\begin{itemize}
  \item přesun vybraných fotografií do existující galerie;
  \item vytvoření nové cílové galerie a~přesun výběru do ní, s~úklidem prázdného nově vytvořeného cíle, pokud přesun selže před potvrzením souborů;
  \item odstranění záznamů vybraných fotografií spolu s~jejich originály a~známými generovanými odvozeninami běžnou službou odstranění;
  \item nastavení vybrané fotografie jako titulního obrázku galerie;
  \item změnu viditelnosti fotografií mezi podporovanými stavy návrhu, veřejné a~soukromé fotografie;
  \item zapnutí nebo vypnutí ochrany NSFW jednotlivých fotografií, je-li známé její schéma;
  \item vytvoření náhledů vybraných fotografií.
\end{itemize}

Pořadí fotografií lze měnit přetažením úchytu řádku a~galerii lze seřadit podle názvu souboru ze sloupce Název. Změny se ukládají běžnou trasou pořadí, nejde o~uspořádání pouze v~prohlížeči.

Je-li povolena volitelná funkce \ui{Správce fotografií}, běžně přihlášený správce může vybírat fotografie a~přímé fyzické podgalerie ve veřejné mřížce galerie. Zaškrtnutí vybírá jednotlivé karty; kliknutí se Shiftem vybírá rozsah; kliknutí s~Ctrl nebo Cmd přepíná kartu; \ui{Vybrat vše}, \ui{Vymazat výběr} a~Ctrl/Cmd+A~podporují větší výběry na viditelné stránce. Karty chytrých galerií nejsou součástí tohoto výběru založeného na souborech. Přetažení nevybrané fotografie začíná výběrem této fotografie, zatímco přetažení existujícího výběru přesouvá vybrané fotografie společně. Viditelná karta podřízené galerie se stává cílem přetažení a~upuštění provádí stejný ověřený přesun fotografií jako výběr cíle. Samotné karty fyzických galerií nejsou zdroji přetažení.

Z~panelu nástrojů může správce odstranit smíšený výběr nebo z~něj vytvořit skutečnou fyzickou galerii. Vytváření nabízí prohledávatelný výběr nadřazené galerie, kopíruje vybrané fotografie do nového kořene a~vybrané stromy fyzických podgalerií pod něj, přičemž zdroje ponechává beze změny. Server zapisuje aktuální doprovodné soubory galerií, odmítá rekurzivní cíle a~existující cílové složky, před kopírováním rezervuje každý cílový kořen, znovu sestaví řádky kopírovaných galerií ze souborového systému a~při selhání odstraní nový částečný výsledek. Smíšené odstranění odstraňuje vybrané fotografie běžným úklidem obrázků a~na vybrané stromy galerií uplatňuje nastavená pravidla obnovitelného Koše, pokud jsou dostupná.

Přesun a~kopírování do existující galerie zůstávají akcemi pouze pro fotografie. Je-li vybrána fyzická galerie, před použitím těchto prvků její výběr zrušte. Výběry fotografií lze také sdílet nebo stahovat. \ui{Sdílet} načte povolené verze obrázků zobrazitelné v~prohlížeči jako objekty \codeword{File} a~otevře nativní nabídku Web Share, pokud prohlížeč a~zařízení podporují sdílení více souborů. Dostupné přijímající aplikace volí prohlížeč, takže PHP Gallery nemůže vynutit Instagram Story, Reel ani jiný konkrétní cíl. Není-li nativní sdílení dostupné, nemůže-li přijmout soubory nebo selže-li po přípravě, správce spustí ověřené stažení ZIP vybraných fotografií jako náhradní postup. Stejný koncový bod ZIP tedy může sloužit výslovnému stažení vybraných fotografií i~náhradě sdílení. Běžní návštěvníci tyto prvky nikdy nedostávají a~server pro každou změnu nebo požadavek ZIP znovu ověřuje přihlášený účet, token CSRF, vlastnictví zdrojů, vlastnictví přímé podřízené galerie a~cílovou galerii.

\subsection{Nahrávání a~vyhledání souborů}
\index{v00260@vyhledání souborů}

Vyhledání v~souborovém systému je užitečné, pokud již byly fotografie zkopírovány na server přes FTP, hostingového správce souborů nebo obnovenou zálohu. Obrazovka vyhledání ukazuje nalezený obsah a~umožňuje správci přidat složky a~soubory do katalogu. Vyhledávání probíhá jako úloha prohlížeče podložená relací v~malých dávkách AJAX místo skenování celého stromu při počátečním požadavku stránky administrace. Současná výchozí hodnota prozkoumá až 80 adresářů v~dávce, přijímá omezenou velikost dávky až 300 a~stav opuštěné úlohy vyprší po dvou hodinách. Velký strom složek tak zůstává na pomalejším hostingu svižný a~současně ukazuje průběh a~konečný seznam výsledků, se kterými lze pracovat.

Skener prochází pod nakonfigurovaným kořenem galerií a~nenásleduje symbolické odkazy na adresáře. Skryté adresáře a~známé názvy generovaných adresářů či mezipamětí, například \filepath{cache}, \filepath{thumbs}, \filepath{thumbnail}, \filepath{thumbnails}, \filepath{preview} a~\filepath{previews}, ignoruje. Nová kandidátní větev musí obsahovat alespoň jeden podporovaný obrázek někde pod sebou. Prázdné strukturální rodičovské složky se přesto mohou stát galeriemi, pokud jsou potřebné k~zachování hierarchie nad potomky obsahujícími obrázky. Složka obsahující pouze doprovodný soubor \filepath{gallery.json} a~žádný podporovaný obrázek v~celé své větvi se v~současném postupu hlásí jako obsahující pouze metadata, nikoli importuje jako galerie. Již indexované složky galerií se vylučují ze seznamu kandidátů a~kandidát, jehož zobrazovaný název by duplikoval existující název sourozence, se nevytváří tiše.

U~každého nespravovaného výsledku může správce importovat strom složek na místě, přesunout podporované fotografie z~vybraných nespravovaných složek do existující galerie nebo odstranit vybrané nespravované složky z~disku. Import rozvíjí výběr v~pořadí od rodiče k~potomkům a~může zahrnout chybějící předky a~potomky obsahující obrázky. Přesun předchází přepsání existujícího cílového souboru volbou deterministického jedinečného názvu, po úspěšných přesunech skenuje cílovou galerii a~odstraňuje zdrojové adresáře pouze tehdy, když jsou prázdné. Akce odstranění z~disku specifická pro vyhledání je záměrně omezena na bezpečné nespravované cesty pod kořenem galerií, kromě kořene samotného, a~odmítá složky již zastoupené databázovými galeriemi. Pro importovanou galerii CMS použijte běžný postup odstranění galerie.

\subsubsection{Doprovodný soubor metadat gallery.json}
\index{g00042@gallery.json}
\index{d00034@doprovodná metadata}

PHP Gallery uchovává aktivní katalog a~vztahy v~databázi, zatímco volitelný soubor \filepath{gallery.json} u~složky galerie zachovává užitečná upravitelná metadata spolu se souborovým systémem. Běžná uložení galerie zapisují doprovodný soubor zpět do složky galerie. Současný výstup obsahuje název, popis, normalizované štítky galerie, viditelnost, pořadí, stav hlasování a~zobrazení názvů souborů, udržovaný jazyk obsahu a~překlady, pokud jsou dostupné, ruční data galerie, vybraná přepsání karty, počtu, lightboxu a~mřížky, meze responzivních náhledů, pravidla přístupu a~uvádění v~seznamech, odkazy na titulní obrázek a~nastavené podklady značky galerie, jsou-li dostupná příslušná schémata. Zápis doprovodného souboru používá stejné kontroly schématu jako odpovídající funkce galerie.

Při vyhledání a~opravě chybějících rodičů může PHP Gallery číst podporované hodnoty z~doprovodného souboru místo odvozování všeho z~názvu složky. Importované řádky konzervativně začínají jako nezveřejněné, pokud není viditelnost uvedena. Současný import obnovuje pole výslovně podporovaná importérem, včetně názvu, popisu a~lokalizace, viditelnosti, ručních dat, hlasování a~prezentace názvů souborů, vybraných přepsání karty, počtu, lightboxu, mřížky a~náhledů, pravidel přístupu a~uvádění, pořadí a~cest podkladů značky. Doprovodný soubor proto považujte za přenositelná metadata složky galerie, nikoli úplnou zálohu databáze: vztahový stav, přihlašovací údaje, protokoly, hlasy, stav hry, definice chytrých galerií a~další záznamy vlastněné databází nadále vyžadují běžné zálohování.

\securitynote{Akce \ui{Odstranit z~disku} na obrazovce vyhledání je skutečným rekurzivním odstraněním nespravovaných kandidátních složek ze souborového systému. Před potvrzením zkontrolujte vybrané relativní cesty. Ochrana známých cest importovaných galerií je dodatečnou bezpečnostní hranicí, nikoli náhradou zálohy.}

Nahrávání v~prohlížeči je pro běžnou práci nejpřívětivější. Může nahrávat do existující galerie nebo vytvořit podřízenou galerii a~nahrát volitelné fotografie ve stejném postupu postranního panelu. Vyberte cíl, zvolte soubory a~sledujte průběh. Aplikace před přijetím každého výsledku kontroluje cíl a~informace o~souboru.

Samotná velikost komprimovaného souboru nedokazuje, že lze obrázek bezpečně dekódovat. Než GD nebo volitelný Imagick alokuje úplný rastr, PHP Gallery kontroluje formát, rozměry, výpočty počtu pixelů, nastavený a~běhový rozpočet paměti a~limity zpracování. Soubor, který nelze bezpečně přijmout, je odmítnut před potvrzením originálu či generovaných odvozenin v~cílové galerii. Tato pravidla platí pro klasický příjem i~příjem připravený prohlížečem; netvrdí, že každý externí dekodér RAW má stejný paměťový model uvnitř procesu.

Nastavení nahrávání rozlišují běžný serverový postup od přípravy s~pomocí prohlížeče. \ui{Povolit všechny formáty podporované serverem} ponechává výběr v~souladu s~rozpoznanými schopnostmi serveru. \ui{Upřednostnit telefonem vykreslené JPG/PNG/WebP, bez RAW/DNG} žádá mobilní prohlížeče o~vykreslenou podobu obrázku, pokud je to možné; hodí se to, pokud postup iPhone ProRAW/DNG vytváří nevhodné barvy. Pokyn prohlížeči není zaručeným převodníkem, proto server nadále ověřuje skutečně nahraná média.

Příprava s~pomocí prohlížeče se povoluje nezávisle a~v~postupu nahrávání je ve výchozím stavu zaškrtnutá. Prohlížeč může před odesláním připravit optimalizované náhledy a~omezené dávky ZIP. Pokud příprava selže před začátkem jakéhokoli zápisu na serveru, běžné výběry obrázků mohou automaticky přejít na klasické serverové nahrávání. Správci mohou omezit výchozí a~maximální počty pracovních procesů prohlížeče, obrázky v~dávce, poměr ZIP k~limitu nahrávání PHP, maximální počet bajtů dávky ZIP a~velikost zdrojového bloku při obnově náhledů v~prohlížeči. Normalizátor také vynucuje tvrdý limit pracovních procesů, aby náhodné nastavení nemohlo vytvořit neomezené souběžné požadavky na sdíleném hostingu.

Současné výchozí hodnoty zpracování v~prohlížeči jsou 8 pracovních procesů, cílových 80~\% zjištěného limitu nahrávání PHP pro každou nekomprimovanou dávku ZIP, nejvýše 8 obrázků v~dávce ZIP a~absolutní limit dávky 24~MiB. Ověřené rozsahy jsou 1--32 pracovních procesů, 10--95~\% limitu nahrávání PHP, 1--64 obrázků v~dávce a~1--128~MiB pro absolutní limit ZIP. Obnova náhledů v~prohlížeči ve výchozím stavu používá zdrojový blok ZIP 512~MiB, omezený na 16~MiB--3~GiB, a~tvůrce zdroje běžně omezuje blok na 96 položek s~interním tvrdým limitem 512. Jde o~bezpečnostní stropy a~výchozí hodnoty, nikoli doporučení zvýšit pomalý sdílený hosting na maximum.

Je-li nahrávání s~pomocí prohlížeče povoleno, výběr přijímá také archivy ZIP, například export iCloud Photos. Prohlížeč archiv prozkoumá místně, přidá podporované obrázky JPEG, PNG, WebP a~GIF do běžné fronty přípravy a~přeskočí složky, skrytá metadata a~nepodporované položky. Vybraný ZIP se nikdy nenahrává do PHP. Podporují se nekomprimované archivy a~archivy Deflate; šifrované, vícediskové, ZIP64, poškozené, podezřele komprimované nebo nadměrně velké archivy se odmítají před zahájením jakéhokoli serverového nahrávání. Import ZIP nemá náhradní klasický postup PHP, proto při výběru archivu ponechte přípravu v~prohlížeči povolenou.

Globální volba automatického přejmenování při nahrávání může po příjmu přes prohlížeč, API nebo WebDAV spustit stejná pravidla pojmenování po skenování jako Přejmenování médií. Ponechte ji vypnutou, pokud jsou externí názvy souborů součástí vašeho archivačního pravidla; povolte ji, pokud mají všechny způsoby příjmu sjednotit názvy podle normalizovaného schématu galerie.
\subsection{Generovaná metadata a~automatizace}
\index{m00083@metadata AI}
\index{o00126@OpenAI}
\index{a00013@automatizace nahrávání}
\index{w00269@WebDAV}

Volitelná fronta analýzy obrázků, místní nebo založená na pracovním procesu, může zaznamenávat vizuální metadata odděleně od lidského názvu, popisu, štítků a~polí EXIF. Správci mohou metadata zkontrolovat v~editoru fotografie a~použít je jako prohledávatelné podklady bez tichého nahrazení redakčního textu. Nakonfigurovaná pomoc OpenAI může vygenerovat popis fotografie či galerie, upravit existující text, shrnout souvislosti, popsat viditelný obsah z~bezpečných generovaných náhledů nebo navrhnout překlad. Tyto akce vkládají kontrolovatelné návrhy; text automaticky nezveřejňují ani neukládají, dokud správce neprovede běžnou operaci uložení.

Pro opakované přenosy vytvořte token automatizace nahrávání omezený na zamýšlenou galerii a~odesílejte podporované soubory do koncového bodu nahrávání. Tokeny lze nezávisle zrušit a~je nutné s~nimi zacházet jako s~hesly. Nezpracovaný klíč API se zobrazuje pouze při vytvoření; poté se ukládá jen jeho hash, proto nový klíč při vytváření zkopírujte do zamýšleného klienta. Globální \ui{Správce API} uvádí přihlašovací údaje automatizace napříč galeriemi a~poskytuje centrální postup zrušení bez odhalení původních tajných hodnot. Ovládací prvky jednotlivých galerií zůstávají dostupné pro stejné přihlašovací údaje s~vymezeným rozsahem.

Doprovodná aplikace pro Windows používá stejný kontrakt nahrávání omezený na galerii. Má dva běžné režimy nahrávání. \ui{Sledování složky} ignoruje soubory přítomné již při zahájení sledování, čeká na ustálení nových souborů, místně si pamatuje potvrzená nahrání a~opakovaně zkouší dočasná selhání s~prodlužováním prodlevy. Volitelně může prostřednictvím SimConnect bezprostředně před každým sledovaným nahráním připojit aktuální zeměpisnou šířku, délku a~výšku kamery nebo letadla Microsoft Flight Simulator a~sledovaný zdroj může odstranit pouze po potvrzení úspěšného nahrání galerií. \ui{Ruční nahrávání} přijímá výběr souborů nezávisle na sledované složce a~může buď generovat responzivní náhledy na PC pomocí pracovních procesů, nebo požádat server galerie o~jejich vytvoření. Práce s~náhledy a~síťové nahrávání probíhají souběžně v~navazujících fázích místo předchozí přípravy celého velkého výběru. Podpora oznamovací oblasti umožňuje pokračování sledování nebo ruční úlohy se skrytým hlavním oknem.

Stejná aplikace Windows může volitelně provozovat pracovní proces metadat AI s~nízkou prioritou. Převezme jednu ověřenou analytickou úlohu, stáhne pouze převzatý podklad, vytvoří interní prohledávatelná metadata pomocí vybrané generace místního backendu a~modelu a~vrátí výsledek do stávajícího koncového bodu automatizace nahrávání. Tento pracovní proces je oddělený od lidských popisků. Akce \ui{Vynutit opětovné generování metadat AI} na úrovni galerie vymaže uložené generované řádky a~starý stav fronty pro danou větev a~zařadí novou práci; náročná analýza nadále probíhá v~pracovním procesu Windows.

Mobilní přihlašovací údaje WebDAV jsou omezené na galerii a~používají autentizaci HTTP Basic a~neuhádnutelný token cesty. Vytvoření připojení zobrazí vygenerované heslo pouze jednou; server poté uchovává jen hash hesla. Stránka administrace uvádí označení připojení, cílovou galerii, adresu URL serveru a~čas posledního použití a~umožňuje připojení zrušit. Koncový bod WebDAV podporuje minimální chování \codeword{OPTIONS}, \codeword{PROPFIND}, \codeword{MKCOL} a~\codeword{PUT} potřebné pro klienty přenosu fotografií, například PhotoSync. Současný příjem WebDAV přijímá JPG, PNG, GIF a~WebP. Je-li tato možnost dostupná, nastavte mobilního klienta na převod HEIC do JPEG před přenosem. Nahraná média poté vstoupí do stejného postupu skenování, volitelného automatického přejmenování, databázového zápisu a~vytváření náhledů jako ostatní nahrávání.
\section{Generování náhledů a~veřejné vykreslování}
\label{sec:thumbnails}
\index{n00094@náhledy}
\index{w00270@WebP}
\index{j00059@JPEG}
\index{n00095@náhledy obrázků}
\index{r00199@rychlost stránky!n00094@náhledy}

\subsection{Co je náhled}
\index{n00092@náhled!d00026@definice}

Náhled je menší zobrazovací kopie fotografie. Původní fotografie může být několik tisíc pixelů široká a~zabírat mnoho megabajtů. Karta galerie na telefonu může potřebovat jen několik set pixelů. Posílání plného originálu pro každou malou kartu by návštěvníky nutilo čekat na detaily, které v~této velikosti nemohou vidět.\footnote{Generování náhledů zachovává zdrojovou fotografii jako podklad pro archivaci a~plné prohlížení. Generovaná kopie je nahraditelnou pomůckou pro procházení.}

Představte si galerii s~60 fotografiemi, z~nichž každá má původní soubor o~velikosti 12~megabajtů. Načtení všech originálů pro první mřížku by mohlo vyžadovat stovky megabajtů. Malé náhledy mohou tuto první fázi prohlížení výrazně zmenšit. Originál zůstává dostupný pro povolené plné prohlížení nebo stažení, zatímco mřížka dostává úsporné náhledy.\footnote{Skutečná úspora závisí na obsahu obrázku, kvalitě, formátu, rozměrech a~kandidátovi vybraném prohlížečem. Příklad ukazuje měřítko, neslibuje pevný kompresní poměr.}

Náhledy plní několik každodenních účelů:

\begin{itemize}
  \item mřížky galerií se zobrazí dříve;
  \item posouvání využívá méně paměti a~síťové kapacity;
  \item mobilní návštěvníci dostávají obrázek bližší velikosti obrazovky;
  \item titulní obrázky galerií a~výsledky vyhledávání zůstávají kompaktní;
  \item velký prohlížeč může začít vhodným náhledem, zatímco se připravuje větší zdroj;
  \item opakované návštěvy mohou znovu použít soubory již uložené v~mezipaměti prohlížeče.\footnote{Mezipaměť prohlížeče uchovává nedávno použité webové prostředky na zařízení návštěvníka podle pravidel serveru a~prohlížeče. Opětovné použití může pozdější návštěvu výrazně zrychlit.}
\end{itemize}

\subsection{Proč existuje několik velikostí náhledů}
\index{n00092@náhled!v00248@velikosti}
\index{p00140@poměr pixelů zařízení}

Jediná malá velikost nemůže vyhovovat každé obrazovce. Úzká karta na telefonu může být ostrá s~náhledem o~300~pixelech. Široká karta na počítači může potřebovat 600 nebo 800~pixelů. Obrazovka s~vysokou hustotou bodů používá více obrazových pixelů pro stejnou fyzickou plochu a~může jí prospět větší kandidát.

PHP Gallery může podle konfigurace a~limitů zdrojového obrázku připravit běžné velikosti, například 300, 600, 800, 960, 1280 a~1600~pixelů. Číslo označuje maximální delší stranu generované kopie. Fotografie na výšku i~na šířku zachovávají původní poměr stran.\footnote{U~fotografie na šířku bývá delší stranou šířka. U~fotografie na výšku bývá delší stranou výška. Kratší strana se vypočítá z~původního poměru stran.}

Menší kandidáti slouží mřížkám a~titulním obrázkům. Větší kandidáti slouží širokým kartám, displejům s~vysokou hustotou bodů, prezentaci vyhledávání, náhledům na mapách a~velkému prohlížeči. Prohlížeč může vybírat z~dostupné skupiny místo přijímání jedné pevné velikosti pro každou situaci.

\subsection{Proč jsou dostupné JPEG a~WebP}
\index{n00092@náhled!f00037@formáty}
\index{j00059@JPEG!p00152@použití pro náhledy}
\index{w00270@WebP!p00152@použití pro náhledy}

JPEG je široce kompatibilní a~běžně známý formát fotografií. WebP často poskytuje podobnou vizuální kvalitu s~menším počtem přenesených bajtů. PHP Gallery může generovat oba, aby kompatibilní prohlížeč mohl používat WebP a~JPEG zůstal spolehlivou alternativou.\footnote{Účinnost komprese se mění podle fotografie a~nastavení kvality. Jemné textury, šum, přechody a~ostré grafické detaily mohou vést k~různým výsledkům.}

Návštěvník formát ručně nevybírá. Stránka prohlížeči sdělí, které verze existují, a~prohlížeč zvolí formát, kterému rozumí. Původní fotografie zachovává vlastní zdrojový formát a~zůstává oddělená od těchto kopií pro procházení.

\subsection{Kdy se náhledy vytvářejí}
\index{n00092@náhled!g00043@generování}
\index{n00092@náhled!u00238@údržba}

Náhledy lze připravit při importu nebo nahrání fotografií, znovu vytvořit pomocí údržby v~administraci nebo generovat chráněným opravným postupem, když veřejná stránka zjistí chybějící očekávanou kopii. Příprava při nahrávání poskytuje pozdějším návštěvníkům nejplynulejší zážitek, protože potřebné soubory již existují před první návštěvou.\footnote{Velké importy mohou vyžadovat značný čas zpracování a~úložiště. Spuštění přípravy jako záměrné úlohy administrace dává vlastníkovi lepší možnost sledovat dokončení.}

Vlastník galerie by měl zkontrolovat údržbu náhledů po:

\begin{itemize}
  \item importu velkého existujícího stromu složek;
  \item změně povolených velikostí nebo kvality náhledů;
  \item obnovení zálohy obsahující originály a~databázová data, ale neúplnou mezipaměť náhledů;
  \item přesunu instalace mezi servery s~odlišnou podporou formátů obrázků;
  \item opakovaných náhradách původním souborem nebo chybějících náhledech ve Stavu systému.
\end{itemize}

Generované náhledy jsou reprodukovatelná data mezipaměti a~lze je obnovit ze zdrojových fotografií. Samotné zdroje vyžadují trvalou ochranu zálohami.

\subsection{Zobrazovací předlohy DNG a~RAW}
\index{d00030@DNG}
\index{f00038@fotografie RAW}

Fotografie DNG mohou zůstat archivním zdrojem, zatímco PHP Gallery vytvoří zobrazovací předlohu pro prohlížeč pro generování náhledů a~veřejné prohlížení. Zobrazovací předloha je generovaná odvozenina, nikoli náhrada původního DNG. Není-li dostupný potřebný místní převodní nástroj nebo převod selže, zprávy nahrávání a~náhledů v~administraci označí neúspěšnou odvozeninu DNG. Údržba náhledů může předlohu vytvořit později; odstranění generovaných odvozenin neodstraňuje zdrojový DNG.

\ui{Diagnostika běhového prostředí} zpřístupňuje aktivní pravidla převodu DNG spolu se skutečnými schopnostmi serveru pro GD, Imagick/ImageMagick a~formáty. Pravidla zdroje mohou používat automatickou náhradu nejprve úplným dekódováním RAW a~poté vloženým náhledem fotoaparátu, výslovně upřednostnit úplné dekódování RAW nebo upřednostnit vložený náhled. Zpracování barev může vynutit sRGB bezpečné pro prohlížeče, co nejvěrněji zachovat dekódovaný vzhled nebo použít vyvážení bílé fotoaparátu, pokud je dostupné. Tyto volby ovlivňují pouze generované odvozeniny pro prohlížeč. Nikdy nepřepisují archivní zdroj DNG a~jejich dostupnost stále závisí na obrazových delegátech a~paměťových limitech hostingu.

\subsection{Co návštěvníci dostávají}
\index{n00092@náhled!v00258@výběr prohlížečem}

Galerie posílá prohlížeči seznam kopií náhledů, které pro kartu skutečně existují. Prohlížeč zvažuje dostupný prostor na stránce, požadavky obrazovky na ostrost, podporovaný formát a~vlastní rozhodování o~načítání. Poté vyžádá vhodný soubor.

Předpokládejme například fotografii s~kopiemi o~300, 600, 800 a~960~pixelech:

\begin{itemize}
  \item malá karta na telefonu může dostat kopii o~300~pixelech;
  \item střední karta může dostat kopii o~600~pixelech;
  \item velká karta nebo karta na displeji s~vysokou hustotou může dostat kopii o~800 nebo 960~pixelech;
  \item velký prohlížeč může vyžádat ještě větší náhled připravený pro tento účel.
\end{itemize}

Přesná volba se může mezi prohlížeči a~obrazovkami lišit. Tato pružnost je užitečná, protože server nemusí odhadovat jedinou univerzální velikost.

\subsection{Responzivní výběr prohlížečem}
\index{r00192@responzivní vykreslovač}
\index{s00222@srcset}
\index{n00092@náhled!r00191@responzivní režim}

\ui{Responzivní výběr prohlížečem} je výchozí režim v~\ui{Administrace > Vzhled > Rozvržení}. Stránka ihned zpřístupní dostupné kandidáty náhledů a~prohlížeč vybere toho, který nejlépe odpovídá vykreslené kartě a~hustotě displeje. Pro tento výběr není potřebný JavaScript.\footnote{Příslušný webový standard nazývá seznam kandidátů \emph{srcset}. Vlastníci galerií se mohou spolehnout na vizuální chování bez znalosti tohoto pojmu značkování.}

Responzivní výběr obvykle přenese jeden vybraný náhled pro každou načtenou kartu a~je běžnou volbou, pokud je nejdůležitější předvídatelné nativní chování prohlížeče.

\begin{tabularx}{\textwidth}{>{\RaggedRight\arraybackslash}p{0.30\textwidth} X}
\toprule
\textbf{Situace} & \textbf{Pravděpodobný výsledek}\\
\midrule
Úzká mobilní karta & Obvykle postačuje menší kandidát náhledu.\\
Široká karta na počítači & Prohlížeč může vybrat středního nebo většího kandidáta.\\
Displej s~vysokou hustotou bodů & Pro stejnou velikost CSS může být vybrán kandidát s~vyšším rozlišením.\\
Nedostupný JavaScript & Responzivní výběr kandidátů nadále funguje.\\
\bottomrule
\end{tabularx}

Kopie o~300~pixelech slouží jako náhradní, pokud je dostupná. Její přítomnost nenutí moderní prohlížeč stáhnout ji jako první, protože větší možnosti se nabízejí současně.

\subsection{Postupné zostřování náhledů}
\index{p00162@progresivní vykreslovač}
\index{p00150@postupné zostřování}
\index{n00092@náhled!p00161@progresivní režim}

\ui{Postupné zostřování náhledů} je volitelný režim v~\ui{Administrace > Vzhled > Rozvržení}. Galerie nejprve zobrazí malý náhled, obvykle o~300~pixelech, a~později nahradí příslušné karty většími kopiemi. Na pomalejším připojení se sbírka může stát dříve rozpoznatelnou za cenu možného přenosu malého i~náhradního náhledu.

\subsubsection{Co se děje na obrazovce}

\begin{enumerate}
  \item Stránka vyhradí konečnou geometrii karty.
  \item Malé náhledy začnou vyplňovat viditelné karty.
  \item Příslušné karty se zařadí do fronty pro většího kandidáta.
  \item Větší kopie se načítá, zatímco malá zůstává viditelná.
  \item Karta se nahradí, jakmile je větší kopie připravená.
\end{enumerate}

Větší vylepšení se řadí do fronty místo zahájení pro všechny karty současně a~viditelné karty mají přednost před vzdáleným obsahem. Současná progresivní volba platí pro karty fotografií v~otevřené galerii. Titulní obrázky galerií, buňky koláží, karty galerií na domovské stránce a~obrázky administrace nadále používají responzivní výběr prohlížečem. Rozhraní administrace označuje postupné zostřování jako Beta.

\subsubsection{Kompromis při přenosu}

Progresivní režim může pro jednu fotografii přenést dva soubory: první malou kopii a~ostřejší náhradu. Responzivní režim často přenese jednu kopii vybranou prohlížečem. Progresivní režim tedy upřednostňuje dřívější vizuální zaplnění, zatímco responzivní režim obvykle snižuje celkový přenos pro stejnou kartu.

\subsection{Který režim má vlastník zvolit?}
\index{n00092@náhled!v00259@výběr vykreslovače}

Zvolte \ui{Responzivní výběr prohlížečem} pro spolehlivý výchozí režim, účinnou nativní volbu prohlížeče a~obvykle jediný přenos náhledu na načtenou kartu.

Vyzkoušejte \ui{Postupné zostřování náhledů}, pokud je velmi důležitá časná odezva stránky, mnoho návštěvníků používá pomalá připojení a~galerii vyhovuje nejprve malá vizuální podoba. Než jej zvolíte jako běžnou prezentaci, otestujte jej na reprezentativních galeriích a~zařízeních.

Užitečné otázky:

\begin{itemize}
  \item Vidí návštěvníci první fotografie dříve?
  \item Vypadá malá verze během zostřování přijatelně?
  \item Kolik dodatečných dat se přenese při typické návštěvě?
  \item Probíhá zostřování na telefonech a~displejích s~vysokou hustotou pohodlným tempem?
  \item Vyhovuje zvolený režim pravděpodobné kvalitě připojení publika?
\end{itemize}

\subsection{Priority načítání a~stabilní rozvržení}
\index{o00119@odložené načítání}
\index{s00223@stabilita rozvržení}
\index{n00092@náhled!o00119@odložené načítání}

První viditelné fotografie dostávají při načítání přednost, protože utvářejí první dojem návštěvníka. Fotografie níže na stránce používají odložené načítání, tedy prohlížeč může počkat, až se přiblíží viditelné oblasti. Nepoužívá tak kapacitu připojení na spodní část dlouhé galerie, zatímco návštěvník ještě sleduje horní část.\footnote{Načasování odloženého načítání částečně náleží prohlížeči. Prohlížeče zvažují posouvání, vzdálenost od viditelné oblasti, stav připojení a~vlastní pravidla implementace.}

Galerie zná šířku a~výšku indexovaných fotografií. Podle těchto proporcí vyhrazuje stabilní tvary karet ještě před dokončením načítání obrázků. Text, tlačítka a~okolní karty proto při zobrazování náhledů zůstávají na svých místech.

Návštěvníci upřednostňující omezený pohyb dostávají klidnou prezentaci bez povinné nepřetržité animace načítání. Odkazy fotografií zůstávají použitelné po celou dobu načítání a~zostřování.

\subsection{Kvalita náhledů a~úložiště}
\index{n00092@náhled!k00073@kvalita}
\index{n00092@náhled!u00242@úložiště}

Vyšší kvalita zachovává více jemných detailů a~může vytvářet větší soubory. Nižší kvalita šetří místo a~přenos za cenu viditelnější komprese. Nejlepší nastavení závisí na tématu, použití obrazovky, kapacitě hostingu a~publiku. Detailní listoví, vlasy, text a~jemná architektura mohou kompresi odhalit snadněji než velká měkká pozadí.

Každá další povolená velikost a~formát zabírají úložiště, protože vytvářejí další kopii každé fotografie. Přínosem je přesnější doručení různým obrazovkám. Stav systému a~statistiky úložiště pomáhají vlastníkovi pochopit výslednou mezipaměť. Velký zdrojový archiv může vytvořit rozsáhlou sbírku náhledů, ale tyto soubory zůstávají reprodukovatelné z~originálů.

\subsection{Kompatibilita náhledů a~režimy údržby}
\index{k00064@kompatibilita náhledů}
\index{u00240@údržba náhledů}
\index{o00114@obnova náhledů v prohlížeči}

Panel údržby má dvě pravidla kompatibility. \ui{Moderní} generuje náhledy WebP všude, kde server umí zapisovat WebP. \ui{Starší kompatibilita} generuje vedle WebP také náhradní náhledy JPEG. Po trvalém přepnutí instalace na Moderní volba \ui{Odstranit starší náhledy JPG} odstraňuje pouze generované soubory náhledů JPEG; originály, odvozeniny WebP a~zobrazovací předlohy DNG se zachovávají.

Úplná kontrola \ui{Zkontrolovat chybějící náhledy} prozkoumá importované fotografie a~zaznamená galerie s~chybějícími, zastaralými, metadatově neúplnými nebo poměrově nesprávnými generovanými variantami, aniž by něco vytvářela. Uložený výsledek umožňuje volbu \ui{Vytvořit chybějící náhledy}, která se zaměřuje pouze na hlášené obrázky. \ui{Vytvořit všechny náhledy} v~Údržbě zahrnuje všechny galerie; stejná akce v~editoru Obrázky galerie zahrnuje tuto galerii a~při volbě \ui{Zahrnout podgalerie} také její potomky. Tyto operace administrace zůstávají hlavním záměrným postupem údržby. Veřejné procházení nespouští neomezenou hromadnou opravu, ale současný vykreslovač má malou chráněnou přípravu na pozadí pro konkrétní karty, které musely použít povolené původní médium, protože očekávaná odvozenina chyběla.

\ui{Obnovit databázi náhledů} uvádí databázová metadata do souladu s~existujícími generovanými soubory a~odstraňuje generované soubory s~nesprávným poměrem stran místo jejich zveřejnění. \ui{Odstranit všechny náhledy} je samostatně potvrzovaná destruktivní operace mezipaměti; zdrojové fotografie zůstávají nedotčené a~odvozeniny lze později obnovit.

\ui{Obnova náhledů v~prohlížeči} je u~těchto akcí obnovy vybrána jako výchozí, je-li dostupné zpracování v~prohlížeči. Pro generování na serveru zrušte její zaškrtnutí. Server posílá původní soubory ověřenému prohlížeči v~omezených zdrojových blocích ZIP, prohlížeč generuje varianty náhledů a~připravené dávky ZIP náhledů se nahrávají zpět k~ověření a~instalaci serverem. Limity zdrojových bloků a~dávek se nastavují spolu s~nahráváním v~prohlížeči, takže postup nezávisí na jediném nadměrném požadavku PHP.

\subsubsection{Chráněná veřejná příprava náhledů}
\index{p00169@předběžná příprava náhledů}
\index{n00092@náhled!a00010@automatická oprava}

Pokud server vykreslil veřejně povolenou fotografii pomocí náhradního původního média, může tuto přesnou fotografii a~chybějící podporované velikosti označit jako kandidáta přípravy na pozadí. Prohlížeč čeká na nečinnost, zastaví se při skrytí stránky, posílá nejvýše dva kandidáty v~požadavku, provede nejvýše čtyři požadavky při jednom načtení stránky a~odděluje je přibližně 1,8~sekundy. Server nezávisle omezuje normalizovaný odeslaný seznam na 24 ID obrázků, ale skutečný rozpočet generování je menší: anonymní požadavek zpracuje nejvýše dva zdrojové obrázky přibližně za čtyři sekundy, zatímco ověřený požadavek správce nejvýše čtyři přibližně za osm sekund.

Kandidát nese token HMAC vázaný na ID obrázku, ID galerie, název souboru, relativní cestu a~složku galerie, které server původně vykreslil. Před generováním koncový bod znovu načte obrázek a~galerii, ověří token, zopakuje aktuální kontroly přístupu návštěvníka, normalizuje požadované velikosti podle nastavené sady náhledů, konzultuje metadata náhledů, pokud jsou dostupná, a~přeskočí varianty, které již lze vykreslit. Prodleva 60~sekund pro každý obrázek brání opakovaným nákladným pokusům. Jediný globální zámek pracovního procesu bez čekání znamená, že souběžný veřejný provoz nemůže nahromadit několik úloh dekódování obrázků; zaneprázdněný požadavek jednoduše oznámí aktivní jiný pracovní proces přípravy. Příprava je ve výchozím stavu povolena interním nastavením aplikace \setting{thumbnail_background_warmup_enabled} a~zapisuje omezené diagnostické události do protokolu administrace, je-li dostupné jeho úložiště.

\securitynote{Příprava na pozadí nemění soukromou nebo jinak nepřístupnou fotografii na veřejný cíl opravy. Podepsaný kandidát dokládá, že server obrázek vykreslil jako kandidáta náhradního zobrazení, a~koncový bod před otevřením zdroje či zápisem odvozenin nadále opakuje aktuální pravidla přístupu.}

\subsection{Pokud náhled chybí}
\index{n00092@náhled!h00051@chybějící}
\index{n00092@náhled!o00112@obnova}

Chybějící náhled se podle obrázku a~dostupných souborů může projevit pomalejším náhradním originálem, prázdným náhledem nebo upozorněním údržby. Otevřete nástroje údržby náhledů, zkontrolujte ovlivněnou galerii či velikosti a~obnovte chybějící varianty. Při opakovaném selhání obnovy zkontrolujte oprávnění zápisu a~podporu formátů obrázků.

Po opravě otevřete galerii s~prázdnou mezipamětí prohlížeče a~ověřte, že titulní obrázky, karty a~velký prohlížeč dostávají očekávané náhledy.
\section{Lightbox, mapy, EXIF a~hlasování}
\label{sec:lightbox}
\index{l00075@lightbox}
\index{e00036@EXIF}
\index{g00045@GPS}
\index{h00048@hlasování}

\subsection{Asynchronní metadata lightboxu}

Lightbox je velký prohlížeč fotografií, který se otevírá nad stránkou galerie. Umožňuje návštěvníkům soustředit se na jednu fotografii a~zachovává ovládání přechodu dopředu a~zpět, zavření, celé obrazovky, prezentace, mapy a~další dostupné prvky.

Galerie připravuje prohlížeč, když se stane užitečným, a~získává informace o~blízkých fotografiích ve zvládnutelných skupinách. Návštěvník tak může procházet velkou galerii, aniž by první stránka musela obsahovat všechny názvy, náhledy, mapové body a~hlasy najednou.\footnote{Současný prohlížeč běžně žádá o~informace ve skupinách blízkých fotografií. Kontroly přístupu se opakují pro každý požadavek, takže pravidla soukromých galerií nadále platí.}

Šipky doleva a~doprava přecházejí o~jednu fotografii zpět nebo dopředu. \codeword{Shift+Left} a~\codeword{Shift+Right} posouvají o~deset pozic pro rychlejší průchod velkou sadou výsledků. Skok sleduje stejnou stabilní seřazenou posloupnost jako běžná navigace, na koncích přechází na opačný konec a~funguje i~v~případě, že aktivní sadou výsledků je chytrá galerie. Přeložená nápověda klávesnice tyto zkratky uvádí; samostatná grafická tlačítka skoku se nepřidávají.

Prohlížeč obvykle začíná vhodně velkým náhledem. Větší nebo původní zdroj lze použít tam, kde to vyžaduje akce a~pravidla přístupu. Blízké náhledy se mohou nenápadně připravovat, aby se další fotografie zobrazila plynule.

Při přípravě nové fotografie prohlížeč používá stejnou oblast průběhu, která hlásí načítání plné kvality. Rychlý pohyb dopředu nebo zpět vždy přidělí správu ukazatele nejnovější požadované fotografii, takže starší požadavek náhledu či metadat nemůže ponechat prohlížeč trvale zaneprázdněný ani nahradit aktuální výběr. Pokud fotografii nelze zobrazit, oblast průběhu ukáže napravitelnou chybu; návštěvník může pokračovat ovládáním dopředu, zpět, zavřením nebo opětovným otevřením. Tyto stavy načítání nemění pravidla přístupu galerie, chráněných médií, chytrých galerií, map, prezentace ani provozu bez JavaScriptu.

\subsection{Přibližování a~posouvání fotografie}
\index{l00075@lightbox!p00180@přiblížení}
\index{o00132@ovládání přiblížení}
\index{g00044@gesto sevření prstů}

Rozsah přiblížení lightboxu je 100--400~\%. Použijte ovládací prvky minus a~plus nebo výběrem procentního údaje obnovte 100~\%. Klávesové zkratky jsou \codeword{+} nebo \codeword{=} pro přiblížení, \codeword{-} nebo \codeword{\_} pro oddálení a~\codeword{0} pro obnovení. Ukazatele fokusu a~aktuální procento zůstávají dostupné asistivním technologiím.

Přibližování kolečkem myši a~trackpadem sleduje ukazatel: část fotografie pod kurzorem při změně měřítka zůstává na místě, i~během několika rychlých kroků. Přiblížení sevřením prstů používá střed gesta. Po zvětšení lze fotografii přetahovat vodorovně i~svisle. Při 100~\% na dotykových zařízeních běžné vodorovné přejetí nadále mění fotografie. Kolečko s~modifikátorem Control/Command zůstává prohlížeči k~dispozici pro přiblížení stránky.

Při 100~\% je fotografie přizpůsobena prohlížeči a~vycentrována. Celoobrazovkový režim používá stejné přibližování a~posouvání a~jeho tlačítko Zavřít i~další ovládací prvky zůstávají při zvětšení obrázku klikatelné.

Přiblížení nad 100~\% okamžitě vyžádá chráněný zdroj plné kvality aktivní fotografie. Stávající náhled zůstává viditelný, dokud větší obrázek není připraven, a~aktuální přiblížení i~poloha posunutí se zachovávají. V~běžném, celoobrazovkovém i~mobilním celoobrazovkovém režimu přenos zobrazuje stejný ukazatel průběhu s~procentem a~počtem bajtů. Změna fotografie, zavření prohlížeče nebo zahájení novějšího požadavku plné kvality ruší předchozí přenos. Neúspěšný nebo zrušený požadavek plné kvality ponechá chráněný náhled na místě.

Navigace, zahájení prezentace, zavření a~opětovné otevření vracejí prohlížeč na vycentrovaných 100~\%. Vstup do celé obrazovky nebo její opuštění zachovává aktuální měřítko, pokud se fotografie nezměnila, a~přizpůsobuje posunutí nové viditelné oblasti. Preference omezeného pohybu odstraňují krátký přechod mezi jednotlivými kroky přiblížení.

Přiblížení ovlivňuje pouze prezentaci. Nemění uloženou fotografii ani neobchází oprávnění médií. Pravidla hesel, odkazů pro sdílení, NSFW, chytrých galerií, map, hlasování, stránkování a~metadat lightboxu nadále platí.

\subsection{Pravidla EXIF a~map}

EXIF jsou informace, které mnoho fotoaparátů a~telefonů ukládá uvnitř fotografie. Mohou zahrnovat datum pořízení, fotoaparát, objektiv, expozici, rozměry, orientaci a~souřadnice GPS. PHP Gallery je může používat pro návrhy dat, mapy, technické podrobnosti, pomoc s~řazením a~kontrolu duplicit.

Mapy se načítají, když o~ně návštěvník požádá, což šetří práci návštěvníkům zajímajícím se pouze o~fotografie. Galerie může zobrazit polohy jednotlivých fotografií i~přehled dostupných bodů. Nastavení rodičů a~potomků určují, které galerie tyto informace zpřístupňují. Globální nastavení veřejného zobrazení EXIF/GPS je náhradou pro galerie s~přepsáním nastaveným na dědění. Administrace také ukazuje, kolik galerií nyní má výslovné přepsání, a~umožňuje všechna tato přepsání obnovit na dědění jedním potvrzeným uložením nastavení. Obnovení chrání kontrola schématu: nelze-li ověřit stav přepsání, aplikace změnu odmítne a~odkáže správce na Stav systému místo odhadování.

Výběr \ui{Otevřít fotografii} ve značce mapy galerie otevře fotografii v~aktuálním lightboxu, pokud patří do aktivní galerie. Funguje to i~tehdy, když fotografie leží mimo aktuálně zobrazenou stránku stránkování: prohlížeč vyžádá jen omezené okno metadat kolem vybrané fotografie. V~celoobrazovkovém režimu rozdělené mapy mapa zůstává otevřená, zatímco se mění fotografická část. U~běžného mapového překryvu se mapa zavře a~odhalí vybranou fotografii v~prohlížeči. Pokud rozšířený prohlížeč nedokáže výběr vyřešit, stejný ovládací prvek přejde na kanonickou stránku fotografie; neodkazuje přímo na nezpracovaný soubor obrázku.

\securitynote{Údaje GPS mohou odhalit citlivá místa. Před zveřejněním osobních fotografií by správci měli zkontrolovat zděděná nastavení map, metadata obrázků a~veřejný přístup.}

\subsection{Letové mapy a~navigační data}
\index{s00207@SimBrief}
\index{l00074@letové mapy}
\index{n00101@navigační data}

Sbírky zaměřené na létání mohou využít volitelnou integraci SimBrief k~načtení nejnovějšího provozního letového plánu pro ID nebo jméno pilota. Kompaktní editor přijímá jediný viditelný identifikátor: samotné číslice vybírají postup ID pilota, ostatní text postup jména pilota. Stávající integrace odesílající oddělená pole zůstávají kompatibilní, přičemž při zadání obou má ID pilota přednost. Editor vytváří upravitelný návrh popisu Markdown, ukládá OFP s~galerií a~zapisuje dostupné souřadnice trasy do záznamu letové mapy galerie. Neúspěšný požadavek SimBrief ponechává běžné pole popisu i~ruční editor trasy dostupné. Návrh před vytvořením je soukromý a~dočasný, jak popisuje oddíl o~vytváření galerií.

Ruční text trasy se vyhodnocuje přednostně offline. Nejprve se kontrolují importované databázové řádky OurAirports a~poté malý přibalený náhradní soubor CSV. Dále lze použít platný výsledek poskytovatele z~mezipaměti; pokud správce propojil nakonfigurovaný účet Navigraph a~vzdálené vyhledávání je povoleno, neznámý identifikátor lze vyřešit prostřednictvím Navigraph a~uložit do mezipaměti pod známým cyklem AIRAC. Správce navigačních dat ukazuje počty v~místní databázi a~přibalených datech, nabízí tester vyhledávání identifikátorů, například letišť, radionavigačních prostředků a~bodů, a~může obnovit místní sadu letišť a~radionavigačních prostředků OurAirports. Úložiště OAuth a~účtů Navigraph je volitelné a~lze je odpojit bez odstranění místních dat.

Veřejné vykreslování map nikdy nevolá SimBrief, Navigraph ani jinou živou plánovací službu. Vykresluje pouze souřadnice již uložené s~galerií. Uložené souřadnice OFP SimBrief tedy zůstávají použitelné, i~pokud je vzdálený poskytovatel později nedostupný. Ruční text trasy může také zcela obejít vyhledávání identifikátorů pomocí výslovných bodů \codeword{NAME@latitude,longitude}.

\subsection{Hlasování}

Veřejné karty a~ovládací prvky lightboxu sdílejí stav hlasů obrázku. Současný veřejný hlas je vratným kladným hlasem: opětovné stisknutí aktivního kladného hlasu jej odstraní. Zobrazené skóre je součtem uložených hlasů obrázku; současná implementace nemá samostatný koncový bod hlasování o~celé galerii. Přihlášené hlasy jsou spojené s~účtem, anonymní hlasy používají hash návštěvníka chránící soukromí. Nové kladné hlasy mají omezenou četnost; zrušení existujícího hlasu je povoleno okamžitě. Serverem vykreslené formuláře umožňují přímé chování, zatímco JavaScript odesílá běžné interakce asynchronně a~synchronizuje skóre mezi odpovídajícími zobrazeními. Server ověřuje identitu obrázku, pravidla galerie, přístup, připravenost schématu, CSRF a~integritu požadavku.

\subsection{Hra s~fotografiemi}
\index{h00050@hra s fotografiemi}
\index{p00143@porovnávání obrázků}

Hra s~fotografiemi je volitelný veřejný režim porovnávání. Správce ji povolí pro větev galerií; způsobilé veřejné fotografie z~galerie a~jejích potomků se pak mohou účastnit podle běžných pravidel přístupu a~viditelnosti. Veřejná hra ukazuje dvě fotografie vedle sebe, zaznamenává předložený pár a~umožňuje návštěvníkovi vybrat preferovaný obrázek před přechodem na další porovnání. Galerie může také představovat nejlepší výsledky nashromážděné hrou.

Protože předložení nového páru vytváří trvalý stav hry ještě před výběrem vítěze, Hra s~fotografiemi vyžaduje známé a~dostupné databázové schéma. Nelze-li ověřit požadované úložiště, zápis se odmítne místo tichého přechodu na nezaznamenané porovnání. Globální přepínač funkce může hru skrýt bez odstranění existujících herních dat nebo nastavení povolení jednotlivých galerií.

\section{Detektor duplicitních fotografií}
\label{sec:duplicates}
\index{d00027@detektor duplicitních fotografií}
\index{z00277@záznamy duplicit}
\index{s00206@SHA-256}

Detektor duplicitních fotografií v~administraci ve výchozím stavu analyzuje vybranou větev galerií a~může rozšířit rozsah na celý web, když správce povolí odpovídající volbu. Hlásí přesné i~možné duplicitní páry, poskytuje kontextové odkazy galerií a~fotografií, zaznamenává rozhodnutí kontroly a~předává potvrzené odstranění běžné službě odstranění obrázků galerie.\footnote{Detektor zůstává nástrojem kontroly. Podklady podporují úsudek správce, protože shodná velikost souboru a~fotografická metadata mohou při některých postupech popisovat odlišné fotografie.}

\subsection{Kategorie podkladů}

\subsubsection{Přesné a~možné nálezy}

Přesná duplicita obsahuje stejný obsah souboru. Aplikace jej rozpoznává pomocí kontrolního součtu, který funguje jako velmi podrobný digitální otisk vypočítaný ze souboru. Shodné otisky poskytují silný důkaz, že dva soubory obsahují stejné bajty.

Možná duplicita má podpůrné podobnosti bez přesné shody otisků. Detektor může porovnat velikost souboru a~dostupné údaje fotoaparátu, například rozměry, čas pořízení, fotoaparát, objektiv, expozici, clonu, ohniskovou vzdálenost, ISO a~orientaci. Dvě samostatně exportované kopie jedné fotografie se mohou lišit jako soubory, ale zachovávat dost společných informací, aby si zasloužily kontrolu.

Chybějící údaje fotoaparátu jednoduše poskytují detektoru méně vodítek. Výsledky zůstávají návrhy k~lidské kontrole. Otevřete oba odkazy fotografií, porovnejte jejich kontext galerie a~kvalitu a~odstraňte pouze kopii, jejíž odstranění odpovídá záměru sbírky.

\subsection{Trvalé záznamy kontroly}

Záznamy náležejí ověřenému správci. Záznamy párů ukládají kanonické identifikátory obrázků s~nižším identifikátorem ve sloupci pro nižší hodnotu a~vyšším ve sloupci pro vyšší hodnotu. Záznamy galerií platí pro přesnou uloženou galerii. Rozhodnutí o~nadřazené a~podřízené galerii zůstávají nezávislá.

Dostupné akce kontroly zahrnují:

\begin{itemize}
  \item ignorování vybraného páru;
  \item ignorování všech nálezů z~přesné galerie levého obrázku;
  \item ignorování všech nálezů z~přesné galerie pravého obrázku;
  \item vymazání záznamů aktuálního správce;
  \item odstranění potvrzené duplicity pomocí ověřeného běžného odstranění obrázků.
\end{itemize}

Akce AJAX obnovují fragment detektoru ve stávajícím pravém postranním panelu administrace. Přímé požadavky POST používají přesměrování jako náhradní postup. Serverem spravované skenovací úlohy zachovávají rozsah mezi pokračovacími požadavky a~odstranění vyřazuje ovlivněný obrázek z~aktuálních výsledků úlohy.
\section{Řízení přístupu a~zabezpečení}
\label{sec:security}
\index{z00271@zabezpečení}
\index{c00020@CSRF}
\index{a00009@autentizace}
\index{n00110@NSFW}

\subsection{Vícevrstvé vyhodnocování přístupu}

Soukromí se vyhodnocuje před sestavením chráněného výstupu. Viditelnost galerie a~zděděná pravidla přístupu se kombinují s~viditelností jednotlivých fotografií, stavem hesla a~odkazu pro sdílení a~pravidly věkového omezení, pokud jsou nastavena. Stejné oprávnění platí pro stránky galerií, náhledy, větší náhledy, původní soubory, metadata lightboxu, mapy, výsledky vyhledávání a~stahování.

Operace administrace vyžadují ověřený účet. Formuláře obsahují soukromý token požadavku a~server před změnou trvalých dat ověřuje cílovou galerii, fotografii a~požadovanou operaci.\footnote{Technický název soukromého tokenu formuláře je token CSRF. Vlastníci galerií se s~ním obvykle setkávají jen prostřednictvím běžných formulářů administrace.}

\subsection{Hesla a~odkazy pro sdílení}
\index{o00122@ochrana heslem}
\index{o00118@odkazy pro sdílení}

Chráněné galerie ukládají hashe hesel a~po úspěšném ověření udělují přístup omezený na relaci. Odkazy pro sdílení používají stav tokenu generovaný serverem a~pravidla vypršení platnosti. Obnova hesla používá jednorázové hashované tokeny, volitelný e-mail pro obnovu, nastavenou dobu platnosti a~vybraný poštovní transport. Používejte HTTPS, silné přihlašovací údaje správce, omezené databázové účty a~chráněné tajné údaje pošty a~API.

\subsection{Účty návštěvníků a~registrace pouze na pozvání}
\index{u00237@účty návštěvníků}
\index{p00154@pozvánky návštěvníků}
\index{s00220@správa účtů návštěvníků}
\index{p00182@přihlášení návštěvníka}

Účty návštěvníků jsou volitelnou vrstvou identity pro přizpůsobené funkce galerie. Jsou oddělené od účtu správce a~samy o~sobě neodemykají soukromou galerii ani galerii chráněnou heslem. O~tom, co návštěvník smí vidět, nadále rozhodují stávající hesla galerií, odkazy pro sdílení, přístup k~chytrým galeriím a~oprávnění médií.

Celý subsystém účtů návštěvníků má hlavní přepínač v~\ui{Administrace > Funkce}. Tato hlavní funkce je ve výchozím stavu vypnutá. Když je vypnutá, veřejná galerie neukazuje ovládací prvky Přihlášení/Účet návštěvníka, trasy návštěvníků jsou nedostupné a~položka \ui{Účty návštěvníků} je skrytá v~nabídce Účet administrace. Existující identity návštěvníků, oblíbené fotografie a~soukromé sbírky zůstávají uložené; přepínač řídí dostupnost místo odstraňování dat.

Po povolení hlavní funkce může správce otevřít \ui{Účet > Účty návštěvníků}. Tato stránka zachovává samostatný přepínač režimu \ui{Účty návštěvníků (pouze na pozvání)}. Přepínač režimu řídí přihlášení a~pozvánky návštěvníků uvnitř povoleného subsystému, zatímco hlavní přepínač Funkcí řídí, zda je subsystém vůbec zpřístupněn.

Správce spravuje identity návštěvníků v~\ui{Účet > Účty návštěvníků}. Stávající postup pouze na pozvání zůstává dostupný: správce může volitelně svázat pozvánku se zamýšlenou e-mailovou adresou a~zkopírovat její tajný odkaz zobrazený pouze jednou. Otevření odkazu nevytváří účet. Příjemce zadá e-mailovou adresu, obdrží ověřovací zprávu nakonfigurovaným poštovním transportem, potvrdí ověření v~prohlížeči a~poté zvolí heslo s~alespoň 15~znaky. Neexistuje obecná stránka \ui{Registrace} ani \ui{Založit účet}.

Stejná stránka administrace může také vytvořit účet návštěvníka přímo. Zadejte e-mail návštěvníka a~buď poskytněte dočasné heslo odpovídající pravidlům hesel návštěvníků, nebo heslo ponechte prázdné, aby galerie vygenerovala silné heslo. Dočasné heslo se zobrazí jednou po vytvoření. Pokud se odešle volitelné oznámení o~vytvoření účtu, zpráva obsahuje důvěryhodnou adresu přihlášení a~pokyny, nikdy však dočasné heslo; heslo předejte samostatně důvěryhodným kanálem. Při zapnuté hlavní funkci zůstává přímé vytváření účtů dostupné i~při vypnutém podřízeném režimu veřejného rozhraní návštěvníků, aby bylo možné připravit účty předem, ale tito návštěvníci se nemohou přihlásit, dokud není tento režim povolen.

Přímo vytvořený návštěvník nemůže dočasné heslo používat k~trvalému přihlašování. Po první úspěšné kontrole přihlašovacích údajů se prohlížeč přesměruje na soukromou stránku prvního přihlášení a~musí zvolit jiné heslo splňující pravidla. Dokud nahrazení neuspěje, galerie nezavede běžnou identitu návštěvníka, nevydá \ui{Zapamatovat si mě} ani neumožní oprávnění pro oblíbené fotografie a~sbírky. Běžná obnova zapomenutého hesla může dočasný přihlašovací údaj také nahradit prostřednictvím nezávisle ověřeného postupu obnovy. Souběžné přihlášení správce ve stejném prohlížeči zůstává nezávislé.

Existující účty návštěvníků jsou na stránce administrace uvedeny se stavem účtu a~bezpečnostními ovládacími prvky. Aktivního návštěvníka lze pozastavit volbou \ui{Pozastavit} nebo odhlásit všude; pozastaveného návštěvníka lze znovu aktivovat volbou \ui{Obnovit}. Pozastavení blokuje autentizaci návštěvníka a~prostřednictvím centrální služby životního cyklu zneplatní jeho stávající relaci, zapamatované přihlášení, obnovu hesla, čekající ověření či změnu e-mailu a~připravená oprávnění sdílení sbírek. Obnovení účet znovu aktivuje, ale neoživuje žádný přihlašovací údaj ani relaci z~doby před pozastavením, takže se návštěvník musí znovu běžně ověřit. Vynucené nahrazení hesla při prvním přihlášení zůstává povinné i~přes pozastavení a~obnovení. \ui{Odhlásit všude} ponechává aktivní účet aktivní a~současně ruší jeho oprávnění přihlášení návštěvníka na všech zařízeních. Tyto bezpečnostní prvky zůstávají správci dostupné, když je hlavní funkce povolena, ale podřízený režim veřejného rozhraní návštěvníků vypnut. Vypnutí hlavní funkce skryje i~správu návštěvníků. Žádný z~těchto přepínačů či bezpečnostních prvků neodstraňuje oblíbené fotografie, soukromé sbírky, položky sbírek, fotografie, galerie, chytré galerie, odkazy pro sdílení galerií ani autentizaci správce.

Odstranění správcem zůstává samostatnou destruktivní akcí s~výslovným potvrzením. Odstraňuje identitu návštěvníka a~jeho stav autentizace, životního cyklu, oblíbených fotografií a~sbírek prostřednictvím připravených databázových vztahů. Neodstraňuje fotografie, galerie, chytré galerie, odkazy pro sdílení galerií ani účty a~relace správců.

Veřejné záhlaví ukazuje \ui{Přihlášení} pouze tehdy, když jsou účty návštěvníků povoleny. Po přihlášení ukazuje \ui{Účet}. Stránka účtu zobrazuje ověřený e-mail, odkazy na soukromé stránky \ui{Oblíbené} a~\ui{Sbírky}, \ui{Změnit heslo}, \ui{Změnit e-mail}, \ui{Odstranit účet} a~akci \ui{Odhlásit se} určenou pouze návštěvníkovi. Odhlášení návštěvníka nebo změna jeho oprávnění neodhlašuje správce, který je současně ověřený ve stejné relaci prohlížeče.

\index{o00111@oblíbené}
Přihlášený návštěvník může označit povolenou fotografii jako oblíbenou malým srdcem na kartě galerie nebo v~lightboxu. Stejná fotografie se synchronizuje mezi kartou a~lightboxem a~soukromá stránka \ui{Oblíbené} uvádí uložené fotografie, které návštěvník nadále smí otevřít. Odebrání z~oblíbených odstraní pouze uložený odkaz návštěvníka; fotografii neodstraňuje ani nemění.

Oblíbená fotografie nikdy neuděluje přístup. Pokud se zdrojová galerie později stane soukromou, vyprší oprávnění hesla či sdílení nebo fotografie jinak přestane být pro požadavek povolená, uložená oblíbená položka tato pravidla neobchází a~zdrojová metadata se na stránce Oblíbené nezobrazí. Prohlížeč současně přihlášený jako správce a~návštěvník stále potřebuje nezávislé oprávnění návštěvníka ke zdroji, než může návštěvník uložit obrázek mezi oblíbené.

\index{s00203@sbírky návštěvníků}
Přihlášený návštěvník může také vytvářet soukromé pojmenované \ui{Sbírky}. Sbírka je seřazený seznam odkazů na existující fotografie; není galerií, nekopíruje fotografie a~nemění hierarchii galerií vlastněnou správcem. Z~povolené karty galerie, karty chytré galerie nebo soukromé stránky \ui{Oblíbené} může návštěvník použít \ui{Přidat do sbírky} a~vybrat jednu ze svých sbírek. Soukromá oblast Sbírky podporuje vytváření a~přejmenování sbírky, odebírání fotografií, posouvání viditelných fotografií nahoru či dolů a~odstranění samotné sbírky.

Vlastnictví sbírky a~přístup k~fotografii jsou samostatná rozhodnutí. Uložení fotografie ukládá pouze její interní odkaz a~pořadí. Sbírka neukládá heslo galerie, token sdílení galerie, cestu souborového systému, cestu náhledu ani zapamatované rozhodnutí o~přístupu. Při každém otevření sbírky galerie znovu vyhodnocuje aktuální oprávnění ke zdrojové galerii a~fotografii. Pokud dříve uložená fotografie již není povolená, vynechá se bez odhalení názvu, názvu souboru, názvu galerie, náhledu, údajů EXIF nebo důvodu odmítnutí. Uložený odkaz může zůstat, aby se fotografie znovu zobrazila, pokud se později vrátí běžný přístup ke zdroji.

Toto pravidlo platí i~při současném přihlášení stejného prohlížeče jako správce a~návštěvníka. Viditelnost správce se nestává oprávněním sbírky návštěvníka a~sbírku nelze nikdy použít jako druhou cestu kolem hesel galerií, uděleného sdílení, věkových omezení nebo kontrol poskytování médií. Oblíbené fotografie a~sbírky jsou nezávislé: fotografie může být v~jednom, obou nebo ani jednom a~odstranění sbírky nikdy neodstraňuje oblíbenou položku ani původní fotografii.

\index{s00204@sdílení sbírek}
Návštěvník může z~vlastní sbírky vytvořit jeden odkaz pro sdílení neuvedený v~seznamech a~určený pouze ke čtení. Odkaz vyprší po 30~dnech a~lze jej kdykoli nahradit nebo zrušit. Úplný tajný odkaz se zobrazí pouze bezprostředně po vytvoření či nahrazení; později stránka sbírky ukazuje jen aktivní sdílení a~čas vytvoření a~vypršení. Nahrazení sdílení znepřístupní předchozí odkaz. Zrušení sdílení odstraní pouze oprávnění sdílení a~neodstraňuje sbírku, její položky, oblíbené fotografie ani fotografie.

Příjemce nepotřebuje účet návštěvníka. Otevření tajného odkazu jej ověří, zavede jen úzké oprávnění relace pro tuto jednu sbírku a~přesměruje na čistou adresu bez tajného údaje. Odkaz je opakovaně použitelný až do vypršení nebo zrušení, takže jej automatické skenery odkazů v~poště či chatu nespotřebují. Sdílená stránka není uvedená v~seznamech, nelze ji ukládat do mezipaměti a~je označena jako neurčená pro indexování vyhledávači.

Sdílení uděluje přístup pouze ke kontejneru sbírky. Nikdy neodemyká zdrojovou galerii ani fotografii. Při každém sdíleném zobrazení se každý uložený odkaz fotografie znovu kontroluje vůči aktuálnímu oprávnění příjemce ke zdrojové galerii. Zdroj chráněný heslem zůstává skrytý, dokud jej příjemce nezávisle neodemkne běžným mechanismem galerie; obdobně jej může běžně povolit nezávislé existující oprávnění sdílení galerie. Ani při současném přihlášení stejného prohlížeče jako správce se při plnění sdílené sbírky nepoužívá výjimka správce. Přímé trasy náhledů a~médií nadále vynucují vlastní existující pravidla zdrojů. Nepřístupné položky se vynechávají bez odhalení skrytého názvu, názvu souboru, cesty, galerie či důvodu.

Pozastavení vlastníka sbírky ruší aktivní sdílení sbírek a~obnovení účtu je neoživuje. \ui{Odhlásit všude} ruší pouze oprávnění přihlášení návštěvníka a~neruší jinak platné sdílení. Odstranění vlastníka nebo sbírky zneplatní její sdílení bez odstranění zdrojových fotografií. Globální přepínač funkce Účty návštěvníků také řídí sdílení sbírek a~zůstává ve výchozím stavu vypnutý.

Volba návštěvníka \ui{Zapamatovat si mě} používá vyhrazený rotující přihlašovací údaj návštěvníka místo tokenu trvalého přihlášení správce. Obnovení tohoto údaje přihlašuje pouze návštěvníka a~nepočítá se jako nedávné opětovné ověření heslem. Požadavky na obnovu zapomenutého hesla vracejí obecnou odpověď a~používají dvoukrokové potvrzení obnovy, takže automatické skenery odkazů nemohou obnovit heslo otevřením URL.

Citlivé změny účtu vyžadují nedávné prokázání hesla návštěvníka. Návštěvník obnovený pouze pomocí \ui{Zapamatovat si mě} je proto před změnou hesla, zahájením změny e-mailu nebo odstraněním účtu požádán o~potvrzení aktuálního hesla návštěvníka. Návratová cesta je omezena na tyto akce účtu a~nemůže přesměrovat na libovolnou stránku.

Změna hesla návštěvníka používá stejná pravidla jako aktivace a~obnova. Úspěšná změna podle bezpečnostního životního cyklu zneplatní oprávnění starého hesla, relace návštěvníka a~zapamatované přihlašovací údaje a~poté požádá o~nové přihlášení. Uložené oblíbené fotografie zůstávají přiřazené stejnému účtu návštěvníka. Správce přihlášený ve stejné relaci prohlížeče zůstává přihlášený.

Změna e-mailu návštěvníka probíhá po etapách. Odeslání nové adresy nenahrazuje aktuální ověřenou adresu. Galerie nejprve odešle ověřovací odkaz na kandidátní adresu v~rámci limitů proti zneužití pošty návštěvníků. Otevření odkazu jej pouze ověří a~připraví krátkodobý stav potvrzení; účet nemění. Konečná změna nastane až po odeslání samostatného chráněného potvrzovacího formuláře návštěvníkem. Úspěšná změna návštěvníka odhlásí, aby mohl při dalším přihlášení použít nový ověřený e-mail.

\ui{Odstranit účet} je trvalá akce pouze pro návštěvníka. Po nedávném potvrzení hesla musí návštěvník odstranění výslovně potvrdit. Operace odstraňuje účet návštěvníka a~jeho data, například oblíbené fotografie a~soukromé sbírky, podle připraveného databázového životního cyklu. Neodstraňuje fotografie galerií, galerie, chytré galerie ani jejich existující odkazy pro sdílení a~neruší souběžné přihlášení správce.

Zabezpečení návštěvníků a~soukromé stránky účtu, oblíbených fotografií, sbírek a~životního cyklu se poskytují jako soukromé odpovědi bez možnosti ukládání do mezipaměti. Odkazy pozvánek, ověření, obnovy a~změny e-mailu používají nakonfigurovaný důvěryhodný původ aplikace. Pošta ověření, obnovy a~změny e-mailu návštěvníků nadále podléhá rozpočtům proti zneužití pro adresu návštěvníka, návštěvníka a~celý systém před použitím nakonfigurovaného transportu PHP mail nebo SMTP.

Sdílení sbírek neuvedené v~seznamech a~pouze ke čtení je dostupné jen uvnitř funkce Účty návštěvníků. Veřejné vyhledávání sbírek, veřejné profily návštěvníků, úpravy či spolupráce příjemců, nahrávání a~komentáře návštěvníků, otevřená registrace, sociální přihlášení návštěvníků, TOTP, přístupové klíče a~přihlášení magickým odkazem zůstávají nedostupné.

\subsection{Účet správce a~nastavení pošty}
\index{u00236@účet správce}
\index{s00213@SMTP}
\index{p00181@přihlášení Google}

Oblast účtu podporuje změny uživatelského jména a~hesla, volitelnou adresu pro obnovu, dobu platnosti obnovy hesla a~testovací zprávu. Pošta obnovy může používat PHP \codeword{mail()} nebo ověřený transport SMTP s~nastavením hostitele, portu, šifrování, uživatelského jména a~hesla. S~přihlašovacími údaji SMTP zacházejte jako s~údaji databáze a~nikdy je nezahrnujte do exportů ani diagnostických snímků obrazovky.

Přihlášení Google je volitelné. Správce nastaví klienta OAuth a~povolenou URI zpětného volání či přesměrování a~poté propojí identitu Google pouze během již ověřené běžné relace s~heslem správce. Nepropojená identita Google se na přihlašovací stránce odmítá a~nesmí vytvořit ani převzít profil správce. Po propojení se prostřednictvím tlačítka Google může k~profilu přihlásit pouze tato uložená identita Google. Odpojení vyžaduje aktuální heslo správce. Externí přihlášení nenahrazuje přihlášení heslem; pokud konfigurace OAuth nebo úložiště propojení identity chybí či je nelze ověřit, běžná autentizace heslem zůstává nezávislá.

Přihlašovací formulář přijímá uživatelské jméno správce nebo nastavený e-mail pro obnovu. Pokusy o~autentizaci jsou omezovány podle návštěvníka i~normalizovaného zadaného identifikátoru, je-li dostupné schéma omezení četnosti. Tabulka omezení ukládá hashované subjekty místo nezpracovaných adres IP či zadaných identifikátorů. Současná pravidla omezují opakovaná neúspěšná přihlášení v~15minutovém okně a~požadavky obnovy hesla v~delších omezených oknech; úspěšné přihlášení vymaže příslušný stav omezení přihlášení.

Obnova hesla záměrně vrací obecný výsledek požadavku, aby veřejný formulář nepotvrzoval existenci zadaného uživatelského jména či e-mailu. Odkazy obnovy používají úložiště jednorázových hashovaných tokenů a~nastavenou platnost 15 až 1440~minut. Dokončení obnovy změní heslo, označí token jako použitý a~zruší tokeny trvalého přihlášení tohoto účtu. Stránka Účet může odeslat testovací e-mail obnovy prostřednictvím nastaveného PHP \codeword{mail()} nebo transportu SMTP.

Volba \ui{Zůstat přihlášený} používá hashovaný token prohlížeče, aby přihlášení správce přečkalo běžný úklid relací PHP na sdíleném hostingu bez uložení nezpracovaného trvalého přihlašovacího údaje do databáze. Trvalé přihlášení zůstává volitelné. Ruší se příslušnými bezpečnostními postupy účtu a~odhlášení a~může být nedostupné, pokud nelze ověřit stav jeho schématu, aniž by se vypnula aktuální běžná relace PHP.
\subsection{Kontrakt připravenosti databáze}
\label{sec:db-readiness-contract}
\index{s00226@Stav systému!p00188@připravenost databáze}
\index{d00024@databáze!p00187@připravenost}
\index{o00120@odpověď 503}

Funkce závislé na databázovém schématu používají jednotný kontrakt připravenosti. Důležité je rozlišovat schéma, které bylo úspěšně prozkoumáno a~zjištěno jako starší, od schématu, které nyní nelze prozkoumat.

\begin{longtable}{>{\RaggedRight\arraybackslash}>{\hspace{0pt}}p{0.12\textwidth} >{\RaggedRight\arraybackslash}>{\hspace{0pt}}p{0.20\textwidth} >{\RaggedRight\arraybackslash}>{\hspace{0pt}}p{0.35\textwidth} >{\RaggedRight\arraybackslash}>{\hspace{0pt}}p{0.17\textwidth}}
\toprule
\textbf{Stav} & \textbf{Význam} & \textbf{Chování aplikace} & \textbf{Akce vlastníka}\\
\midrule
\endfirsthead
\toprule
\textbf{Stav} & \textbf{Význam} & \textbf{Chování aplikace} & \textbf{Akce vlastníka}\\
\midrule
\endhead
Dostupné & Požadované úložiště bylo nalezeno a~ověřeno. & Běžné čtení a~zápis dané schopnosti. & Žádná.\\
Chybějící & Průzkum uspěl, ale požadovaný objekt chybí. & Může zůstat dostupný dokumentovaný postup pro starší schéma; jinak funkce požaduje migraci. & Zálohujte a~použijte čekající migraci.\\
Neznámé & Schéma nebylo možné spolehlivě prozkoumat. & Chráněné čtení může být bezpečně odmítnuto nebo se může vynechat volitelný výstup pouze pro čtení. Závislé zápisy jsou blokované. & Opravte připojení k~databázi, vybrané schéma nebo oprávnění metadat.\\
Vypnuté & Volitelná funkce je vypnuta. & Funkce se nepoužívá. & Žádná, pokud nemá být funkce povolena.\\
\bottomrule
\end{longtable}

\ui{Stav systému} a~\ui{Diagnostika běhového prostředí} používají tyto stavy jednotně. Výsledek \ui{Chybějící} obvykle ukazuje na potřebu migrace; \ui{Neznámé} je provozní problém a~nesmí být považován za důkaz, že tabulka nebo sloupec nikdy neexistovaly. Diagnostika může ukazovat název schopnosti, stav, ověřené názvy databázových objektů, navržené kontroly a~odkaz na požadavek. Vylučuje text SQL, nezpracované databázové výjimky, přihlašovací údaje, tokeny, cookies a~soukromé cesty souborového systému.

\subsection{Výjimky citlivé na zabezpečení}
\index{o00123@ochrana NSFW}
\index{o00113@obnova hesla!p00188@připravenost databáze}
\index{p00181@přihlášení Google!p00188@připravenost databáze}
\index{o00118@odkazy pro sdílení!p00188@připravenost databáze}

Některé schopnosti mají přísnější nebo konkrétnější chování:

\begin{description}
  \item[Ochrana NSFW.] Pokud průzkum potvrdí starší schéma bez polí NSFW, web zachová chování před zavedením funkce až do migrace. Jsou-li pole \ui{Neznámé}, veřejné požadavky, které by mohly odhalit média s~věkovým omezením, se bezpečně odmítají a~úpravy NSFW jsou pozastaveny.
  \item[Přístup ke galerii.] Databáze se považuje za starší formát před ochranou pouze tehdy, když průzkum uspěje a~chybějí všechna základní pole přístupu. Smíšené uspořádání je neúplnou migrací, proto chráněné trasy zůstávají nedostupné až do opravy. Starší hodnota zveřejnění \codeword{draft} se zapisuje pouze tehdy, je-li známo, že ji definice pole vyžaduje; současná schémata používají \codeword{unpublished}.
  \item[Zrušení odkazu pro sdílení.] Zrušení může používat menší známou sadu polí potřebnou k~vymazání ověřovacího hashe, protože tato operace omezuje přístup. Přesto se zastaví, pokud jsou její vlastní požadovaná pole \ui{Neznámé}.
  \item[Trvalé přihlášení správce.] Chybějící tabulka zapamatovaného přihlášení vypíná trvalé tokeny prohlížeče, ale nebrání běžnému přihlášení heslem vytvořit aktuální relaci PHP.
  \item[Obnova hesla.] Obnova vyžaduje pole e-mailu pro obnovu i~úložiště tokenů obnovy. Chybějící úložiště poskytne pokyny k~migraci; \ui{Neznámé} úložiště vyvolá dočasné selhání místo nesprávného výsledku účtu či tokenu.
  \item[Propojení identity Google.] Externí přihlášení vyžaduje platnou konfiguraci OAuth i~známé úložiště propojení identity. Pokud úložiště chybí nebo je neznámé, přihlášení heslem zůstává nezávisle dostupné.
\end{description}

\subsection{Zápisy, údržba a~přihlašovací údaje}
\index{d00024@databáze!b00014@bezpečnost změn}
\index{n00097@nahrávání!p00188@připravenost databáze}
\index{a00003@aktualizace!p00188@připravenost databáze}

Operace měnící soubory, přihlašovací údaje nebo trvalé záznamy provádějí kontrolu připravenosti před prvním krokem měnícím cíl. Patří sem změny galerií a~fotografií, nahrávání, zápisy WebDAV, migrace galerií, údržba náhledů, oprava a~úklid databáze a~aktivace aktualizace aplikace. Výsledek \ui{Neznámé} zastaví operaci před změnou cíle. Potvrzený výsledek \ui{Chybějící} může použít kompatibilní postup staršího schématu nebo počáteční přípravu migrace pouze tam, kde je toto chování dokumentováno.

Zrušení přihlašovacích údajů následuje stejný omezující princip jako zrušení odkazu pro sdílení: může použít menší známou sadu polí, pokud tato pole postačují k~vypnutí přístupu. Vytvoření nebo důvěřování přihlašovacímu údaji nadále vyžaduje úplné úložiště daného autentizačního postupu.

Je-li zápis blokován jako \ui{Neznámé}, před opakováním opravte průzkum databáze. Přejde-li stav na \ui{Chybějící}, zálohujte a~použijte požadovanou migraci. Přejde-li na \ui{Dostupné}, operaci jednou zopakujte a~zkontrolujte výsledky souborového systému i~databáze. Obnovte existující obnovitelnou migrační či aktualizační úlohu místo zahájení duplicitní úlohy.

\subsection{Volitelné funkce}
\index{m00078@mapy!p00188@připravenost databáze}
\index{h00050@hra s fotografiemi!p00188@připravenost databáze}
\index{h00048@hlasování!p00188@připravenost databáze}
\index{o00126@OpenAI!p00188@připravenost databáze}
\index{s00207@SimBrief!p00188@připravenost databáze}
\index{t00233@telemetrie!p00188@připravenost databáze}

Volitelné prezentační funkce mohou omezit své chování bez vypnutí hlavní galerie, pokud tím nemohou odhalit chráněný obsah nebo udělit přístup. Například mapu lze vynechat, pokud nelze číst její volitelné úložiště, ale uložení přepsání GPS či lightboxu stále vyžaduje známé příslušné pole.

Hlasování i~Hra s~fotografiemi zapisují trvalý stav. Hra zaznamenává zobrazený pár před výběrem vítěze, proto je načtení nového páru také hranicí zápisu a~vyžaduje známé úložiště. Fronty AI, údržba telemetrie, navigační tokeny, uložená data tras SimBrief a~podobné volitelné zápisy následují stejné pravidlo. Funkce může vrátit výsledek pouze pro čtení bez volitelného uložení jen tam, kde rozhraní toto omezení jasně sděluje.

\subsection{Provozní zabezpečení}

Doporučené postupy:

\begin{itemize}
  \item Poskytujte aplikaci prostřednictvím HTTPS.
  \item Omezte přihlašovací údaje databáze na schéma aplikace a~požadované operace.
  \item Uchovávejte \filepath{config.php}, zálohy, mezipaměti a~generované archivy mimo veřejné vyhledávání všude, kde to uspořádání hostingu dovoluje.
  \item Používejte migrace a~aktualizace prostřednictvím ověřených postupů údržby.
  \item Kontrolujte auditní protokoly a~bezpečnostní diagnostiku.
  \item Testujte chráněné galerie prostřednictvím anonymních relací návštěvníků.
  \item Udržujte souběžné zálohy souborového systému a~databáze.
\end{itemize}
\section{Vzhled, prezentace a~lokalizace}
\label{sec:theme}
\index{v00268@vzhled}
\index{l00076@lokalizace}

Oblast Vzhled řídí identitu webu, barvy, rozměry, šířku stránky, pozadí, vlastní značku, podklady ikony webu, jazyk, prezentaci karet galerií, výchozí nastavení lightboxu, prezentaci GPS, oblíbenou navigaci a~veřejné vykreslování náhledů. Nastavení používají normalizované skalární hodnoty nebo cílené modely služeb.

Editor Vzhledu nepoužívá samostatný přepínač světlého a~tmavého režimu. Veřejnou barevnou paletu přímo určují uložené barevné volby a~volitelný motiv Vlastního CSS. Nastavuje také rozvržení karet galerií a~podgalerií a~umožňuje nastavit až tři zkratky v~záhlaví. Každá může cílit na hlavní stránku, jednu galerii nebo zůstat prázdná. Anonymní návštěvníci vidí pouze nastavené zkratky galerií, které jsou stále veřejné a~uvedené v~seznamech; duplicitní nebo odstraněné galerie se při uložení nastavení vyřadí. Ponechání všech tří pozic prázdných skryje tlačítka oblíbených galerií. Karty galerií mohou zobrazovat také nastavitelný odznak počtu obsažených fotografií. Je-li odznak pro galerii účinný, otevření galerie zachovává stejný počet větve viditelný v~úvodním panelu: počítá obrázky v~aktuální galerii a~jejích přístupných podgaleriích podle stejných pravidel veřejnosti a~přístupu jako karta. Prezentace jednotlivé galerie může dědit výchozí hodnoty webu nebo přepsat volby nabízené editorem galerie.

\subsection{Přehled ovládacích prvků Vzhledu}
\index{v00268@vzhled!o00131@ovládací prvky}
\index{s00227@stránkování!n00100@nastavení Vzhledu}
\index{i00054@ikona webu}
\index{v00253@vlastní CSS}

Současná stránka Vzhled je rozdělena na \ui{Vzhled}, \ui{Značka a~média}, \ui{Rozvržení}, \ui{Jazyk} a~\ui{Vlastní CSS}. Následující přehled zaznamenává ovládací prvky viditelné vlastníkovi a~normalizaci vynucovanou aktuálním sestavením.

\begin{longtable}{>{\RaggedRight\arraybackslash}>{\hspace{0pt}}p{0.20\textwidth} >{\RaggedRight\arraybackslash}>{\hspace{0pt}}p{0.27\textwidth} >{\RaggedRight\arraybackslash}>{\hspace{0pt}}p{0.41\textwidth}}
\toprule
\textbf{Oblast} & \textbf{Ovládací prvek} & \textbf{Současné chování}\\
\midrule
\endfirsthead
\toprule
\textbf{Oblast} & \textbf{Ovládací prvek} & \textbf{Současné chování}\\
\midrule
\endhead
Vzhled & Identita webu a~barvy & Název webu a~barvy zvýraznění, tmavého zvýraznění, pozadí stránky, panelu, panelu otevřené galerie, názvu v~záhlaví a~názvu galerie. \uv{Tmavé zvýraznění} je barva při najetí či obrysu, nikoli volba tmavého režimu.\\
Vzhled & Rohy a~písmo & Poloměr rohů je 0--32~pixelů. Režim písma je \ui{Klasické patkové} nebo \ui{Čisté bezpatkové}. Živý náhled se při úpravě těchto prvků aktualizuje.\\
Vzhled & Šířka veřejné stránky & \ui{Výchozí}, \ui{Širší}, \ui{Plná šířka} nebo \ui{Vlastní}. Vlastní šířka je omezena na 1024--2048~pixelů a~bez platné hodnoty je výchozí 1440.\\
Vzhled & Značka GPS & Je-li funkce map GPS povolena, značky na kartách fotografií a~jejich podklad lze zobrazit či skrýt samostatně. Velikost značky je 14--48~pixelů, výchozí 26. Posuvník podkladu se zobrazuje jako 0--48~pixelů, ale současný serverový normalizátor považuje 0 nebo jinou nekladnou hodnotu za nenastavenou a~obnoví 22~pixelů; pozadí skryjte samostatným zaškrtávacím polem podkladu.\\
Vzhled & Veřejné stránky štítků & Vlastní sloupce a~řádky stránek štítků používají stejné bezpečné maximální hodnoty jako jiné veřejné mřížky: 1--12 sloupců a~1--50 řádků. Karty výsledků štítků mohou používat svislý nebo vodorovný návrh karty galerie nezávisle na běžných výpisech galerií.\\
Vzhled & Štítky úvodního panelu galerie & Řazení podle využití nebo abecedy. Buď zobrazit všechny ihned, nebo zobrazit 1--200 štítků před ovládáním rozbalení na místě; výchozí práh je 20. Volitelné posouvání používá 1--12 viditelných řádků, výchozí 5, a~aktivuje se jen tehdy, když vykreslené štítky skutečně přesáhnou nastavené řádky.\\
Rozvržení & Zkratky záhlaví & Tři seřazené pozice. Pozice může být prázdná, odkazovat na hlavní stránku nebo na galerii nalezenou výběrem galerií. Veřejní návštěvníci vidí pouze veřejné cílové galerie uvedené v~seznamech.\\
Rozvržení & Návrh karty galerie & Globální výchozí hodnota je \ui{Svislý systém} nebo \ui{Vodorovný systém}. Fyzická galerie může tuto hodnotu dědit nebo přepsat, je-li dostupné schéma prezentace.\\
Rozvržení & Odznak počtu fotografií & Globální odznak obsažených fotografií je ve výchozím stavu povolen. Jednotlivé galerie jej mohou dědit nebo uložit přepsání. Je-li povolen pro otevřenou galerii, úvodní panel také ukazuje počet obrázků její větve včetně přístupných podgalerií.\\
Rozvržení & Veřejný vykreslovač náhledů & \ui{Responzivní výběr prohlížečem} je výchozí; \ui{Postupné zostřování náhledů} je volitelný režim beta pro karty fotografií v~otevřené galerii. Titulní obrázky galerií a~koláže nadále používají responzivní výběr prohlížečem.\\
Rozvržení & Stránkování & Globální stránkování je ve výchozím stavu vypnuté. Sloupce mají výchozí hodnotu 3 a~omezení 1--12; řádky mají výchozí hodnotu 3 a~omezení 1--50. Účinný počet položek na stránku je součin sloupců a~řádků.\\
Rozvržení & Mřížka hlavní stránky & Úvodní stránka má vlastní mřížku s~1--12 sloupci a~1--50 řádky. Nemění mřížky uvnitř galerií.\\
Rozvržení & Obnovení mřížek galerií & \ui{Obnovit všechny vlastní mřížky galerií} vymaže databázová přepsání, obnoví výchozí stav dědění podgalerií a~odstraní odpovídající klíče mřížky z~doprovodných souborů \filepath{gallery.json}, aby pozdější vyhledání nemohlo znovu importovat zastaralá nastavení mřížek.\\
Rozvržení & Výchozí lightbox & \ui{Jeden obrázek}, \ui{Pás fotografií} nebo \ui{3D karusel}, je-li funkce režimů lightboxu povolena. Klasický režim jednoho obrázku je náhradou a~jediným zpřístupněným režimem při vypnuté funkci.\\
Značka & Banner veřejného záhlaví & Volitelný banner JPG/PNG/GIF/WebP na úrovni Vzhledu, nejvýše 8~MB. Nahrazuje viditelný název webu, pokud není nastaven banner konkrétní galerie.\\
Značka & Oddělovač záhlaví & Volitelný oddělovač JPG/PNG/GIF/WebP, nejvýše 8~MB. Šířka 0 následuje responzivní šířku stránky; kladná šířka je omezena na 160--3840~pixelů. Výška má výchozí hodnotu 72 a~omezení 8--512~pixelů. Normální režim zachovává poměr stran a~výšku považuje za maximum; roztažení přesně používá nastavený rámec a~může obrázek záměrně deformovat.\\
Značka & Ikona webu & Nahrajte JPG/PNG/GIF/WebP, přetáhněte čtvercový výřez a~zvětšete jej 1--3krát. Uložení vytváří varianty PNG ikony webu o~32, 48 a~180~pixelech.\\
Značka & Pozadí Vzhledu & Zachovává původní nahraný soubor a~volitelně generuje odvozeninu WebP, je-li dostupná podpora GD/WebP. Optimalizovaná nejdelší strana je 1024--3840~pixelů, výchozí 1920, s~kroky posuvníku 128~pixelů. Lze ji obnovit či odstranit bez odstranění originálu.\\
Značka & Viditelnost a~náhrada pozadí & Viditelnost pozadí je 0--100~procent, výchozí 65. Náhradním zdrojem galerie může být žádný, nový soubor nahraný do Vzhledu, existující obrázek galerie nebo koláž z~veřejných galerií. Samostatné akce obnovují všechny volby pozadí galerií, odstraňují pozadí Vzhledu nebo ikonu webu.\\
Vlastní CSS & Přednastavený motiv nebo nahrání & Předvolby jsou soubory CSS nyní přítomné v~\filepath{custom_css/}; výběr zkopíruje soubor do \filepath{public/assets/custom.css}. Ručně nahraný soubor \filepath{.css} nahrazuje tento aktivní soubor. Aktivní vlastní styl se načítá po vestavěném CSS a~uložených volbách Vzhledu.\\
Vlastní CSS & Akce obnovení & \ui{Obnovit vlastní CSS} odstraní aktivní vlastní styl a~vymaže značku předvolby. \ui{Obnovit podle CSS} vymaže uložená vizuální přepsání palety, písma, poloměru, zdroje pozadí a~GPS, aby se znovu uplatňovalo podkladové CSS; strukturální nastavení šířky stránky se záměrně zachovávají.\\
Jazyk & Jazyk administrace a~veřejnosti & Jazyk administrace se ukládá pro relaci a~prohlížeč správce; veřejný jazyk je anonymní výchozí hodnota celého webu. Výběr návštěvníka je nezávislé přepsání popsané níže.\\
Jazyk & Údržba balíčků & Podporované balíčky ukazují počty řetězců a~pokrytí. Editor JSON ověřuje objekt s~textovými hodnotami, importuje a~exportuje JSON a~může vymazat diagnostiku chybějících klíčů z~ověřených relací. Angličtina zůstává zdrojovým a~náhradním katalogem.\\
\bottomrule
\end{longtable}

\subsection{Režimy procházení lightboxu}
\index{l00075@lightbox!r00193@režimy procházení}
\index{p00133@pás fotografií}
\index{a00000@3D karusel}

Editor Vzhledu určuje výchozí veřejnou prezentaci lightboxu. \ui{Jeden obrázek} zachovává klasický prohlížeč. \ui{Pás fotografií} přidává pod aktivní fotografii náhledy blízkých fotografií. \ui{3D karusel} umísťuje malou sadu sousedních fotografií za aktivní fotografii a~poskytuje při procházení více vizuálních souvislostí.

Fyzická galerie může dědit výchozí hodnotu Vzhledu nebo uložit vlastní přepsání režimu lightboxu. Prezentace chytré galerie může také určit účinný režim lightboxu ze svého nastavení prezentace a~výchozích hodnot webu. Režimy sdílejí stejný autorizovaný postup metadat a~navigace lightboxu, takže celá obrazovka, prezentace, přiblížení, mapy, hlasování a~postupná kvalita nadále následují běžná pravidla. Vypnutí globální funkce režimů lightboxu skryje tyto ovládací prvky procházení bez změny uložených fotografií.

\subsection{Výběr jazyka rozhraní}
\label{sec:language-selection}
\index{j00058@jazyk}
\index{l00076@lokalizace!v00257@výběr jazyka}
\index{p00175@překladové balíčky}

PHP Gallery odděluje jazyk rozhraní programu od textu, který píšete sami. Změna jazyka rozhraní mění tlačítka, položky nabídek, nápovědu, stavové zprávy, dialogy a~další formulace aplikace. Názvy a~popisy galerií a~fotografií mohou mít volitelné verze v~udržovaných jazycích, ale upravují a~ukládají se odděleně; štítky a~další obsah vlastníka zůstávají beze změny.

Otevřete \ui{Vzhled > Jazyk}. Jsou dostupné dva nezávislé výběry:

\begin{itemize}
  \item \ui{Jazyk rozhraní administrace} mění jazyk správce v~aktuálním prohlížeči.
  \item \ui{Jazyk veřejných návštěvníků} nastavuje výchozí jazyk rozhraní veřejných stránek.
\end{itemize}

Volitelný \ui{Výběr jazyka návštěvníka} umožňuje jednotlivým návštěvníkům přepsat veřejnou výchozí hodnotu ve vlastním prohlížeči. Správce vybírá, které udržované jazyky se ve veřejném výběru objeví; vypnutí funkce vrací návštěvníky k~veřejnému jazyku celého webu. Jazyky administrace a~veřejnosti se nemusejí shodovat.

Veřejné stránky používají kompaktní jazykový prvek s~názvy v~daných jazycích a~místně poskytovanými ikonami vlajek. Volba jazyka ponechává návštěvníka na stejné veřejné stránce a~pamatuje si volbu pro pozdější návštěvy v~tomto prohlížeči. \ui{Použít výchozí jazyk webu} odstraní osobní přepsání. Odpovídající parametr \codeword{?lang=} zůstává dostupný pro přímé odkazy a~použití bez JavaScriptu.

\subsubsection{Přizpůsobení výběru návštěvníka}

Výběr má pět předvoleb: \ui{Klasický}, \ui{Plné pilulky}, \ui{Obrysy}, \ui{Jemné karty} a~\ui{Minimální}. \ui{Nastavení > Obecné} poskytuje základní volby předvolby a~vlajek, zatímco \ui{Vzhled > Jazyk} obsahuje podrobné návrhové prvky a~živý náhled. Tato nastavení vzhledu nemění nabízené jazyky, výchozí jazyk webu, jazyk administrace ani osobní volbu návštěvníka.

\subsubsection{Vícejazyčný text galerií a~fotografií}
\index{l00076@lokalizace!o00117@obsah galerie}
\index{l00076@lokalizace!n00107@názvy fotografií}

Editory galerií a~fotografií zachovávají existující název a~popis jako zdrojový obsah. Volbou \ui{Jazyk aktuálního názvu a~popisu} označte text jako anglický, český, německý či švédský nebo ponechte \ui{Neuvedeno}. Existující instalace se automaticky neoznačují jako anglické.

Ovládacím prvkem zeměkoule otevřete \ui{Další jazyky} a~zadejte volitelné překlady názvů a~popisů. Názvy a~popisy galerií používají náhradní zdroj nezávisle. Uložený jazykový řádek fotografie se naopak považuje za jednu variantu popisku: prázdná pole zůstávají skrytá místo míchání zdrojového jazyka pod přeloženým textem fotografie. Ponechání obou přeložených polí prázdných odstraní jazykový řádek. Ukládání probíhá běžným formulářem galerie či fotografie a~zůstává ve stávajícím postranním panelu administrace.

Veřejné stránky používají udržovaný jazyk vybraný návštěvníkem. Odpovídající text se ukazuje na kartách, v~popisech galerií, překryvech fotografií, metadatech lightboxu, výsledcích vyhledávání, alternativním textu a~metadatech SEO. Chybějící pole galerií používají zdrojový obsah; přítomná varianta popisku fotografie zobrazuje pouze uložená přeložená pole. Překlady nemění URL, identifikátory v~URL, složky, názvy souborů, pořadí, viditelnost, hesla, pravidla NSFW ani oprávnění médií.

Je-li nakonfigurována textová pomoc OpenAI, \ui{Navrhnout překlad pomocí OpenAI} odešle zdrojové pole a~vloží kontrolovatelný návrh do cílového pole. Návrh se nikdy automaticky neukládá ani nezveřejňuje. Ruční překlad zůstává plně dostupný bez poskytovatele.

\subsubsection{Dostupné jazyky a~jejich současný stav}
\index{j00058@jazyk!a00005@angličtina}
\index{j00058@jazyk!c00021@čeština}
\index{j00058@jazyk!n00109@němčina}
\index{j00058@jazyk!s00232@švédština}

Angličtina je kanonickým zdrojovým jazykem, výchozím jazykem nové instalace a~náhradním jazykem. Čeština, němčina a~švédština jsou také udržovány jako úplné katalogy. Tyto čtyři jazyky jsou nyní dostupné ve výběru rozhraní.

\begin{longtable}{>{\RaggedRight\arraybackslash}>{\hspace{0pt}}p{0.20\textwidth} >{\RaggedRight\arraybackslash}>{\hspace{0pt}}p{0.10\textwidth} >{\RaggedRight\arraybackslash}>{\hspace{0pt}}p{0.58\textwidth}}
\toprule
\textbf{Jazyk} & \textbf{Kód} & \textbf{Současná role}\\
\midrule
\endfirsthead
\toprule
\textbf{Jazyk} & \textbf{Kód} & \textbf{Současná role}\\
\midrule
\endhead
Angličtina & \codeword{en} & Úplný kanonický zdroj, výchozí jazyk nové instalace a~náhradní jazyk.\\
Čeština & \codeword{cs} & Úplný udržovaný český překlad dostupný ve výběru.\\
Němčina & \codeword{de} & Úplný udržovaný německý překlad dostupný ve výběru.\\
Švédština & \codeword{sv} & Úplný udržovaný švédský překlad dostupný ve výběru.\\
\bottomrule
\end{longtable}

\important{Nyní lze vybírat pouze angličtinu, češtinu, němčinu a~švédštinu. Pokud udržovaný překlad chybí, náhradou zůstává angličtina.}

\subsubsection{Jak fungují překladové balíčky}
\index{p00175@překladové balíčky!n00096@náhradní jazyk}
\index{p00175@překladové balíčky!j00060@JSON}
\index{p00175@překladové balíčky!d00028@diagnostika}

Překladový balíček je jednoduše slovník pojmenovaných zpráv rozhraní. Program požádá o~stabilní klíč zprávy, například \uv{uložit galerii} nebo \uv{odstranit vybrané fotografie}, a~aktivní jazykový balíček dodá text k~zobrazení. Interní klíč zůstává ve všech jazycích stejný; mění se pouze zobrazený text.

Tabulka \ui{Rozpoznané jazykové balíčky} ukazuje pokrytí podporovaných jazyků a~pododdíl \ui{Diagnostika} zaznamenává chybějící klíče pozorované během ověřené relace administrace. Angličtina dodává náhradní formulaci, pokud v~udržovaném překladu chybí klíč.

\ui{Editor balíčků} je určen pro pečlivou údržbu nebo import připraveného překladu. Jazyková data se ukládají jako JSON. Text na levé straně každého páru JSON je v~editoru stabilním klíčem a~nesmí se překládat. Text na pravé straně je hodnota, kterou lze překládat. Zástupné značky jako \codeword{\{count\}}, \codeword{\{gallery\}} nebo \codeword{\{time\}} představují živé hodnoty vkládané PHP Gallery a~musejí zachovat stejná jména ve všech jazycích.

Jazykový balíček lze také exportovat jako JSON, přeložit mimo PHP Gallery a~znovu importovat. Import se přijme pouze tehdy, je-li soubor objektem JSON obsahujícím textové hodnoty, takže překladatel může pracovat na balíčku bez přístupu k~serveru galerie.

Pro vývojáře a~správce jsou udržované katalogy v~\filepath{app/lang/}. JSON je kanonický upravitelný formát. Slovníky PHP pro angličtinu, češtinu, němčinu a~švédštinu zůstávají pouze kompatibilními náhradami pro instalace bez odpovídajícího souboru JSON. Serverem vykreslené stránky i~JavaScript prohlížeče používají stejný model aktivního jazyka s~anglickou náhradou, takže překlad není nutné duplikovat v~samostatných katalozích frontendu a~backendu.

Vlastní CSS podporuje vizuální doladění konkrétní instalace. CSS generované Vzhledem zpřístupňuje ověřené proměnné a~vypočítané selektory prostřednictvím vyhrazené trasy. Vlastní značka galerie může přepsat vybrané podklady webu a~zachovat účinné náhradní chování.

\subsection{Štítky úvodního panelu galerie}
\index{s00230@štítky!u00246@úvodní panel galerie}
\index{v00268@vzhled!s00231@štítky galerie}

Pododdíl \ui{Vzhled > Štítky galerie} řídí kompaktní sbírku štítků pod názvem otevřené galerie. Oblast štítků využívá celou šířku úvodního panelu, takže se štítky mohou zalamovat až k~pravému okraji místo omezení běžnou čitelnou šířkou odstavce. Přímé štítky galerie a~štítky shromážděné z~obsaženého obsahu zůstávají samostatnými významovými skupinami, ale každá skupina následuje stejná globální pravidla zobrazení.

Stejný pododdíl také řídí veřejnou cílovou stránku štítku. \ui{Galerie na řádek} a~\ui{Řádky na stránku} určují mřížku galerií výsledků štítku a~kapacitu stránky, zatímco \ui{Návrh karty galerie} vybírá svislou nebo vodorovnou prezentaci karet těchto výsledků. Hodnoty ovlivňují veřejné stránky otevřené přes štítek; nemění běžné stránky galerií. Nejsou-li vlastní hodnoty uloženy, náhradou zůstává globální mřížka a~návrh karet Vzhledu. Globální přepínač stránkování nadále určuje, zda se dlouhý výsledek štítku rozdělí na stránky. Z~\ui{Upravit štítky} volba \ui{Nastavit zobrazení štítků} otevírá přímo tento pododdíl a~\ui{Spravovat metadata štítků} poskytuje zpětný kontextový odkaz.

Ve výchozím stavu prohlížeč nejprve ukáže 20 štítků. Limit lze změnit od 1 do 200 posuvníkem nebo přesným číselným polem. Existuje-li více štítků, objeví se \ui{Zobrazit všechny štítky} a~pomocí JavaScriptu rozbalí existující prvky na místě; nepoužívá se opětovné načtení stránky ani dodatečný požadavek serveru. Stejný prvek může sbírku znovu sbalit. Správci, kteří nechtějí postupné odhalování, mohou povolit \ui{Zobrazit každý štítek okamžitě}. Všechny štítky jsou ve všech režimech nadále součástí serverem vykresleného HTML, takže vypnutí JavaScriptu ponechá úplnou sbírku viditelnou a~použitelnou místo trvalého skrytí obsahu.

Řazení štítků lze nastavit jako \ui{Nejpoužívanější první} nebo \ui{Abecedně}. Výchozí je \ui{Nejpoužívanější první}. Využití je celkový počet přímých přiřazení štítku galeriím a~fotografiím v~celé instalaci. Shodné počty využití se rozlišují přirozeným abecedním pořadím bez rozlišení velikosti písmen a~poté stabilním identifikátorem. Přímé štítky galerie se řadí uvnitř své skupiny a~obsažené štítky uvnitř své; skupiny se nemíchají.

Dlouhé sbírky štítků mohou také používat responzivní vnitřní posuvník. Tato volba je ve výchozím stavu povolena a~má pět viditelných řádků. Práh řádků lze nastavit od 1 do 12 posuvníkem nebo přesným číselným polem. Prohlížeč změří řádky po skutečném zalomení štítků při aktuální šířce úvodního panelu a~povolí posouvání jen tehdy, když naměřený počet přesáhne nastavený práh. Změna velikosti stránky vyvolá nové měření. Vypnutí posuvníku odstraní omezení výšky a~umožní přirozený růst úvodního panelu.
\section{Stahování, sdílení a~přenos}
\label{sec:sharing}
\index{s00224@stahování}
\index{z00279@ZIP}
\index{w00269@WebDAV}

Stahování galerií připravuje povolený obsah jako archivy ZIP a~průběžně odesílá dokončený archiv prohlížeči. Běžné stažení galerie prochází povolený podstrom galerií; stahování chytré galerie znovu vyhodnocuje uložený dotaz a~aktuální pravidla přístupu; stažení výběru Správce fotografií balí pouze aktuálně vybrané fotografie s~ověřeným vlastnictvím; správce může také požádat o~archiv všech galerií. Názvy položek archivu se normalizují, aby bránily procházení mimo určenou cestu a~zachovávaly deterministické cesty.

V~moderním prohlížeči stahování veřejné či chytré galerie nejprve vyžádá soukromý omezený manifest. Prohlížeč poté žádá o~každý povolený originál nezávislým koncovým bodem médií a~sestavuje ZIP místně, ukazuje průběh a~umožňuje zrušení či opakování bez udržování jediného dlouhého archivačního požadavku PHP. Každý požadavek zdroje opakuje aktuální kontroly galerie, příslušnosti k~chytré galerii, viditelnosti návštěvníka, verze média a~umístění uvnitř povolené cesty; samotný identifikátor obrázku nikdy neuděluje přístup. Limity počtu souborů, souhrnných bajtů zdrojů, duplicitních názvů, paměti a~ZIP64 udržují postup omezený. Přímé odkazy a~klienti bez JavaScriptu zachovávají stávající náhradní ZIP na serveru.

Generované ZIP galerií, všech galerií a~výběrů používají podpis obsahu odvozený z~příslušných metadat galerií a~obrázků, takže stále aktuální archiv z~mezipaměti lze použít znovu a~změněná sbírka vytváří novou identitu mezipaměti. Tento \emph{podpis obsahu} jsou metadata zneplatnění mezipaměti, nikoli veřejný autorizační token HMAC pro stahování. Oprávnění před poskytnutím archivu kontroluje řadič stahování.

\subsection{Migrace mezi galeriemi}
\index{m00088@migrace galerie}
\index{m00087@migrace API}

Oblast API editoru galerie podporuje oba směry přenosu mezi kompatibilními instalacemi PHP Gallery. \ui{Exportovat tuto galerii do jiné instance} považuje aktuální galerii za zdroj a~používá klíč API omezený na přijímající nadřazenou galerii. \ui{Importovat do této galerie z~jiné instance} považuje aktuální galerii za přijímajícího rodiče a~používá klíč API vytvořený pro vzdálenou zdrojovou galerii. Import vytváří novou podřízenou galerii pod tímto rodičem; nikdy nepřepisuje ani nepřesouvá vybranou přijímající galerii.

\ui{Zahrnout podgalerie} je ve výchozím stavu povoleno. Při výběru verzovaný manifest popisuje zdrojovou galerii a~úplnou hierarchii potomků v~pořadí s~rodiči prvními. Cíl znovu vytvoří tuto hierarchii pod novým importovaným kořenem a~zachovává podporovaná metadata galerií a~obrázků, překlady, stav letové mapy, originály, existující náhledy a~podklady značky galerie. Vypnutí volby přenáší pouze vybranou zdrojovou galerii. Současná pravidla kompatibility vyžadují přesně stejnou verzi aplikace na zdroji i~cíli.

Přijímající server vytváří deterministický plán balíčků omezený svými limity nahrávání. Každý originál zůstává seskupený s~již generovanými náhledy, zatímco podklady galerie lze balit nezávisle. Položky ZIP používají stabilní klíče podkladů, ukládají existující bajty bez opětovné komprese fotografií a~před instalací se kontrolují proti velikostem a~kontrolním součtům SHA-256 manifestu. Přesný rozsah klíče API, struktura stromu galerií, cílové cesty, připravenost databáze a~úplnost podkladů se ověřují na příslušných hranicích změn.

Správce vybírá interval opětovného připojení od 5 do 300~sekund. Po vypršení časového limitu, obnovení stránky nebo přerušení spojení se aktivní strana dotáže trvalé cílové úlohy, které klíče podkladů již jsou přítomné, a~přeskočí dokončené balíčky místo slepého opětovného přenosu. Mapování složek a~přijaté podklady se ukládají průběžně, aby mohl stejný přenos bezpečně pokračovat. Dokončení zůstává samostatným krokem a~odmítá se, dokud nejsou přítomné všechny deklarované podklady. Zrušení postupu v~prohlížeči zastaví další požadavky; přihlašovací údaje přesto zrušte, jakmile dočasný migrační klíč již není potřeba.

\subsection{Mobilní přenos WebDAV}
\index{w00269@WebDAV!m00089@mobilní nahrávání}

\ui{Galerie > Mobilní nahrávání} vytváří připojení kompatibilní s~WebDAV pro jednu vybranou galerii. Zadejte rozpoznatelné označení, vyberte galerii, vytvořte připojení a~zkopírujte vygenerovanou URL serveru, uživatelské jméno a~heslo do mobilní aplikace. Heslo se zobrazuje pouze jednou a~ukládá jako hash. Existující řádky ukazují čas posledního použití a~lze je ihned odstranit z~administrace.

V~klientech, například PhotoSync, použijte typ cíle WebDAV. Současný serverový příjem přijímá JPG, PNG, GIF a~WebP, proto tam, kde je to dostupné, povolte v~klientovi převod HEIC do JPEG. Nahrávání WebDAV neobchází pravidla příjmu galerie: přihlašovací údaj je omezen na galerii, každý zápis vyžaduje autentizaci, schéma galerií a~obrázků musí být známé a~přijatá média se následně skenují, volitelně automaticky přejmenují a~zpracují běžným postupem náhledů.

\subsection{Veřejné vyhledávání obsahu a~nástroje vybraných fotografií}
\index{r00194@robots.txt}
\index{s00208@sitemap.xml}
\index{o00116@obrazová mapa webu}
\index{s00221@správce fotografií}

Veřejná odpověď \codeword{/robots.txt} popisuje hranice pro roboty, zatímco \codeword{/sitemap.xml} uvádí způsobilé veřejné URL galerií a~obrázků s~metadaty obrázků a~informacemi o~poslední změně. Nezveřejněný, soukromý, heslem chráněný, věkově omezený nebo jinak nepovolený obsah se těmito výstupy nezveřejňuje. Kanonické URL, sociální metadata a~strukturovaná data obrázků používají stejný model veřejných cest zohledňující přístup.

Správce fotografií je překryv veřejného zobrazení galerie pro přihlášeného správce, nikoli funkce návštěvníka. Správce může vybrat fotografie a~fyzické podgalerie, odstranit smíšený výběr nebo jej zkopírovat do nové fyzické galerie pod vybraným rodičem. Výběry pouze fotografií lze také přetáhnout na viditelnou podřízenou galerii, přesunout či kopírovat prohledávatelným výběrem cíle, sdílet nativní nabídkou zařízení, je-li podporována, nebo stáhnout. Nativní sdílení přechází na ZIP vybraných fotografií, pokud prohlížeč nemůže sdílet připravené soubory. Chytré galerie jsou vyloučeny z~operací fyzických podstromů. Samotný výběr je jen stavem prohlížeče a~neuděluje oprávnění. Každý požadavek změny či ZIP znovu kontroluje běžně přihlášený účet, token CSRF, zdrojovou galerii, vlastnictví výběru a~případný cíl. Běžní návštěvníci tyto prvky nedostávají. Správce fotografií je záměrně nezávislý na schopnosti Přímá správa: vypnutí Přímé správy odstraní prvky úprav, přidávání, odstraňování a~pořadí přímo na veřejné stránce, zatímco nezávisle povolený Správce fotografií může nadále poskytovat výběr, odstranění, přesun, kopírování, vytváření fyzických galerií a~export výběru.
\section{Údržba, aktualizace a~obnova}
\label{sec:maintenance}
\index{u00238@údržba}
\index{a00003@aktualizace}
\index{z00275@záloha}

\subsection{Migrace}

Změny schématu jsou v~souborech PHP s~časovým razítkem pod \filepath{database/migrations/}. Migrační nástroj předem ověřuje čekající definice, používá je v~pořadí názvů souborů a~zaznamenává úspěšné verze. Nové změny schématu dostávají nový migrační soubor, aby existující instalace zachovávaly deterministickou historii.\footnote{Některé migrace kromě SQL zahrnují opravné callbacky. Testy kompatibility chrání scénáře částečných aktualizací, v~nichž starší migrační nástroj narazí na novější definice migrací.}

Kontroly připravenosti databáze se pamatují pouze pro aktuální požadavek PHP. To běžně snižuje počet opakovaných dotazů na metadata, ale migrace může změnit odpověď během stále běžícího procesu administrace, instalátoru, aktualizátoru nebo příkazového řádku. Migrační nástroj proto maže zapamatované odpovědi po příkazech měnících schéma a~po úspěšných opravných callbacích. Kontrola po migraci pak čte nový databázový katalog místo použití odpovědi z~doby před migrací. Běžné migrační příkazy měnící pouze data mezipaměť připravenosti nemažou, protože nemohou vytvořit ani odstranit tabulku, pole či index.

Verze~0.106 přidala tabulku \codeword{maintenance\_jobs} běžnou migrací. Pokud dosud nebyla použita, spusťte ji obvyklým způsobem ze stránky správy galerií. Vyžaduje běžné oprávnění instalace vytvářet tabulky, ale nevyžaduje databázové triggery, uložené rutiny, \codeword{SUPER}, změnu globální serverové proměnné ani samostatný zásah správce databáze. Dokud migrace není dostupná a~ověřená, Centrum údržby odmítá zahájit úlohu místo pokusu o~údržbu bez trvalých kontrolních bodů. Migrace kooperativních galerií a~identity náhledů používají stejný běžný postup aktualizace.

\subsection{Jak galerie kontroluje připravenost databáze}
\index{d00024@databáze!p00187@připravenost}
\index{p00166@průzkum schématu}
\index{m00086@migrace!h00052@chybějící objekty}
\index{s00226@Stav systému!p00165@průzkum databáze}
\index{r00202@řešení problémů!p00186@připojení k databázi}

Společné stavy připravenosti a~akce vlastníka definuje oddíl~\ref{sec:db-readiness-contract}. Údržba používá stejný kontrakt při kontrole tabulek, polí a~indexů potřebných pro obnovu hesla, ochranu NSFW, automatizaci nahrávání, metadata náhledů a~další funkce závislé na schématu.

Během upgradu může být stav \ui{Chybějící} očekávaný, pokud nové soubory aplikace dorazí před spuštěním jejich migrace. \ui{Neznámé} naopak znamená, že PHP Gallery nemohlo spolehlivě prozkoumat vybranou databázi. Při \ui{Chybějící} zálohujte a~použijte čekající migraci. Při \ui{Neznámé} před změnou schématu zkontrolujte připojení k~databázi, vybrané schéma a~oprávnění průzkumu metadat.

Po opravě obnovte \ui{Stav systému} a~znovu otestujte ovlivněnou funkci. Při žádosti o~pomoc poskytovatele hostingu nebo správce uchovejte případný odkaz na požadavek společně s~časem a~verzí aplikace.

\subsection{Koš galerií a~obnova}
\index{g00040@galerie!k00071@koš}
\index{g00040@galerie!o00112@obnova}
\index{k00071@koš!a00011@automatické mazání}

Otevřete \ui{Údržba > Koš} pro kontrolu obnovitelných odstranění galerií. Každý aktivní řádek ukazuje původní cestu, čas odstranění, počty galerií a~fotografií, uloženou velikost, stav životního cyklu a~termín automatického vymazání, pokud jsou tato pravidla aktivní. Problematické položky zůstávají viditelné ke kontrole; PHP Gallery neodhaduje, že nejednoznačný stav souborového systému a~databáze lze bezpečně zničit.

Hlavní přepínač Koše řídí budoucí odstranění galerií a~je ve výchozím stavu povolen. Vypnutí neskrývá ani nemaže existující položky: obnova a~ruční trvalé odstranění zůstávají dostupné. Automatické mazání je samostatná dobrovolná volba, ve výchozím stavu vypnutá. Uchování má výchozí dobu 30~dnů, dávka údržby je omezená a~automatické odstranění se pozastavuje, kdykoli je Koš vypnutý. Zapnutí automatického mazání po pozastavení poskytne aktuálním obnovitelným položkám nové úplné období uchování před možným začátkem plánovaného odstranění.

Migrace \filepath{202609070001_gallery_trash_bin.php} a~\filepath{202609070002_gallery_trash_state_machine.php} přidávají trvalé záznamy operací. Je-li Koš povolen, odstranění kontroluje schéma před prvním přesunem v~souborovém systému či odstraněním databázového řádku. Potvrzené chybějící nebo neznámé úložiště zastaví obnovitelné odstranění a~odkáže správce k~použití či kontrole migrací; tiše nepřejde na trvalé odstranění. Plánovaná údržba webu konzervativně uvádí do souladu zastaralé rozpracované operace a~spouští vypršelé automatické mazání pouze tehdy, když jsou povoleny Koš i~automatické mazání.

Obsah Koše je v~chráněném trvalém úložišti \filepath{data/gallery-trash/} místo mezipaměti nebo aktivního stromu galerií. Balíčky nasazení obsahují soubor pravidel přístupu adresáře, ale vylučují uložený obsah Koše. Zahrňte tento běhový adresář a~databázi do plánování záloh a~obnovy; ani jedna část samostatně netvoří úplnou obnovitelnou položku.

\subsection{Centrum údržby}
\index{c00017@centrum údržby}
\index{u00238@údržba!c00016@centrální postup}

Otevřete \ui{Údržba > Centrum údržby} pro řízený postup \ui{Analyzovat}, \ui{Zkontrolovat}, \ui{Provést} a~\ui{Ověřit} napříč oblastmi běžné údržby. Analýza je pouze pro čtení obsahu galerií, telemetrie, protokolů, Koše, mezipaměti, náhledů a~databáze; zapisuje jen trvalý kontrolní bod Centra údržby potřebný k~pokračování průzkumu. Kontrola ukazuje povinné, výchozí a~volitelné úkoly před zahájením vybraného úklidu. Hloubkové ověření médií je volitelné a~nevybírá se automaticky.

Zkontrolovaný plán je svázán s~aktuální verzí aplikace, revizí použitých migrací, revizí registru Centra údržby a~obsahem plánu. Změní-li se některý vstup, provádění zastaralý plán odmítne a~požádá o~novou analýzu. Server při provádění znovu kontroluje závislosti úkolů, stav schopností a~způsobilost databázových tabulek; zaškrtávací pole prohlížeče jsou vstupem prezentace, nikoli oprávněním.

Provádění postupuje přes omezené kontrolní body, takže můžete stránku obnovit či zavřít a~později pokračovat ve stejné úloze. Trvalý nárok na změny vlastní pouze jedna běžící nebo pozastavená centrální úloha a~dvě karty prohlížeče nemohou současně posouvat jednu úlohu. Automatická údržba webu ustoupí před první změnou, dokud je tento nárok držen. Veřejné požadavky galerie pouze pro čtení zůstávají dostupné.

Pozastavení a~zrušení jsou kooperativní: projeví se mezi omezenými operacemi a~nevracejí již potvrzenou práci. Již běžící databázový \codeword{ANALYZE TABLE} nebo \codeword{OPTIMIZE TABLE} se musí dokončit, než může pokračovat zrušení. Centrum zpracovává nejvýše jednu aktuálně způsobilou tabulku na krok prohlížeče, přeskakuje chráněné, neznámé, nadměrné nebo nově nezpůsobilé tabulky a~zaznamenává nejistý atomický výsledek místo slepého opakování tabulky po ztracené odpovědi.

Závěrečná zpráva odděluje odstraněné řádky a~soubory, analyzované či optimalizované tabulky, přeskočení, upozornění a~chyby. Fyzické úložiště před a~po hlásí pouze tehdy, když závěr podporují inventarizační podklady. Obecný postup nikdy nepoužívá migrace, neopravuje schéma, neinstaluje aktualizace, neodstraňuje archivy protokolů administrace, nevyprazdňuje všechny položky Koše ani neobnovuje média; pro tyto záměrné operace použijte specializované stránky.

Úklid stahování a~mezipaměti postupuje přes mezipaměť manifestů, starší mezipaměť podkladů stahování a~mezipaměť generovaných ZIP v~samostatných omezených krocích. Tyto mezipaměti jsou nezávislé optimalizace: není-li některá na aktuálním hostingu dostupná, centrum zachová již dokončený úklid, zaznamená bezpečné upozornění a~pokračuje ostatními oblastmi a~pozdějšími úkoly. Upozornění může určit operaci a~stabilní kategorii chyby, ale neodhaluje nezpracovanou zprávu výjimky, SQL, přihlašovací údaj, soukromou cestu ani zásobník volání. Po aktualizaci měnící tento kontrakt provádění je starší připravený plán záměrně zastaralý; před pokračováním spusťte znovu Analýzu.

\subsection{Údržba náhledů}

Nástroje údržby kontrolují generované varianty, určují chybějící či zastaralé odvozeniny a~obnovují vybrané velikosti. Obnova s~pomocí prohlížeče může podle nastavení připravovat zdrojové bloky. Příprava na pozadí řeší malé mezery veřejného vykreslování omezenými požadavky, zatímco výslovná údržba administrace zůstává hlavní cestou rozsáhlých oprav.

\subsection{Ochrana požadavků SEO a~hygiena veřejných dotazů}
\index{o00124@ochrana požadavků SEO}
\index{s00218@spam parametrů dotazu}
\index{k00061@kanonická URL}

Oblast údržby obsahuje volitelnou ochranu veřejných požadavků SEO, ve výchozím stavu povolenou. Pro anonymní veřejné požadavky GET udržuje výslovný seznam povolených parametrů dotazu pro každou trasu. Známé analytické parametry, například pole UTM a~běžné identifikátory kliknutí na reklamu, ignoruje, protože nemění vykreslený obsah. Neočekávané parametry se před běžným veřejným vykreslovačem odmítají odpovědí 404, \codeword{noindex,nofollow} a~zákazem mezipaměti, aby spam parametrů generovaný roboty nenásobil indexovatelné URL galerií.

Trasy administrace a~strojové koncové body, například nastavení, WebDAV, automatizace nahrávání, přijímače migrací galerií a~koncový bod cronu údržby, jsou mimo tuto ochranu veřejných dotazů. Přihlášení správci mají také výjimku, aby pravidla veřejných robotů neblokovala diagnostické a~správcovské postupy. Protokolování odmítnutých dotazů lze povolit nezávisle a~má vzorkovací limit na den UTC, aby záplava robotů nezaplnila protokol administrace. Stejná služba dodává náhradní kanonické URL, pokud veřejná stránka ještě nevytvořila vlastní kanonickou značku.

\subsection{Údržba databáze}

Průzkum databáze inventarizuje objekty schématu, odkazy, kategorie vlastnictví, pravidla uchování a~jistotu úklidu. Logický úklid používá povolená pravidla s~deterministickým výběrem, omezenými dávkami a~auditními záznamy. Oprava schématu poskytuje plán nanečisto a~idempotentní migrační postup. \codeword{ANALYZE TABLE} může obnovit statistiky optimalizátoru vybraných tabulek. Fyzický \codeword{OPTIMIZE TABLE} je dostupný pouze jako náhled následovaný samostatně potvrzenou operací vybraných tabulek; nikdy není automatickým krokem úklidu.\footnote{Hodnoty alokace databázového stroje, například odhady volných dat, popisují stav stroje a~nezaručují přesný počet bajtů souborového systému uvolněných po optimalizaci.}

\subsection{Statistiky úložiště a~Úplná zpráva}
\index{s00225@statistiky úložiště}
\index{u00243@úplná zpráva}
\index{z00282@zpráva galerie}

Podrobné \ui{Statistiky úložiště} se počítají pouze na žádost správce, takže běžný přehled při každé návštěvě neprovádí úplné skenování souborového systému odvozenin. Zobrazení Soubory odděluje indexované zdrojové fotografie od generovaných náhledů JPEG/WebP a~zobrazovacích předloh DNG pro prohlížeč, ukazuje rozložení typů a~velikostí zdrojových souborů, největší galerie a~zdroje, neznámé velikosti a~upozornění skenování. Snímek v~mezipaměti se při změně dat galerií označí jako zastaralý. Obnovení snímku skenuje očekávané generované soubory v~malých dávkách řízených prohlížečem. Stejná oblast odkazuje na podrobnosti úložiště databáze a~její údržbu.

Karta Databáze čte statistiky úložiště z~informačního schématu odděleně od skenování souborového systému. Pokud možno rozlišuje tabulky vlastněné Gallery od celé vybrané databáze a~ukazuje největší tabulky místo vynucení úplného průzkumu tabulek v~běžném požadavku přehledu. Výslovná akce \ui{Přepočítat statistiky databáze} spouští obnovu statistik databázového stroje pro dostupné tabulky a~hlásí částečná selhání v~protokolu administrace místo předstírání úplného přepočtu.

\ui{Údržba > Úplná zpráva} vytváří jedinou samostatnou zprávu HTML zahrnující úložiště, strukturu a~velikost databáze, strukturu galerií, statistiky obrázků a~EXIF, přehled GPS, telemetrii, provozní protokoly, stav funkcí a~nastavení a~diagnostiku běhového prostředí. Práce s~obrázky a~generovanými médii probíhá v~omezených dávkách řízených prohlížečem, aby zpráva zůstala praktická na sdíleném hostingu. Dočasné časové limity, omezení četnosti a~serverová selhání se opakují s~krátkými rostoucími prodlevami, zatímco souhrny EXIF a~GPS s~mnoha jedinečnými hodnotami zůstávají omezené kvůli úložišti relace a~paměti velkých knihoven. Vybrané období telemetrie ovlivňuje pouze oddíl telemetrie; statistiky obrázků vždy pokrývají celou knihovnu. Souhrny míst GPS používají místní přibližné shlukování a~offline referenční data; pravděpodobné snímky simulátorů či her se vylučují z~četnosti skutečných míst, pokud klasifikaci podporují dostupné podklady.

Dokončená zpráva se vrací prohlížeči ke stažení a~neukládá se jako soubor na serveru. Pro uchování kopie PDF otevřete generované HTML a~použijte funkci prohlížeče Tisk nebo Uložit jako PDF; zpráva obsahuje tiskové styly. S~každou exportovanou zprávou zacházejte jako se správcovskými daty, protože může obsahovat podrobné provozní statistiky a~statistiky sbírky.

\subsection{Telemetrie a~ovládání soukromí}
\index{t00233@telemetrie}
\index{s00215@soukromí!t00233@telemetrie}

Volitelná oblast \ui{Telemetrie} zaznamenává provozní měření s~řízeným soukromím a~představuje je jako zprávy administrace. Dostupné prvky zahrnují hlavní přepínač telemetrie, měření veřejného využití, výkon, události mezipaměti a~databáze, respektování Do Not Track, vyloučení správců, maximální započítávanou dobu prohlížení fotografie a~uchování surových, hodinových a~denních dat. Zprávy zahrnují například anonymní relace, zobrazení stránek, otevření fotografií a~omezený čas prohlížení, chyby klientů, naměřené bajty médií, zastoupení prohlížečů a~aktivitu mezipaměti a~databáze, pokud jsou tyto kategorie povoleny.

Subsystém telemetrie je navržen tak, aby neukládal nezpracované adresy IP, řetězce user-agent prohlížečů, URL odkazujících stránek, jména, e-mailové adresy, identifikátory účtů, těla požadavků ani přesné polohy návštěvníků. Exporty obsahují souhrnnou anonymní telemetrii a~hashe relací místo nezpracovaných identifikátorů návštěvníků. Údržba shrnuje a~maže telemetrii podle nastavených pravidel uchování a~obrazovka administrace může exportovat samostatnou zprávu HTML pro offline kontrolu. Před povolením sběru na veřejném webu zkontrolujte ovládání soukromí.

\subsection{Diagnostika běhového prostředí, integrita a~měření galerie}
\index{d00029@diagnostika běhového prostředí}
\index{k00067@kontrola integrity}
\index{m00080@měření galerie}

\ui{Diagnostika běhového prostředí} poskytuje kopírovatelný pohled na prostředí a~podklady pro podporu pro PHP, třídy připravenosti databáze, limity hostingu, GD, Imagick/ImageMagick a~příslušnou podporu formátů médií. Zpřístupňuje také dříve popsaná pravidla převodu DNG. Panely připravenosti oddělují závislosti prezentace a~zpráv, autentizace a~zabezpečení a~závislosti citlivé na změny, aby správce rozlišil volitelnou chybějící schopnost od podmínky, která musí blokovat zápis.

Stejný omezený model hlásí připravenost revizí úprav galerií a~nevyřešené přesuny obrázků. Chybějící nebo neznámé požadované schéma je stavem vyžadujícím akci a~zastavuje příslušnou změnu před změnou cílových souborů či řádků. Souhrny čekajících přesunů obsahují pouze omezené identity operací a~galerií a~bezpečné kategorie stavů; pro výslovnou operaci použijte dokumentovaný postup obnovy z~příkazového řádku místo očekávání, že diagnostická stránka přesune soubory.

Samostatný přepínač \ui{Vývojová diagnostika} je pouze pro správce. Při povolení pro přihlášeného správce přidává měřicí nástroje prohlížeče pro přednačítání, mezipaměť, paměť, síť a~časování snímků ve veřejném zobrazení, lightboxu a~celé obrazovce. Nezapíná diagnostiku běžným anonymním návštěvníkům. Při běžné správě jej ponechte vypnutý, pokud nezkoumáte vykreslování nebo výkon; měření galerie také závisí na tomto kontextu vývojové diagnostiky.

Postup \ui{Integrita} v~administraci kontroluje soubory jádra spravované integritou proti instalovanému manifestu a~pomáhá určit neúplné či neočekávané změny nasazení. Jde o~ověřovací nástroj za běhu; souběžné zálohy jsou stále nutné, protože manifest integrity neobsahuje uživatelské fotografie ani nenahrazuje média obnovy.

Je-li dostupná vývojová diagnostika, přihlášený správce může pro aktuální galerii zahájit \ui{Měření načítání veřejné galerie}. Prohlížeč několikrát načte galerii ve skrytém rámci stejného původu s~anonymním náhledem, spojí čítače vykreslování PHP s~časováním navigace prohlížeče a~poté nabídne protokol měření jako JSON. Při porovnávání dvou změn použijte stejnou galerii, režim prezentace a~podmínky sítě a~mezipaměti; měření je nejužitečnější pro porovnání před a~po, nikoli jako univerzální skóre.

\subsection{Plánovaná údržba webu}
\index{p00135@plánovaná údržba}
\index{c00019@cron!u00241@údržba webu}

Běžná údržba může probíhat v~omezených částech místo jediného dlouhého požadavku. Služba údržby řetězí bezpečné webové požadavky, je-li dostupný provoz, a~lze ji také volat z~cronu hostingu příkazem CLI uvedeným v~administraci. Může kontrolovat stav náhledů, obnovovat metadata, zpracovávat úklid, shrnovat telemetrii a~archivovat způsobilé dny protokolu administrace podle nastavených pravidel. Skrytý token webového cronu lze obměnit; obměna zneplatňuje dříve zveřejněné URL webového cronu. Údržba je ve výchozím stavu nedestruktivní a~nenahrazuje výslovné potvrzení optimalizace databáze ani jiných závažných akcí.
\subsection{Historie a~archivy protokolů administrace}
\index{p00163@protokoly administrace}
\index{a00008@archiv protokolů}

Každodenní správa vytváří užitečnou historii: nahrávání, změny galerií, bezpečnostní rozhodnutí, výsledky údržby, upozornění a~další provozní události mohou zanechat záznam v~protokolu administrace. Historie je cenná, když chcete pochopit, co se stalo, kdy se to stalo a~který účet nebo požadavek byl zapojen. Na aktivní instalaci je také běžné, že se tabulka značně zvětší. Verze~\version{} proto zachází s~péčí o~protokoly jako se samostatným výslovným úkolem údržby místo umožnění neočekávaného odstranění auditní historie obecným úklidem databáze.

Obrazovka protokolu administrace může zobrazovat jednotlivé položky nebo sbalit opakované události do shrnutí. Skupina může představovat opakovaná upozornění náhledů, výsledky nahrávání nebo zprávy údržby; její otevření ukáže jednotlivé záznamy. Filtry a~ovládání stavu postupů pomáhají zúžit provozní kontrolu, zatímco omezené exporty uchovávají vybranou historii bez načtení celé tabulky jediným požadavkem. Velké seznamy, rozbalení jednotlivých výskytů a~exporty se zpracovávají v~omezených částech, takže požadavek úplné historie nevyžaduje, aby prohlížeč nebo PHP držely celou tabulku v~paměti najednou.

Vlastník může zvolit, jak dlouho mají nedávné protokoly zůstávat v~aktivní databázi. Volba \emph{Navždy} vypíná automatickou archivaci. Volba konečné aktivní doby starší záznamy okamžitě nezahazuje. Postup údržby nejprve vybere úplné kalendářní dny mimo aktivní období uchování, zapíše chráněný archiv obsahující strojově čitelný JSON i~čitelnou statickou zprávu HTML, ověří očekávaný počet záznamů v~archivu a~teprve poté odstraní stejné řádky z~aktivní tabulky. Archiv je uložen pod \filepath{data/admin-log-archives/}, což jsou data aplikace, nikoli veřejná oblast galerií, a~má dodatečné ochranné pravidlo webového serveru.

V~údržbě archivů lze pokračovat. Před povýšením dokončeného archivu zapisuje dočasné soubory, zaznamenává malý stavový soubor a~používá zámek, aby dva požadavky údržby nezpracovávaly stejnou historii současně. Pokud práci zastaví limity hostingu, přerušené spojení nebo jiný problém, obrazovka administrace může ukázat neúplný stav a~další pokus údržby jej může ověřit, pokračovat nebo obnovit. Archiv se nikdy nepovažuje za důvod odstranění aktivních záznamů jen proto, že existuje soubor s~podobným názvem; kontroly počtu řádků a~metadat archivu jsou součástí bezpečnostní hranice.

Správci by nadále měli zahrnovat databázi a~adresář \filepath{data/admin-log-archives/} do souběžných záloh. Aktivní databáze obsahuje nedávnou provozní historii, zatímco archivní adresář obsahuje starší historii vybranou výslovnými pravidly uchování. Zachování obou poskytuje vlastníkovi souvislý záznam a~usnadňuje zkoumání obnovené instalace po aktualizaci, migraci nebo neočekávané události.
\subsection{Úplný přehled současných funkcí}
\label{sec:complete-feature-reference}
\index{p00171@přehled funkcí}
\index{s00217@soupis funkcí}

Tento přehled je kontrolním seznamem úplnosti příručky pro verzi~\version{}. Uvádí současné schopnosti produktu zpřístupněné veřejnými trasami, navigací administrace, editory galerií a~fotografií, nástroji údržby a~volitelnými integracemi. Interní pomocné funkce se neuvádějí jako samostatné funkce, pokud pouze implementují jeden z~níže uvedených řádků.

\begin{longtable}{>{\RaggedRight\arraybackslash}>{\hspace{0pt}}p{0.23\textwidth} >{\RaggedRight\arraybackslash}>{\hspace{0pt}}p{0.66\textwidth}}
\toprule
\textbf{Funkce} & \textbf{Implementované chování}\\
\midrule
\endfirsthead
\toprule
\textbf{Funkce} & \textbf{Implementované chování}\\
\midrule
\endhead
Galerie založené na souborovém systému & Vnořené fyzické složky zrcadlené databázovými záznamy, prakticky neomezená hierarchie, vyhledání a~import, ruční vytváření, fyzické přesuny, čisté veřejné cesty, doprovodná metadata a~stabilní pořadí sourozenců a~fotografií.\\
Metadata galerie & Název, popis, udržované překlady, ruční datum a~rozsah, návrhy dat EXIF, identifikátor v~URL, název složky, rodič, pořadí, štítky, titulní fotografie, nahraný náhled galerie, vlastní značka, pozadí a~přepsání prezentace.\\
Přístup ke galerii & Veřejná, nezveřejněná, soukromá, uzamčení heslem, odkazy pro sdílení s~vypršením i~bez něj, vyhodnocování zděděného přístupu, přístup správce a~hranice potvrzení NSFW / 18+.\\
Hromadná správa galerií & Odstranění zohledňující větve, změny viditelnosti, přepsání map EXIF/GPS, hlasování, zobrazení názvů souborů a~hry s~fotografiemi, obnova cest, operace náhledů a~skenování a~postupy změny pořadí sourozenců a~řazení podle data.\\
Chytré galerie & Uložené vnořené booleovské dotazy nad fyzickými galeriemi, štítky, daty, EXIF/GPS, textem, metadaty AI, stavem duplicit, vlastnostmi souborů a~soukromým redakčním hodnocením; veřejné, neuvedené, kořenové a~připojené umístění; pořadí nad a~pod obsahem pro každého rodiče; ochrana před cykly; nezávislá přepsání prezentace; úplné procházení stránkovaného lightboxu; omezené stahování ZIP.\\
Editor fotografie & Název, popis a~překlady, viditelnost, NSFW, pořadí, soukromé redakční hodnocení, štítky, meze náhledů jednotlivých fotografií, kontrola EXIF/GPS, interních metadat AI a~nástroje návrhů OpenAI.\\
Hromadná správa fotografií & Přesun do existující či nové galerie, odstranění, přiřazení titulního obrázku, viditelnost, NSFW, vytváření náhledů a~pořadí přetažením či podle názvu souboru.\\
Štítky & Sdílené štítky galerií a~fotografií, normalizované opakovaně použitelné názvy a~identifikátory, veřejné popisy, počty využití, vážené návrhy editoru, přejmenování, úpravy a~odstranění v~administraci, veřejné cílové stránky štítků, vlastní nastavení mřížky a~karet štítků, řazení, rozbalování a~posouvání štítků úvodního panelu.\\
Veřejná domovská stránka a~stránky galerií & Hierarchické karty, drobečková navigace, titulní obrázky a~koláže, stránkování, přímé mřížky fotografií, počty galerií, data, popisy, vlastní značka a~pozadí, oblíbené zkratky záhlaví, kontextové prvky úprav, odstraňování a~pořadí pro přihlášené správce, plovoucí návrat nahoru u~dlouhých seznamů, přizpůsobivé veřejné cesty a~náhradní parametry dotazu při vypnutém přepisování.\\
Veřejné vyhledávání & Volitelné vyhledávání webu a~větve v~názvech, popisech, názvech souborů, metadatech štítků, textu udržovaných jazyků a~dostupném interním prohledávatelném textu AI; minimální délka dotazu a~omezený počet výsledků; pravidla přístupu a~uvádění použitá před vrácením výsledků.\\
Robots, mapa webu a~SEO & \codeword{robots.txt} zohledňující přístup, \codeword{sitemap.xml} zohledňující obrázky, kanonická, sociální a~strukturovaná metadata a~výše popsaná ochrana parametrů dotazu specifická pro trasu.\\
Veřejná autorizace médií & Vyhrazené trasy originálů, náhledů, titulních obrázků, značky, pozadí a~ikony webu znovu vyhodnocují viditelnost galerií a~obrázků a~chráněný přístup místo odhalení libovolných cest souborového systému.\\
Generování náhledů & Responzivní velikostní varianty, pravidla kompatibility JPEG/WebP, responzivní a~progresivní vykreslovače vybrané galerie, odložené načítání a~chování poblíž viditelné oblasti, meze velikostí galerií a~fotografií, zobrazovací předlohy DNG, sladění metadat, postupy obnovy v~administraci a~prohlížeči a~podepsaná veřejná příprava náhradních náhledů s~kontrolou přístupu a~malým rozpočtem.\\
Lightbox & Povolená asynchronní metadata, navigace klávesnicí a~dotykem, prezentace a~celá obrazovka, režimy jednoho obrázku, pásu fotografií a~3D karuselu, překryv EXIF a~mapy, stav hlasování, přiblížení 100--400~\% s~posunem, sevřením prstů a~ukotvením kurzoru, okamžité přepnutí na kvalitu originálu při přiblížení a~skoky Shift+šipka o~10 fotografií.\\
EXIF a~GPS & Získávání údajů fotoaparátu, objektivu, expozice, data a~GPS, globální výchozí hodnota s~děděním či vynuceným zapnutím nebo vypnutím galerie, značky fotografií, mapová data galerie, návrhy dat, podklady duplicit a~veřejné potlačení citlivé na soukromí.\\
Hlasování & Vratný kladný hlas pro obrázek se synchronizací skóre, anonymním hashem návštěvníka nebo přihlášeným vlastnictvím, omezením četnosti nových hlasů a~zapnutím či vypnutím pro galerii.\\
Hra s~fotografiemi & Volitelné porovnání párů na úrovni větve se způsobilými veřejnými obrázky, trvalým stavem párů a~vítězů, žebříčkem a~souhrnem výsledků a~bezpečným odmítnutím při neznámém schématu úložiště.\\
Správce fotografií & Správcovský překryv veřejného zobrazení pro přihlášené s~výběrem fotografií a~fyzických podgalerií, obnovitelným odstraněním, vytvořením fyzické galerie pod zvoleným rodičem s~kopírováním podstromů a~vrácením při selhání, akcemi přesunu, kopírování, přetažení a~sdílení pouze fotografií, nativním Web Share více souborů, je-li podporováno, a~náhradním či výslovným stažením ZIP vybraných fotografií.\\
Stahování & Povolené ZIP galerií, chytrých galerií a~vybraných fotografií, ZIP všech galerií pro správce, normalizované cesty archivu, podpisy obsahu mezipaměti, omezený počet a~bajty zdrojů chytré galerie a~bezpečný úklid mezipaměti.\\
Nahrávání v~prohlížeči & Postupy nahrávání do existující či nové galerie, výslovný výběr zpracování prohlížečem nebo serverem, úplné ověření připravených náhledů, omezené pracovní procesy a~dávky se samostatnými nadměrnými balíčky pod tvrdým limitem PHP, pokyny pravidel formátů, místní import ZIP a~automatické přejmenování po skenování.\\
Vyhledání v~souborovém systému & Omezené skenování složek podložené relací; přeskočení symbolických odkazů, skrytých a~generovaných adresářů; import hierarchie s~obrázky na místě; možnost přesunu nespravovaných fotografií do existující galerie nebo odstranění bezpečných nespravovaných kandidátních stromů; samostatné hlášení doprovodných souborů obsahujících jen metadata.\\
Doprovodná metadata souborového systému & Volitelný \filepath{gallery.json} každé galerie zapisovaný z~upravitelného stavu galerie a~čtený při vyhledání či opravě rodičů pro podporovaná metadata; užitečný pro přenositelnost souborového systému, ale není náhradou zálohy databáze.\\
API nahrávání & Klíče API omezené na galerii a~viditelné pouze jednou, centrální správce API, zrušení klíčů, společný postup příjmu a~ověřený koncový bod automatizace.\\
Doprovodná aplikace Windows & Nahrávání ze sledované složky, ruční hromadné nahrávání, volitelné místní generování responzivních náhledů, souběžná a~navazující práce, opakované pokusy a~stavové soubory, oznamovací oblast, volitelné souřadnice kamery SimConnect, odstranění zdroje po potvrzeném úspěchu a~volitelný místní pracovní proces metadat AI.\\
Mobilní WebDAV & Přihlašovací údaje s~jednorázově zobrazeným heslem omezené na galerii, autentizace Basic, metody vyhledání a~zápisu WebDAV požadované mobilními klienty přenosu, příjem JPG/PNG/GIF/WebP, sledování posledního použití a~zrušení.\\
Přejmenování médií & Plán fyzických názvů souborů nanečisto, výchozí \codeword{\{gallery_path\}_\{seq4\}} a~dokumentované zástupné značky, řešení kolizí a~bezpečnosti, kontext pořadí galerie, sladěné aktualizace databáze, veřejných cest, odvozenin a~mezipaměti, omezené použití na celém webu a~integrace automatického přejmenování po nahrání.\\
Organizátor metadat & Pouze náhledové seskupování dat EXIF s~filtry nejstaršího a~nejnovějšího data, podřízené složky ISO \codeword{YYYY-MM-DD}, opětovné použití odpovídajících přímých potomků podle data, potvrzené fyzické přesuny s~odvozeninami a~omezené dávky provádění; seskupování podle polohy zůstává plánované, nikoli aktivní.\\
Detektor duplicitních fotografií & Přesné páry SHA-256 a~možné páry podpořené metadaty, kontextové odkazy, postup odstranění, záznamy zkontrolovaných párů a~přesných galerií pro každého správce a~trvalý stav kontroly.\\
Textová pomoc OpenAI & Šifrovaný klíč API a~nastavení modelu pro účet, volitelný výslovný souhlas s~náhledy, návrhy popisů galerií a~fotografií, úprava textu, shrnutí souvislostí, popis viditelného obsahu z~malých generovaných náhledů, návrhy překladů a~vložení pouze s~ruční kontrolou.\\
Interní metadata AI & Samostatná prohledávatelná strojová metadata a~tabulky fronty, ověřené převzetí úloh a~výsledky pracovního procesu Windows, kontrola v~editoru fotografie pouze pro čtení, podklady chytrých galerií a~vyhledávání, sledování generací modelu a~vynucené opětovné zpracování větve.\\
Letové mapy a~SimBrief & Ruční text trasy, výslovná syntaxe souřadnic, načtení nejnovějšího OFP SimBrief, upravitelný generovaný Markdown, uložená data OFP a~souřadnice a~veřejné vykreslování pouze z~uložených bodů.\\
Navigační data & Import místní databáze OurAirports, přibalený offline CSV, body poskytovatele v~mezipaměti, volitelná integrace Navigraph OAuth, balíčků a~AIRAC, vyhodnocování přednostně offline a~tester vyhledávání v~administraci.\\
Migrace galerie & Přenos API odesláním ze zdroje nebo stažením cílem při stejné verzi, verzované manifesty, kontrolní součty metadat a~podkladů, omezené požadavky jednoho podkladu, trvalý stav cíle, opětovné připojení a~pokračování a~samostatné dokončení.\\
Vzhled a~rozvržení & Přímé barvy palety, režim písma, zaoblené rohy, šířka stránky, mřížky galerií a~karet, stránkování, rozvržení karet, odznak počtu, veřejný vykreslovač náhledů, výchozí lightbox, prezentace značky GPS, tři oblíbené zkratky záhlaví a~živý náhled. Současné rozhraní Vzhledu nemá samostatný přepínač světlého a~tmavého režimu.\\
Značka a~identita v~prohlížeči & Banner záhlaví webu, rozměry a~roztažení oddělovače, optimalizované pozadí a~průhlednost, přepsání galerií, veřejné podklady titulních obrázků, čtvercově oříznutá ikona webu s~generovanými variantami 32/48/180~pixelů a~předvolby a~editor vlastního CSS.\\
Lokalizace & Kanonická anglická náhrada a~udržované české, německé a~švédské balíčky rozhraní dostupné ve výběru, preference prohlížeče jednotlivých návštěvníků, návrhy výběru vlajek a~rodných názvů, společné katalogy prohlížeče a~serveru, diagnostika, editor, import a~export balíčků a~obsah galerií a~fotografií v~udržovaných jazycích.\\
Přepínače funkcí & Globální přepínače zohledňující trasy pro veřejné vyhledávání, režimy lightboxu, Správce fotografií, hlasování o~obrázcích, hru s~fotografiemi, stahování, účty návštěvníků a~soukromé sbírky, mapy galerií a~letů, navigační data, SimBrief, OpenAI, místní metadata AI, API nahrávání, mobilní WebDAV, migrace galerií, Přejmenování médií a~telemetrii.\\
Centrální Nastavení & Řízený postupný Průvodce nastavením s~konečným souhlasem, živým náhledem a~bezpečnými adaptéry vlastníků; prohledávatelný kategorizovaný registr, klávesnicová navigace ve stylu Spotlight, kanonické předávání ukládání schválených globálních hodnot, přímé odkazy specialistů, vlastnictví zohledňující dědění, popisky výchozích hodnot a~zdrojů a~skrytí tajných údajů.\\
Pracovní prostředí galerií v~administraci & Filtr viditelnosti s~rušením výběru skrytých řádků, trvalý stav sbalení větví, Sbalit/Rozbalit vše, změna pořadí hierarchie, veřejné nástroje pořadí viditelné stránky, editory postranního panelu a~kontextové prvky úprav, odstranění a~viditelnosti na veřejné stránce pro přihlášené správce.\\
Autentizace správce & Přihlášení uživatelským jménem nebo e-mailem pro obnovu, ochrana CSRF a~relace, hashované subjekty omezení četnosti, volitelný trvalý token přihlášení, obnova hesla přes PHP mail/SMTP, změny účtu a~profilu chráněné kontrolami aktuálního hesla a~volitelné přihlášení Google, které musí být nejprve propojeno z~ověřené relace s~heslem a~k~odpojení vyžaduje aktuální heslo.\\
Účty návštěvníků, životní cyklus, oblíbené fotografie, soukromé sbírky a~sdílení neuvedené v~seznamech & Přímé vytváření a~odstraňování návštěvníků správcem a~registrace pouze na pozvání, správcovské pozastavení, obnovení a~odhlášení všude s~centrálním rušením oprávnění návštěvníka, vynucené nahrazení dočasných hesel poskytnutých správcem před běžným oprávněním návštěvníka, ověření e-mailu a~aktivace bezpečné vůči skenerům, samostatné přihlášení a~odhlášení návštěvníka, vyhrazený rotující zapamatovaný přihlašovací údaj, obecná obnova hesla bezpečná vůči skenerům, nedávné opětovné ověření heslem pro citlivé akce účtu, změna hesla a~e-mailu, vlastní odstranění návštěvníka, soukromé odpovědi účtu, oblíbených fotografií, sbírek a~životního cyklu bez ukládání do mezipaměti, prvky oblíbených fotografií na povolených kartách a~obrázcích lightboxu, seřazené sbírky návštěvníka s~aktuálním oprávněním zdrojů, jeden zrušitelný odkaz sbírky pouze ke čtení na 30~dnů, neuvedený v~seznamech, s~jednorázově zobrazeným tajným údajem a~čistým přesměrováním, průnik ACL zdroje v~kontextu příjemce bez výjimky správce a~přísné oddělení od oprávnění galerií; bez otevřené registrace, veřejného vyhledávání sbírek a~profilů, spolupráce příjemců, komentářů či nahrávání.\\
Připravenost databáze a~Stav systému & Třístavový průzkum schématu rozlišující Dostupné, Chybějící, Neznámé a~tam, kde je to vhodné, Vypnuté; samostatné panely schopností zabezpečení a~autentizace, citlivosti na změny a~volitelné prezentace a~zpráv; zápisy se při neznámém požadovaném schématu bezpečně odmítají.\\
Údržba databáze & Výslovný soupis schématu, migrací a~odkazů kódu, informace o~úložišti, kandidáti úklidu nanečisto, omezený logický úklid, vyhrazená připravenost starších oprav, vybraný ANALYZE, náhled optimalizace vybraných tabulek a~samostatně potvrzený OPTIMIZE TABLE.\\
Statistiky úložiště & Skenování zdrojů a~generovaných médií na vyžádání, rozložení typů a~velikostí souborů, největší galerie a~soubory, evidence předloh DNG a~náhledů, práce se zastaralými snímky, samostatný pohled využití databáze a~tabulek, výslovný přepočet statistik databáze, odkazy údržby databáze a~omezené dávky obnovy.\\
Úplná zpráva & Samostatný export HTML pokrývající úložiště, databázi, statistiky galerií, obrázků, EXIF a~GPS, stav funkcí a~nastavení, telemetrii, protokoly a~diagnostiku běhového prostředí s~omezeným zpracováním obrázků a~generovaných médií, opakováním dočasně neúspěšných požadavků a~omezenými souhrny mnoha jedinečných hodnot.\\
Telemetrie & Volitelná anonymní měření využití, výkonu, mezipaměti a~databáze s~řízeným soukromím, respektování DNT, vyloučení správců, omezený čas prohlížení, uchování surových, hodinových a~denních dat, shrnování a~mazání, přehled administrace a~export HTML.\\
Protokoly administrace & Strukturované události, filtry a~stav, seskupování opakovaných událostí s~rozbalením členů, omezené exporty a~export ZIP, nastavitelná aktivní doba uchování, chráněné denní archivy JSON+HTML, stav pokračování a~obnovy a~nástroje stahování a~prohlížení archivů.\\
Centrum údržby & Trvalý postup Analyzovat--Zkontrolovat--Provést--Ověřit, plán úkolů autorizovaný serverem, omezené kontrolní body, pozastavení, pokračování a~zrušení, centrální nárok na změny, fyzické kroky databáze po jedné tabulce a~ověřená zpráva před a~po.\\
Plánovaná údržba webu & Denní okno UTC, volitelné pokračování spouštěné požadavky, omezený rozpočet dávky obrázků a~času, podpora CLI a~webového cronu, akce jednoho úseku, obnovení přerušeného stavu, obměnitelný skrytý token webového cronu, úlohy náhledů, metadat, úklidu, telemetrie a~archivace protokolů a~ustoupení, když Centrum údržby vlastní centrální nárok na změny.\\
Diagnostika běhového prostředí & Připravenost PHP, hostingu a~databáze, podpora GD/Imagick/ImageMagick a~formátů, pravidla DNG, kopírovatelné shrnutí podpory, vývojová diagnostika pouze správce pro přednačítání, mezipaměť, paměť, síť a~časování snímků prohlížeče a~vstupy měření pouze při vývoji.\\
Integrita & Ověřování hashů instalovaného manifestu jádra pro spravované soubory aplikace a~generování i~ověření manifestu při vydání; oddělené od zálohy uživatelských dat.\\
Měření galerie & Ověřené měření v~prohlížeči pouze při vývoji používající anonymní náhledová načtení, skryté rámce stejného původu, čítače vykreslování serveru a~časování navigace a~poté exportující JSON pro porovnání před a~po.\\
Aktualizátor & Vyhledání stabilních a~beta vydání, poznámky k~vydání v~mezipaměti, trvalé úlohy s~pokračováním, obnovení stahování Range, ověření archivu a~manifestu, příprava, snímek pro vrácení, aktivace, migrace, úklid, zámky a~obnova zastaralých zámků, opakování či zrušení před aktivací, vrácení souborů po aktivaci, obnova stabilního vydání, čistá přeinstalace či instalace beta, pokračování CLI a~nouzový \filepath{reset.php}.\\
Pomocné nástroje nasazení & Skripty nasazení Windows/Linux, metadata vydání, generování manifestu jádra a~ověření balíčků s~nasazením navrženým tak, aby místně instalovaný spustitelný soubor PHP nebyl bezpodmínečným běhovým požadavkem samotné galerie.\\
\bottomrule
\end{longtable}
\subsection{Aktualizace aplikace}

Kompaktní přehled vydání sdružuje nainstalovanou a~dostupnou verzi, kanál, stav a~hlavní akce; podrobnosti repozitáře a~GitHub API zůstávají ve sbalené diagnostické sekci. Po dokončení úlohy prohlížeč obnoví přehled z~místních metadat přímo na místě, ponechá panel otevřený a~URL beze změny a~skryje již neaktuální akci stabilní aktualizace. Obnova neprovádí další kontrolu GitHubu. Pokud selže, tlačítkem obnovení stavu zopakujte místní načtení bez opakování instalace.

Aktualizátor kontroluje nastavené kanály, ověřuje balíčky, připravuje obsah, ověřuje požadované běhové soubory a~velikosti, uchovává přepsané soubory v~úložišti obnovy a~aktivuje aktualizaci. Obsluha omezení četnosti a~stav opakovaných pokusů brání opakovaným vzdáleným požadavkům během čekacích období poskytovatele. Karta úlohy s~možností pokračování se objevuje ve \ui{Stav} i~\ui{Pokročilé nástroje}, takže instalace beta, obnova nebo přeinstalace nadále ukazuje průběh tam, kde začala. Stránku lze zavřít a~znovu otevřít bez zahození dokončených kontrolních bodů.

Oblast aktualizací v~administraci také ukazuje informace o~vydání v~mezipaměti a~verzované poznámky k~vydání dostupného kanálu, aby si vlastník mohl vydání prohlédnout před zahájením úlohy. Stabilní a~beta kanály zůstávají samostatnými volbami; označení Beta náleží volbě vydání v~rozhraní administrace a~nemění terminologii trvalých úloh ani manifestu aktualizátoru.

Dočasná selhání stahování či hostingu lze opakovat. Balíček, jehož soubory neodpovídají vlastnímu manifestu integrity, je jiný případ: opětovné stažení stejného publikovaného sestavení nemůže změnit tyto bajty. PHP Gallery proto žádá správce o~zrušení připravené úlohy a~výběr novějšího vydání nebo beta kódu místo nabídky nekonečného opakování.

Vyhledání stabilního vydání a~práce s~balíčkem jsou samostatné omezené operace. Požadavek vyhledání může vytvořit trvalou úlohu na pozadí, zatímco pozdější bezpečný požadavek nebo nástroj CLI posouvá krátký úsek přes stahování, ověření archivu, rozbalení, ověření manifestu, přípravu, zálohu pro vrácení, aktivaci, migrace a~úklid. Požadavky Range mohou obnovit přerušené stahování, zámky pracovních procesů brání souběžnému postupu a~zastaralý zámek lze obnovit po vypršení jeho vlastníka. Na nečinném sdíleném hostingu naplánujte \codeword{php scripts/application_update.php} z~cronu, například každých pět minut. Každé spuštění provádí omezené vyhledání nebo posouvá existující práci na pozadí; správnost nezávisí na \codeword{set_time_limit()} ani \codeword{ignore_user_abort()}.

Současná vydání také během přípravy a~aktivace aktualizace uvádějí do souladu malou sadu souborů serverových pravidel vlastněných aplikací. To chrání instalace, jejichž starší aktualizátor přeskočil soubor pravidel: ověřená kopie vydání se porovná a~atomicky zveřejní, metadata vrácení se rozšíří před nahrazením a~počáteční postup podložený migrací může opravit již aktivovaný starší strom. Sladění je omezené na známé cesty pravidel a~nemění média galerií ani uživatelskou konfiguraci.

Před aktivací může správce obnovitelnou úlohu zrušit nebo zopakovat. Po aktivaci může akce obnovy v~administraci obnovit spravované soubory aplikace z~uloženého snímku pro vrácení, ale vrácení souborů nevrací databázové migrace. Obnova stabilního vydání, čistá přeinstalace a~instalace beta zůstávají úlohami s~pokračováním a~vlastním ověřením balíčku. Samostatný nouzový postup \filepath{reset.php} je odděleným nástrojem obnovy pro návrat ke stabilní větvi, když nelze použít běžnou trasu administrace; stále vyžaduje stejně pečlivou zálohu a~kontrolu integrity.

Během krátké kritické části aktivace časná běhová brána bez závislostí brání novým běžným požadavkům načítat částečně nahrazenou aplikaci. Tyto požadavky dostávají soukromou odpověď \codeword{503 Service Unavailable} bez ukládání do mezipaměti, dokud není dokončení aktivace trvale zaznamenáno. Ověřený postup obnovy a~stavu aktualizace zůstává dostupný, aby přerušená aktivace mohla pokračovat. Poškozený či neúplný stav brány bezpečně odmítá přístup místo prezentace instalace jako zdravé.

Stejná časná běhová hranice převádí nezachycená a~zachytitelná fatální selhání PHP na omezené odpovědi \codeword{500}, pokud PHP ještě může řídit odpověď. Produkční hosting musí ponechat \codeword{display\_errors} vypnuté; jinak může samotné běhové prostředí PHP vypsat podrobnosti upozornění či fatálních chyb dříve, než obsluha aplikace naformátuje bezpečnou odpověď.

\important{Před aktualizací aplikace proveďte souběžnou zálohu všech souvisejících částí. Ponechte ji dostupnou, dokud neověříte běhové kontroly, migrace, veřejné galerie, chráněný přístup, nahrávání a~operace administrace.}

\subsection{Manifest integrity}
\index{m00077@manifest jádra}

\filepath{app/core-manifest.json} zaznamenává verzi aplikace a~hashe souborů jádra spravovaných integritou. Místní generátor pod \filepath{scripts/generate_manifest.php} počítá konečné hashe po změnách běhového kódu. Zahrnutí dokumentace a~uživatelských dat se řídí zavedenými pravidly generátoru.
\section{Volitelné souvislosti: jak PHP Gallery funguje}
\label{sec:architecture}
\index{a00007@architektura}
\index{r00201@řadiče}
\index{s00210@služby}
\index{p00139@pohledy}

Následující stránky popisují cestu požadavku používanou aplikací. Jsou důležité především pro údržbu a~vývoj; běžná správa galerií je nevyžaduje.

\subsection{Inicializace, směrování a~předání řadiči}
\index{i00055@inicializace}
\index{s00211@směrování}
\index{p00168@předání řadiči}
\index{p00155@požadavek!z00284@životní cyklus}

Běžné požadavky vstupují přes \filepath{public/index.php}; \filepath{index.php} v~kořeni repozitáře je podporovaný kompatibilní vstup pro jiná uspořádání hostingu. Veřejný vstup načítá \filepath{app/early\_runtime.php} před \filepath{app/bootstrap.php} a~obě uspořádání pak dosahují stejného běžného postupu zpracování požadavku.

Inicializace čte místní konfiguraci a~ustavuje kontext požadavku: přístup k~databázi, kde je potřebný, jazyk, stav relace, bezpečnostní předpoklady a~sdílené běhové služby. Směrování poté převádí příchozí URL nebo náhradní podobu s~parametry dotazu na interní trasu a~parametry, například cestu galerie, číslo stránky, velikost média nebo token sdílení. Dispečer zavolá řadič registrovaný pro tuto trasu.

Pořadí je záměrné. Požadavek musí být klasifikován před rozhodnutím kódu, jak jej obsloužit, a~kontroly přístupu musejí proběhnout před sestavením chráněného výstupu. Zdánlivě přímá URL obrázku či náhledu je tedy stále požadavkem aplikace; řadič médií jej před vrácením bajtů autorizuje.

\begin{lstlisting}
browser request
  -> index.php or public/index.php
  -> app/bootstrap.php
  -> cms_run()
  -> cms_route_from_request()
  -> controller function
  -> service functions
  -> HTML, JSON, media, archive, or redirect response
\end{lstlisting}

\subsection{Řadiče}
\index{r00200@řadič}

Řadiče vlastní koordinaci na úrovni HTTP. Přijímají vyřešenou trasu a~kontext požadavku, ověřují operaci, volají potřebné služby a~vybírají typ odpovědi. Veřejná trasa galerie běžně vrací HTML; akce postranního panelu administrace může vrátit JSON nebo fragment HTML; trasa médií vrací povolený soubor; trasa stahování může průběžně odesílat archiv.

Veřejné vykreslování galerie začíná v~\filepath{app/controllers/public_gallery.php} a~související odpovědnosti jsou rozdělené do cílených souborů řadičů. Nahrávání, údržba, aktualizace, stahování, přístup k~médiím a~funkce administrace mají vlastní vstupní body, zatímco společná doménová pravidla zůstávají ve službách.

Řadiče nejsou zkratkou oprávnění. Trasy náhledů, map, lightboxu, vyhledávání, stahování a~přímých médií opakují příslušné rozhodnutí o~přístupu místo předpokladu, že návštěvník již úspěšně otevřel stránku galerie.

\subsection{Služby}
\index{s00210@služby}

Opakovaně použitelná pravidla aplikace jsou v~\filepath{app/services/}. Služby vyhledávají galerie a~obrázky, vyhodnocují přístup, normalizují nastavení, spravují nahrávání a~změny, připravují manifesty náhledů a~médií, provádějí vyhledávání, udržují metadata a~koordinují další doménovou práci. Řadič by měl mít možnost požádat o~některou z~těchto operací bez nové implementace jejích pravidel.

Oddělení je důležité, protože stejná fotografie může být dosažitelná z~několika tras. Pravidla přístupu, výběr náhledů nebo význam odstranění nesmějí záviset na tom, zda požadavek přišel z~karty galerie, výsledku vyhledávání, lightboxu, nástroje administrace nebo přímého odkazu. Zdrojové fotografie zůstávají ve stromu galerií; služby koordinují databázové záznamy a~generované odvozeniny, které je popisují nebo podporují.

\subsection{Pohledy a~moduly prohlížeče}
\index{p00139@pohledy}
\index{m00090@moduly prohlížeče}

Pohledy vykreslují připravená data do HTML. Odpovídají za strukturu a~prezentaci, nikoli za rozhodování, zda má návštěvník oprávnění nebo jak se generuje náhled. Serverový výstup obsahuje významové odkazy, formuláře, názvy, alternativní text a~počáteční kandidáty médií před zahájením vylepšení JavaScriptem.

Moduly prohlížeče přidávají chování lightboxu, návrhy vyhledávání, hlasování, průběh nahrávání, akce přetažení, vylepšení náhledů, diagnostiku a~postranní panel administrace. Kód zůstává bez frameworku. Veřejné stránky tak zachovávají užitečný základ vykreslený serverem a~ověřené postupy mohou přidávat bohatší interakce na místě.

Postranní panel administrace je dynamický: jeho fragment lze nahradit bez přechodu ze stránky. Moduly vlastnící prvky panelu proto používají delegované vazby událostí nebo vazby zohledňující životní cyklus, aby nově vykreslená kopie nadále fungovala a~starý fragment již nevlastnil posluchače.

\subsection{Veřejný postup vykreslení vybrané galerie}
\index{p00148@postup vykreslení}

Požadavek vybrané galerie nejprve vyhledá galerii a~její účinná pravidla přístupu včetně zděděné ochrany, stavu hesla a~sdílení, stavu zveřejnění a~věkových omezení. Do vykreslení postupuje pouze povolený obsah.

Řadič poté načte přímé podřízené galerie a~aktuální stránku fotografií spolu s~metadaty potřebnými pro názvy, popisy, štítky, data, hlasy, pořadí a~navigaci. Metadata náhledů se připravují v~dávkách a~manifesty médií platné pro požadavek popisují povolené kandidáty dostupné pro každý viditelný obrázek.

Pohled dostane připravený model a~vykreslí název stránky, značku, navigaci, drobečkovou navigaci, karty, stránkování, konfiguraci lightboxu a~podklady zápatí. Moduly prohlížeče po doručení výsledek vylepší. Přístup zůstává odpovědností serveru; JavaScript může interakci zrychlit nebo zpohodlnit, ale nemění nepovoleného kandidáta na povoleného.

Toto rozdělení s~prioritou serveru zachovává užitečné chování i~při blokovaných nebo opožděných skriptech. Přímé odkazy, významový obsah a~počáteční volby obrázků jsou již v~HTML, zatímco JavaScript obsluhuje funkce skutečně vyžadující interaktivního klienta.

\section{Volitelné souvislosti: co databáze uchovává}
\label{sec:database}
\index{d00025@databázové schéma}
\index{m00086@migrace}

Databáze ukládá stav aplikace, který nepatří do názvů souborů obrázků: popisy, soukromí, pořadí, štítky, hlasy, data, veřejné adresy, účty a~historii údržby. Níže uvedený přehled postačuje pro plánování záloh a~jednání o~hostingu; definice jednotlivých sloupců zůstávají v~databázové referenci \cite{database}.

\subsection{Hlavní oblasti}

\begin{longtable}{>{\bfseries}>{\hspace{0pt}}p{0.25\linewidth}>{\hspace{0pt}}p{0.67\linewidth}}
\toprule
Oblast & Odpovědnost za trvalá data\\
\midrule
\endhead
Uživatelé a~autentizace & Identita správce, hashe hesel, e-mail pro obnovu, propojení externích identit, tokeny obnovy a~přihlašovací údaje s~vymezeným rozsahem.\\
Galerie & Cesta složky, hierarchie, veřejná cesta, metadata, viditelnost, přístup, vlastní značka, přepsání prezentace a~provozní pořadí.\\
Obrázky & Vlastnictví galerie, relativní cesta, název souboru, metadata, rozměry, viditelnost, kontrolní součet, údaje EXIF, pořadí a~odvozený stav.\\
Štítky & Opakovaně použitelná identita štítku a~vztahy s~galeriemi a~obrázky.\\
Hlasy & Záznamy hodnocení galerií a~obrázků spojené s~návštěvníkem.\\
Náhledy & Soupis platných generovaných variant, rozměry, formát, stav a~verze odvozeniny.\\
Nastavení & Skalární hodnoty aplikace, vzhledu, funkcí, aktualizací, údržby a~vykreslování.\\
Audit a~telemetrie & Události správce, provozní diagnostika, měření s~řízeným soukromím a~souhrny.\\
Záznamy duplicit & Kanonická rozhodnutí párů a~pravidla potlačení přesných galerií pro každého správce.\\
Úlohy a~tokeny & Omezený stav prohlížeče, detektoru, migrací, WebDAV, obnovy hesla, sdílení a~údržby tam, kde je nutné trvalé uložení.\\
\bottomrule
\end{longtable}

\subsection{Vztahy a~kaskády}

Cizí klíče vyjadřují vlastnictví tam, kde to vývoj schématu a~kompatibilita dovolují. Kaskády odstraňují závislé záznamy postupů při zániku nadřazené identity. Odstranění obsahu nadále probíhá službami aplikace, aby dopady na souborový systém a~databázi zůstávaly sladěné. Omezení jedinečnosti vyjadřují identity, například kanonické duplicitní páry, přesná pravidla správce a~galerie, identifikátory v~URL, tokeny a~varianty metadat.

\subsection{Nastavení aplikace}

Tabulka \codeword{app_settings} ukládá konfiguraci klíč--hodnota. \codeword{app_setting()} čte hodnoty, \codeword{set_app_setting()} ukládá normalizované hodnoty a~cílené služby interpretují doménová nastavení. Konfigurace veřejného jazyka používá \setting{public_language} pro výchozí hodnotu webu, \setting{public_language_selector_enabled} pro funkci přepsání návštěvníkem, která je ve výchozím stavu zapnutá, a~\setting{public_language_selector_languages} pro její seřazený neprázdný seznam jazyků. Veřejný vykreslovač náhledů používá \setting{public_thumbnail_rendering_mode} s~hodnotami \setting{responsive} a~\setting{progressive}.

\section{Volitelná reference pro vývojáře}
\label{sec:development}
\index{v00267@vývoj}
\index{s00229@styl kódu}

Tato reference je pro osoby měnící samotné PHP Gallery. Vlastníci potřebující pouze kontroly zveřejnění a~údržby mohou pokračovat oddílem~\ref{sec:testing}.

\subsection{Uspořádání repozitáře}

\begin{longtable}{>{\ttfamily}>{\hspace{0pt}}p{0.27\linewidth}>{\hspace{0pt}}p{0.65\linewidth}}
\toprule
\textrm{Umístění} & \textrm{Odpovědnost}\\
\midrule
\endhead
app/controllers/ & Řadiče HTTP a~koordinace odpovědí.\\
app/services/ & Opakovaně použitelné doménové a~infrastrukturní služby.\\
app/views/ & Rozvržení a~opakovaně použitelné pomocné funkce pohledů vykreslených serverem.\\
app/lang/ & Překladové katalogy serveru a~prohlížeče.\\
database/migrations/ & Vývoj schématu s~časovými razítky.\\
public/assets/ & JavaScript bez frameworku, CSS a~veřejné podklady.\\
scripts/ & Nástroje migrací, integrity, nasazení a~údržby.\\
tests/ & Samostatné testy PHP a~cílené testy JavaScriptu.\\
winapp/ & Nástroje a~integrace specifické pro Windows.\\
\bottomrule
\end{longtable}

\subsection{Pravidla kódování}

Nové soubory PHP používají přísné typy a~odsazení čtyřmi mezerami. Řadiče koordinují chování požadavků, služby vlastní opakovaně použitelnou logiku. Nový JavaScript a~CSS zůstávají bez frameworku a~patří pod veřejné podklady. Migrace používají názvy souborů s~časovým razítkem na začátku. Existující vysvětlující komentáře, dokumentační řetězce a~PHPDoc se při vývoji chování zachovávají a~opravují.

Samostatné moduly by měly vlastnit jednu soudržnou odpovědnost. Společné kontrakty si zaslouží normalizované vstupy, výslovné výstupy a~cílené testy. Dynamické fragmenty administrace a~veřejnosti vyžadují bezpečnou inicializaci a~úklid podle životního cyklu.

\subsection{Přidání funkce}

Typická změna funkce postupuje takto:

\begin{enumerate}
  \item Prozkoumejte současnou implementaci a~určete vlastnící řadič, služby, pohledy, moduly prohlížeče, trvalé uložení a~testy.
  \item Definujte kontrakty přístupu, CSRF, rozsahu, náhradního chování, migrace a~provozu bez JavaScriptu.
  \item Přidejte cílené služby a~integraci tras.
  \item Pokud jsou nutné změny trvalého schématu, přidejte nové migrace bez změn historie.
  \item Přidejte přeložené řetězce rozhraní.
  \item Přidejte cílené automatické testy a~dokumentované ruční ověření.
  \item Aktualizujte architekturu, mapu kódu, databázi, testování a~uživatelské pokyny podle jejich vlastnictví.
  \item Po konečných změnách běhového kódu obnovte metadata integrity.
  \item Zkontrolujte úplný rozdíl změn a~spusťte příslušnou místní sadu testů.
\end{enumerate}

\section{Kontrola galerie před zveřejněním}
\label{sec:testing}
\index{t00234@testování}
\index{z00273@zajištění kvality}

Testování pro vlastníka galerie znamená projít web jako zamýšlené publikum. Technicky zdravá stránka může stále obsahovat soukromé místo, nejasný název, nesprávné datum, neúmyslné povolení stahování nebo nepohodlné mobilní rozvržení. Krátká lidská kontrola tyto problémy prezentace a~soukromí odhalí.

\subsection{Kontrola zveřejnění očima návštěvníka}
\index{k00070@kontrolní seznam zveřejnění}

Otevřete soukromé okno prohlížeče a~přejděte stejným odkazem, který plánujete sdílet. Zkontrolujte:

\begin{enumerate}
  \item Název a~první odstavec vysvětlují sbírku.
  \item Vybraný titulní obrázek zůstává jasný v~malé velikosti.
  \item Podřízené galerie se zobrazují ve smysluplném pořadí.
  \item Názvy a~popisy fotografií odpovídají viditelným motivům.
  \item Soukromé fotografie a~galerie zůstávají skryté.
  \item Přístup heslem nebo sdílením funguje z~nové relace.
  \item Mapy GPS odhalují pouze zamýšlená místa.
  \item První stránka se na telefonu načítá pohodlně.
  \item Velký prohlížeč, ovládání dopředu a~zpět a~zavření působí přirozeně.
  \item Vyhledávání, štítky, hlasování a~stahování se chovají podle účelu galerie.
\end{enumerate}

Požádejte další osobu o~vyzkoušení galerie bez pokynů. Její první kliknutí a~první otázka často odhalí navigaci či formulaci, která byla zřejmá pouze vlastníkovi.

\subsection{Automatické kontroly pro správce softwaru}

Vývojáři mohou po změně kódu aplikace spustit spustitelné testovací skripty projektu. Souhrnný nástroj pokrývá sadu projektu, cílené skripty ověřují jednotlivé funkce.

\begin{lstlisting}[language=bash]
php tests/run.php
php tests/public_thumbnail_rendering_model_test.php
php tests/public_thumbnail_markup_test.php
php tests/duplicate_photo_detector_test.php
php tests/duplicate_photo_ledger_test.php
php tests/migration_consistency_test.php
node tests/progressive_thumbnail_renderer_test.mjs
\end{lstlisting}

\subsection{Ruční ověření v~prohlížeči}

\subsubsection{Kontroly sítě a~interakcí}

Kontroly softwaru závislé na prohlížeči používají reprezentativní data, anonymní a~ověřené relace, prázdnou mezipaměť, mobilní a~počítačové velikosti zobrazení, provoz bez JavaScriptu tam, kde je relevantní, preference omezeného pohybu a~zpomalování sítě.

U~veřejného vykreslování náhledů prozkoumejte panely Prvky a~Síť. Responzivní režim má zpřístupnit úplnou aktivní sadu kandidátů. Progresivní režim má před aktivací zpřístupnit malého aktivního kandidáta a~neaktivní větší kandidáty a~poté provádět nejvýše dokumentovaný počet souběžných vylepšení příslušných karet. Zkontrolujte stabilitu rozvržení, opětovné použití mezipaměti, náhradu při chybě, interakce lightboxu, mapy, hlasování, chráněná média a~přímé odkazy.

\subsection{Čistota vydání}

Kontrola vydání ověřuje konzistenci běhové verze, kompatibilitu seřazených migrací, hashe manifestu, dokumentaci, překlady, obnovu mezipaměti prohlížeče změnou verzování, testy, vyloučení z~balíčku a~bezpečnost uživatelských dat. Správce začíná z~čisté větve vydání, porovnává ji s~přesným tagem předchozího vydání, kontroluje každý mezilehlý commit a~cestu a~zaznamenává vše záměrně vyloučené. Běhová verze v~\filepath{app/bootstrap.php}, klíč a~tag \filepath{release-metadata.json}, nejnovější nadpis \filepath{PATCH_NOTES.md}, značky dokumentace, zamýšlený název archivu a~konečný tag Git se musejí shodovat.

Nové migrace se kontrolují v~pořadí časových razítek a~zkoušejí testy konzistence migrací a~staršího migračního nástroje. Správce znovu sestaví toto PDF úplnou posloupností \filepath{docs/LATEX_BUILD.md} a~prohlédne titulní stránku, přehled vydání, obsah, záložky, rejstřík a~změněné oddíly funkcí. Každý změněný soubor PHP projde \codeword{php -l}; změněný JavaScript a~testovací podklady Node projdou kontrolou syntaxe; cílené testy funkcí, konzistence překladů, pokrytí dokumentace, úplná sada PHP a~\codeword{git diff --check} musejí uspět.

Teprve po konečné úpravě zdrojů a~dokumentace správce spustí \codeword{php scripts/generate_manifest.php} a~jeho režim \codeword{--check}. Pomocný nástroj nasazení poté vytvoří požadovanou místní složku nebo ZIP; pouze balí soubory a~nespouští PHP ani neopravuje zastaralý manifest. Prohlédněte seznam archivu a~ověřte přítomnost aktuálního PDF, poznámek k~vydání, metadat vydání, migrací a~\filepath{app/core-manifest.json}, zatímco chráněná konfigurace, mezipaměti, protokoly, dočasné soubory, místní nástroje, výstup nasazení a~uživatelská média následují pravidla balení. Po zabalení znovu zkontrolujte manifest, vytvořte commit vydání a~anotovaný tag \codeword{v_<version>} z~ověřeného stromu a~proveďte základní zkoušku upgradu aktualizátorem z~předchozího stabilního vydání. Chybějící místní závislosti se musejí hlásit s~přesnými blokovanými příkazy a~podrobnostmi prostředí místo označení vydání za úplně ověřené.

\section{Jak zajistit svižný dojem galerie}
\label{sec:performance}
\index{v00262@výkon}
\index{r00198@rychlost}

Vnímaná rychlost není jediným číslem. Návštěvníci vnímají první zobrazení stránky, rozpoznatelnost prvních fotografií, plynulost posouvání, rychlost reakce ovládacích prvků a~dobu načítání většího obrázku po vyžádání. Přispívají kvalita připojení, vzdálenost serveru, velikost stránky, pravidla náhledů, podklady značky a~vybraný vykreslovač.

\subsection{Praktické způsoby zrychlení}

\begin{itemize}
  \item Udržujte přiměřený počet karet na stránku pro cílové publikum a~hostingový tarif.
  \item Pro velké veřejné sbírky předem vygenerujte nastavené velikosti náhledů.
  \item Vyhněte se zbytečně velkým úvodním obrázkům, pozadím, logům a~dalším podkladům vzhledu.
  \item Upřednostněte responzivní vykreslovač náhledů, pokud jsou nejdůležitější předvídatelnost a~chování bez JavaScriptu; progresivní vykreslovač použijte, pokud sbírce a~zastoupení prohlížečů vyhovuje jeho postupné zostřování.
  \item Pokud je to praktické, umístěte galerii na hosting geograficky blízko hlavního publika.
  \item Testujte s~prázdnou mezipamětí. Prohlížeč vývojáře s~naplněnou mezipamětí může skrýt nákladné přenosy první návštěvy.
  \item Než budete předpokládat, že každá pomalá galerie má problém s~náhledy, zkontrolujte latenci PHP a~databáze odděleně od doby přenosu obrázků.
\end{itemize}

\subsection{Test pomalého připojení}

Otevřete galerii v~soukromém okně prohlížeče, zvolte reprezentativní mobilní zobrazení, zapněte zpomalení sítě ve vývojářských nástrojích a~obnovte stránku s~prázdnou mezipamětí. Sledujte pořadí příchodu stránky, titulních obrázků, viditelných karet a~pozdějších vylepšení obrázků. Zopakujte s~nastavením vykreslovače nebo velikosti stránky, které chcete porovnat; jinak soubory v~mezipaměti zneplatní spolehlivost porovnání.

\subsection{Technická měření pro správce softwaru}

Profilování veřejného vykreslování a~diagnostika prohlížeče jsou užitečné, když subjektivní zpráva potřebuje čísla. Zaznamenejte alespoň:

\begin{itemize}
  \item čas do prvního bajtu a~dobu vykreslení na serveru;
  \item velikost přenosu HTML;
  \item počet požadavků a~přenesené bajty počátečních náhledů;
  \item čas do ustálení největšího viditelného obsahu;
  \item posuny rozvržení při načítání karet a~médií;
  \item práci hlavního vlákna přičitatelnou modulům prohlížeče;
  \item počet a~velikost vylepšení progresivních kandidátů;
  \item celkové počty databázových dotazů vývojového profileru, pokud jsou dostupné;
  \item tlak dekódování obrázků na paměť při testování velmi velkých originálů nebo výrazného přiblížení lightboxu.
\end{itemize}

Vývojový profiler a~zachycené záznamy měření jsou určeny k~porovnání před a~po napříč režimy vykreslovačů a~rozvrženími galerií. Používejte reprezentativní obrazovky s~vysokou hustotou a~omezená připojení místo spoléhání pouze na rychlý počítač s~naplněnou mezipamětí.
\section{Řešení problémů}
\label{sec:troubleshooting}
\index{r00202@řešení problémů}

\subsection{Aplikace se nespouští}

Ověřte verzi PHP, požadovaná rozšíření, umístění konfigurace, připojení k~databázi, zapisovatelné cesty a~směrování kořene webu. Prohlédněte protokoly PHP a~webového serveru. Spusťte kontroly syntaxe nedávno změněných souborů PHP a~prozkoumejte čekající migrace.

\subsection{Galerie nebo obrázky chybějí}

Ověřte cesty souborového systému, stav vyhledání galerie, databázové záznamy, viditelnost, pravidla přístupu rodiče, stav odemknutí heslem, pravidla NSFW, podporovaný typ souboru a~relativní cestu obrázku. Otestujte jako správce i~anonymní návštěvník.

\subsection{Náhledy chybějí nebo jsou neostré}

Prohlédněte metadata náhledů a~zprávy údržby. Ověřte generované velikosti, podporu formátů, nastavené meze, geometrii zdroje, stav odvozeniny a~oprávnění veřejného koncového bodu. V~responzivním režimu prozkoumejte kandidáta vybraného prohlížečem a~\codeword{sizes}. V~progresivním režimu prozkoumejte malého kandidáta, stav fronty, vybranou náhradu a~výsledek dekódování.

\subsection{Akce administrace opouštějí postranní panel}

Ověřte načtení odpovídajícího modulu JavaScriptu, rozpoznání formuláře delegovanými obsluhami, správnost hlaviček AJAX a~kontraktů odpovědi, přítomnost očekávaných značek životního cyklu ve vráceném fragmentu a~vyloučení formulářů vlastněných funkcemi z~obecných obsluh formulářů administrace. Přímé chování POST může nadále poskytovat dokumentovanou náhradní trasu.

\subsection{Migrace selhává}

Uchovejte ověřenou zálohu databáze a~zaznamenejte název migračního souboru, verzi Gallery, verzi databázového serveru, omezenou zprávu na obrazovce, odkaz na požadavek a~existující stav schématu a~evidence migrací. Zopakujte postup ze stránky správy galerií, aby běžná idempotentní kompatibilní cesta mohla dokončit přerušenou migraci. Verze~\version{} nepotřebuje triggery, uložené rutiny, \codeword{SUPER}, \codeword{log\_bin\_trust\_function\_creators} ani globální změnu databáze. Pokud omezená chyba přetrvává, prozkoumejte Diagnostiku běhového prostředí a~předejte odkaz na požadavek správci softwaru nebo poskytovateli hostingu; do zprávy podpory nevkládejte přihlašovací údaje databáze ani nezpracovaná produkční data.

\subsection{Kompilace příručky PDF selhává}

Spusťte příkazy v~\filepath{docs/LATEX_BUILD.md}. Ověřte instalaci \codeword{pdflatex} a~\codeword{makeindex}. První spuštění MiKTeX může požadovat běžné balíčky. Po neúspěšném přechodu balíčků odstraňte ignorované mezivýsledky sestavení a~zopakujte dokumentovanou posloupnost.

\section{Kontrolní seznam zálohování a~obnovy}
\label{sec:backup}
\index{o00112@obnova}

\begin{enumerate}
  \item Exportujte úplnou databázi aplikace.
  \item Zkopírujte \filepath{galleries/} se zdrojovými soubory a~podklady vlastněnými galeriemi.
  \item Bezpečně zkopírujte místní konfiguraci a~příslušné tajné údaje instalace.
  \item Uchovejte vlastní podklady vzhledu a~vlastní CSS.
  \item Zaznamenejte verzi aplikace a~stav čekajících migrací.
  \item Ukládejte zálohy mimo aktivní webový strom.
  \item Pravidelně testujte obnovu v~izolovaném prostředí.
  \item Po obnově ověřte veřejné trasy, chráněný přístup, přihlášení správce, náhledy, nahrávání a~migrace.
\end{enumerate}

Zálohy vytvořené aktualizátorem uchovávají přepsané soubory aplikace pro obnovu aktualizace. Úplná provozní záloha zahrnuje také databázi, strom galerií, místní konfiguraci a~podklady vlastněné instalací.

\subsection{Ověření sady pro obnovu}
\index{o00112@obnova!o00130@ověření}

Uchovávejte export databáze, původní fotografie, místní konfiguraci, příslušné úložiště \filepath{data/}, obsah Koše, vlastní podklady a~přesnou verzi aplikace společně jako jednu sladěnou sadu pro obnovu. Snímek poskytovatele nebo omezené okno údržby mohou zabránit zápisům rozdělit stav databáze a~souborů. Uložte kopii mimo hosting a~dohodněte, kolik nedávné práce může být ztraceno a~jak rychle musí být služba obnovena.

Vstupní bod CLI \filepath{scripts/recovery.php} může sadu inventarizovat a~ověřit vzorky obnovených souborů a~hashe manifestu v~samostatném umístění. Formát podkladů a~příkazy najdete v~\filepath{docs/RECOVERY_ASSURANCE.md}. Nástroj nespouští databázový dump, neobnovuje produkci ani neověřuje poskytovatele záloh. Jeho syntetické cvičení poškození dokazuje, že kontrola rozpozná poškození testovacích podkladů, nikoli že skutečná záloha dokáže obnovit web.

Před otevřením obnovené instalace izolujte e-mail, plánované úlohy, integrace a~zápis od produkčních cílů. Ověřte obnovenou databázi, vztahy, přihlášení správce, odmítnutí soukromých médií a~obnovu Koše. Zaznamenejte stáří zálohy, uplynulý čas obnovy a~pozorování obsluhy bez kopírování tajných údajů do zpráv. Úplný postup najdete v~\filepath{docs/RECOVERY_OFF_HOST.md}. Vrácení souborů aktualizátorem stále nevrací databázové migrace.

\section{Ověření správcem softwaru a~podklady vydání}
\label{sec:release-evidence}
\index{k00072@kvalifikace vydání}
\index{t00234@testování!d00031@dočasné postupy}

Automatická regrese a~provozní přejímka odpovídají na odlišné otázky. Centrální audit kontroluje kód, kontrakty, syntaxi, dostupné testovací podklady prohlížeče a~konzistenci vydání. Izolované testovací prostředí přidává skutečné průchody databází, HTTP a~Chromium s~generovanými fotografiemi a~účty. Vyžaduje výslovnou konfiguraci určenou pouze pro dočasné testování a~nikdy nesmí používat databázi ani média instalace. Předpoklady najdete v~\filepath{docs/GALLERY_WORKFLOWS.md}; chybějící prostředí je mezerou pokrytí, nikoli důkazem fungující funkce.

Při přípravě vydání postupujte podle \filepath{RELEASE.md}: dokončete zdroje a~dokumentaci, znovu sestavte a~prohlédněte toto PDF, vygenerujte manifest a~inicializujte kvalifikační záznam. Spusťte jeden centrální audit vydání, nikoli posloupnost quick/full/release. S~nastavenými izolovanými předpoklady přijímá \filepath{scripts/gallery_workflow_mysql.php} volbu \codeword{--release} pro vytvoření testovacího prostředí kolem stejného auditu vydání.

Kvalifikační nástroj \filepath{scripts/release_qualification.php} zaznamenává přesný otisk zdroje a~PDF. Jeho příkaz \codeword{render} vytváří omezené náhledy stránek; samotné vykreslení neschvaluje jejich vzhled. Příkaz \codeword{record-audit} přijímá pouze odpovídající centrální zprávu vydání, nikoli starší úplný audit. Změněný zdrojový soubor nebo znovu sestavené PDF vybírají nové čekající podklady místo tichého opětovného použití dřívějšího souhlasu.

Prohlédněte titulní stránku, obsah, rejstřík, změněné stránky, interní odkazy a~reprezentativní postupy prohlížeče. Zaznamenejte, co bylo skutečně zkontrolováno, kým a~s~jakým výsledkem. Nevyřešené kontroly ponechte čekající a~základní test aktualizátoru po zveřejnění oddělený. Zelená automatická zpráva není svolením ke zveřejnění a~nástroj nikdy nevytváří commity ani tagy, nenahrává ani nepublikuje vydání. Příkazy záznamů najdete v~\filepath{docs/RELEASE_QUALIFICATION.md}.

\section{Slovník}
\label{sec:glossary}
\index{s00209@slovník}

\begin{description}
  \item[AJAX] Požadavek prohlížeče používaný k~aktualizaci oblasti stránky při zachování okolního rozhraní.
  \item[Rozsah větve] Galerie a~její potomci, pokud vybraná funkce definuje rekurzivní rozsah.
  \item[Kanonický pár] Dva identifikátory obrázků uložené ve stabilním pořadí od nižšího k~vyššímu.
  \item[Token CSRF] Hodnota vázaná na relaci pro ověření záměrného odeslání formuláře měnícího stav.
  \item[Odvozenina] Generovaná podoba média, například náhled JPEG nebo WebP.
  \item[EXIF] Metadata fotoaparátu a~pořízení vložená v~podporovaných souborech obrázků.
  \item[Manifest integrity] Verzovaný seznam cest jádra a~očekávaných hashů.
  \item[Lightbox] Celoobrazovkový nebo překryvný prohlížeč fotografií s~navigací a~souvisejícím ovládáním.
  \item[Manifest médií] Data vykreslování platná pro požadavek, připravená z~metadat náhledů viditelných obrázků.
  \item[Ochrana NSFW] Pravidla věku a~galerií či obrázků řídící přístup k~omezeným médiím.
  \item[Postupné zostřování] Postup náhledů s~malou kopií jako první, který později aktivuje dekódovaného většího kandidáta.
  \item[Responzivní výběr] Nativní výběr kandidáta prohlížečem z~okamžitě aktivní sady zdrojů.
  \item[Vybraná galerie] Galerie představovaná aktuální veřejnou trasou.
  \item[Postranní panel] Kontextové rozhraní administrace vpravo, načítané a~obnovované požadavky fragmentů.
  \item[Meze náhledů] Minimální a~maximální generované velikosti způsobilé pro výběr se zohledněním galerie.
\end{description}

\clearpage
\section{Zdroje}
\label{sec:references}

\begin{thebibliography}{99}

\bibitem{readme}
Projekt PHP Gallery, \emph{README.md}, místní přehled produktu, průvodce funkcemi, požadavky a~dokumentace instalace, verze \version.

\bibitem{architecture}
Projekt PHP Gallery, \emph{ARCHITECTURE.md}, místní reference architektury aplikace a~životního cyklu subsystémů, verze \version.

\bibitem{database}
Projekt PHP Gallery, \emph{DATABASE.md}, místní reference migrací a~trvalého schématu, verze \version.

\bibitem{codemap}
Projekt PHP Gallery, \emph{CODEMAP.md}, místní mapa vlastnictví zdrojového kódu a~údržby, verze \version.

\bibitem{settingsinventory}
Projekt PHP Gallery, \emph{docs/ADMIN\_SETTINGS\_INVENTORY.md}, kanonický soupis vlastnictví globálního Nastavení, náhradního chování, citlivosti a~migrací, verze \version.

\bibitem{testing}
Projekt PHP Gallery, \emph{TESTING.md}, místní průvodce automatickým a~ručním ověřováním, verze \version.

\bibitem{agents}
Projekt PHP Gallery, \emph{AGENTS.md}, místní pravidla přispívání do repozitáře a~zabezpečení.

\bibitem{phpmanual}
The PHP Documentation Group, \emph{Příručka PHP}, \url{https://www.php.net/manual/en/}.

\bibitem{mysqlmanual}
Oracle Corporation, \emph{Referenční příručka MySQL}, \url{https://dev.mysql.com/doc/}.

\bibitem{mariadbmanual}
MariaDB Foundation, \emph{Dokumentace serveru MariaDB}, \url{https://mariadb.com/docs/server/}.

\bibitem{htmlimages}
WHATWG, \emph{Průběžně aktualizovaný standard HTML: Obrázky}, \url{https://html.spec.whatwg.org/multipage/images.html}.

\bibitem{owasp}
OWASP Foundation, \emph{Pokyny k~zabezpečení webových aplikací}, \url{https://owasp.org/}.

\end{thebibliography}

\clearpage
\phantomsection
\addcontentsline{toc}{section}{Rejstřík}
\printindex

\end{document}
